\chapter{Linear algebra}
\label{vol02:ch09}

% \spn, \col, \nul, \rank, \sgn, \proj are defined in eng-macros.sty.

\epigraph{The column space tells you what is reachable; the null space
tells you what is invisible; everything else in linear algebra
follows from understanding these two subspaces.}{Paraphrase of
the working orientation in Strang, \emph{Introduction to Linear
Algebra}, 5th ed.\ \cite[ch.~3]{text:strang2016}.}

\chapmeta{Half-life: foundational. Archetypes: optimisation
(least-squares fitting as the canonical convex problem), uncertainty
(conditioning as the formal handling of input-error amplification).
Exercise target: 40. Project track: simulation with analysis (a
polynomial fit by three methods with conditioning comparison, plus a
1500-word reflection). Hazard class: none. Reader path default:
standard, with sections explicitly tagged. Named cases: Ariane 5
Flight 501 (1996), as the canonical case of a numerical conversion
that should have been refused; Hubble primary mirror (1990) and
Patriot missile clock drift at Dhahran (1991) as short-form
companion cases of the same precondition-verification discipline.}

\noindent
A linear system is the engineer's working notation for a problem in
which several quantities depend, jointly and proportionally, on
several others. Volume I Chapter 8 fitted a straight line to data by
formulas the reader took on trust; this chapter supplies the
apparatus that earns the formulas. The same apparatus carries the
reader from a hand-solved $3 \times 3$ system to the polynomial fit
of section 9.6, the conditioning analysis of section 9.7, and the
refusal discipline of section 9.8. Two textbooks supplement the
working derivations: Strang \cite{text:strang2016} for the
subspace-first orientation and an undergraduate-accessible
abstract treatment in Axler \cite{text:v2ch09-axler2015}; Trefethen
and Bau \cite{text:trefethen-bau1997} for numerical linear algebra
and Golub and Van Loan \cite{text:v2ch09-golub-vanloan2013} for the
encyclopaedic algorithmic reference; Higham
\cite{text:v2ch09-higham2002} for the backward-stability and
condition-number analysis the failure section depends on.

\section{Vector spaces}\pathtag{core}

A \engterm{vector space} over the real numbers is a set $V$ on which
two operations are defined, vector addition and scalar
multiplication, satisfying the eight axioms recorded in any
linear-algebra text \cite[ch.~3]{text:strang2016}. We do not treat
the axioms as a memorisation exercise. The working content for an
engineer sits in three concepts that follow from them: the
\emph{spanning set}, the \emph{linearly independent set}, and the
\emph{basis}.

The standard concrete example is $\R^{n}$, the set of $n$-tuples of
real numbers under componentwise addition and scalar multiplication.
A vector $\vect{v} \in \R^{n}$ is written as a column,

\[
\vect{v} = \begin{pmatrix} v_{1} \\ v_{2} \\ \vdots \\ v_{n}
\end{pmatrix},
\qquad v_{i} \in \R.
\]

We adopt the column convention throughout, consistent with Volume I
Chapter 8's regression notation. Row vectors appear when we transpose
a column with $\transpose$.

\subsection{Spans, independence, basis}

A \engterm{linear combination} of vectors $\vect{v}_{1}, \ldots,
\vect{v}_{k}$ is any vector of the form $\sum_{i=1}^{k} c_{i}
\vect{v}_{i}$ with scalar coefficients $c_{i}$. The \engterm{span}
$\spn(\vect{v}_{1}, \ldots, \vect{v}_{k})$ is the set of all such
combinations. Spans are subspaces: the span of a finite set of
vectors is closed under addition and scalar multiplication.

A set $\{\vect{v}_{1}, \ldots, \vect{v}_{k}\}$ is \engterm{linearly
independent} when the only solution of $\sum c_{i} \vect{v}_{i} =
\vect{0}$ is $c_{1} = \cdots = c_{k} = 0$. Equivalently, no vector
in the set is a linear combination of the others. When a set fails
this condition it is \engterm{linearly dependent}, and at least one
vector carries no information that the others do not already carry.

A \engterm{basis} for a subspace $W$ is a linearly independent set
that spans $W$. Every basis of a finite-dimensional subspace has the
same number of vectors; that number is the \engterm{dimension}
$\dim W$. The standard basis of $\R^{n}$ is $\{\vect{e}_{1}, \ldots,
\vect{e}_{n}\}$, where $\vect{e}_{i}$ has $1$ in position $i$ and
zeros elsewhere.

The engineering working content of these definitions is direct.
Spanning answers \emph{which targets are reachable} by combining the
chosen building blocks. Independence answers \emph{which building
blocks are wasted}. A basis is the smallest set of building blocks
that still reaches everything; its size is the dimension, and the
dimension is what counts the system's degrees of freedom.

The test for independence is a homogeneous linear system. Stack the
candidate vectors as the columns of a matrix and ask whether
$\mat{A}\vect{c} = \vect{0}$ has a nonzero solution. It does exactly
when the columns are dependent. Section 9.2 makes the same statement
in the language of the null space, and section 9.3 supplies the
elimination that answers it in finite arithmetic.

\begin{example}[Testing independence by hand]
Are $\vect{v}_{1} = (1, 2, 3)^{\transpose}$, $\vect{v}_{2} = (2, 5,
7)^{\transpose}$, and $\vect{v}_{3} = (1, 3, 4)^{\transpose}$
independent? Form $\mat{A} = [\vect{v}_{1} \mid \vect{v}_{2} \mid
\vect{v}_{3}]$ and reduce. Subtracting $2\times$ row 1 from row 2 and
$3\times$ row 1 from row 3,
\[
\begin{pmatrix} 1 & 2 & 1 \\ 2 & 5 & 3 \\ 3 & 7 & 4 \end{pmatrix}
\;\longrightarrow\;
\begin{pmatrix} 1 & 2 & 1 \\ 0 & 1 & 1 \\ 0 & 1 & 1 \end{pmatrix}
\;\longrightarrow\;
\begin{pmatrix} 1 & 2 & 1 \\ 0 & 1 & 1 \\ 0 & 0 & 0 \end{pmatrix},
\]
where the last step subtracts row 2 from row 3. The third pivot is
missing, so the rank is two and the three vectors are dependent. The
free column is the third; setting its coefficient $c_{3} = 1$ and
back-solving the echelon system gives $c_{2} = -1$ and $c_{1} = 1$,
so the dependence is $\vect{v}_{1} - \vect{v}_{2} + \vect{v}_{3} =
\vect{0}$. Substituting confirms $(1,2,3) - (2,5,7) + (1,3,4) =
(0,0,0)$. The echelon form's free-column structure is the reliable
generator of the relation; a dependence guessed from a glance must
always be verified by substitution. Had the third vector been
$(1,3,5)^{\transpose}$ instead, the elimination would have produced a
third pivot and reported three independent vectors.
\end{example}

\begin{example}[A basis from a spanning set]
The four vectors $\vect{a}_{1} = (1,0,1)^{\transpose}$, $\vect{a}_{2}
= (0,1,1)^{\transpose}$, $\vect{a}_{3} = (1,1,2)^{\transpose}$,
$\vect{a}_{4} = (2,1,3)^{\transpose}$ span a subspace of $\R^{3}$.
A spanning set with four vectors in a three-dimensional ambient
space cannot be independent. Reducing $[\vect{a}_{1} \mid \cdots
\mid \vect{a}_{4}]$ finds pivots in columns one and two; columns
three and four are combinations of them ($\vect{a}_{3} = \vect{a}_{1}
+ \vect{a}_{2}$ and $\vect{a}_{4} = 2\vect{a}_{1} + \vect{a}_{2}$).
The pivot
columns $\{\vect{a}_{1}, \vect{a}_{2}\}$ form a basis for the span,
which is the plane $\{(x, y, z) : x + y = z\}$ through the origin.
The dimension is two. The redundant vectors $\vect{a}_{3}$ and
$\vect{a}_{4}$ carry no reach the basis does not already provide.
\end{example}

\begin{mastery}
Vector spaces need not be coordinate spaces. The set of polynomials
of degree at most $n$ is a vector space of dimension $n + 1$, with
the monomial basis $\{1, x, x^{2}, \ldots, x^{n}\}$. The set of
$m \times n$ real matrices is a vector space of dimension $mn$. The
set of solutions of a homogeneous linear ordinary differential
equation of order $k$ is a vector space of dimension $k$ (Chapter 7).
The abstract definition earns its generality here: every theorem
proved from the eight axioms, including the existence and size of a
basis, applies verbatim to function spaces and matrix spaces. The
engineer who has internalised the coordinate case reads the abstract
result as a translation, with the inner product of section 9.5
becoming the $L^{2}$ integral $\inner*{f, g} = \int f g \,\dd x$ and
the Gram-Schmidt procedure producing the Legendre and Chebyshev
polynomials that section 9.6 invokes against ill-conditioning. The
finite-element method of Volume IV is the systematic exploitation of
finite-dimensional subspaces of these infinite-dimensional function
spaces.
\end{mastery}

\subsection{Change of basis}

A vector's coordinates depend on the basis chosen to express it.
The vector itself, an arrow in space or a physical displacement, is
basis-independent; its coordinate column is not. If
$\mathcal{B} = \{\vect{b}_{1}, \ldots, \vect{b}_{n}\}$ is a basis of
$\R^{n}$ and $\vect{x}$ has coordinates $\vect{c}$ in $\mathcal{B}$,
then $\vect{x} = \mat{B}\vect{c}$ where $\mat{B} = [\vect{b}_{1}
\mid \cdots \mid \vect{b}_{n}]$. Recovering the coordinates from the
standard-basis representation is the solve $\vect{c} =
\mat{B}^{-1}\vect{x}$, which section 9.3 performs without ever
forming $\mat{B}^{-1}$ explicitly.

\begin{example}[Coordinates in a non-standard basis]
Express $\vect{x} = (4, 5)^{\transpose}$ in the basis
$\mathcal{B} = \{(1,1)^{\transpose}, (1,-1)^{\transpose}\}$. Solve
$\begin{psmallmatrix} 1 & 1 \\ 1 & -1 \end{psmallmatrix}\vect{c} =
\begin{psmallmatrix} 4 \\ 5 \end{psmallmatrix}$. Adding the two
equations gives $2 c_{1} = 9$, so $c_{1} = 4.5$; subtracting gives
$2 c_{2} = -1$, so $c_{2} = -0.5$. The coordinates are $\vect{c} =
(4.5, -0.5)^{\transpose}$, and the reader checks $4.5(1,1) -
0.5(1,-1) = (4, 5)$. This basis is orthogonal but not orthonormal;
section 9.5 shows that an orthonormal basis makes the coordinate
solve a sequence of inner products with no system to invert.
\end{example}

A change of basis moves vectors; it also moves the matrices that act on
them. A linear map represented by $\mat{M}$ in the standard basis is
represented by $\mat{B}^{-1}\mat{M}\mat{B}$ in the basis $\mathcal{B}$,
because the new matrix must first translate coordinates into the
standard basis ($\mat{B}$), apply the map ($\mat{M}$), and translate
back ($\mat{B}^{-1}$). Two matrices related this way, $\mat{M}$ and
$\mat{B}^{-1}\mat{M}\mat{B}$, are \engterm{similar}, and similarity is
the formal statement that they are the same map seen in two coordinate
systems. Similar matrices share every basis-independent quantity: the
determinant, the trace, the rank, the eigenvalues, the characteristic
polynomial. The whole diagonalisation programme of Chapter 10 is the
search for the basis in which a map's matrix is as simple as possible,
and the eigenvector basis is the one in which it is diagonal.

\begin{example}[A map simplified by a change of basis]
The matrix $\mat{M} = \begin{psmallmatrix} 3 & 1 \\ 0 & 2
\end{psmallmatrix}$ acts on $\R^{2}$. In the basis $\mathcal{B} =
\{(1, 0)^{\transpose}, (1, -1)^{\transpose}\}$, with $\mat{B} =
\begin{psmallmatrix} 1 & 1 \\ 0 & -1 \end{psmallmatrix}$ and
$\mat{B}^{-1} = \begin{psmallmatrix} 1 & 1 \\ 0 & -1
\end{psmallmatrix}$ (the basis is its own inverse map here), the map's
matrix becomes
\[
\mat{B}^{-1}\mat{M}\mat{B} = \begin{pmatrix} 1 & 1 \\ 0 & -1
\end{pmatrix}\begin{pmatrix} 3 & 1 \\ 0 & 2 \end{pmatrix}
\begin{pmatrix} 1 & 1 \\ 0 & -1 \end{pmatrix} = \begin{pmatrix} 3 & 0
\\ 0 & 2 \end{pmatrix}.
\]
The similarity transform has diagonalised the map, and the columns of
$\mat{B}$ are its eigenvectors: $(1, 0)^{\transpose}$ with eigenvalue
$3$ and $(1, -1)^{\transpose}$ with eigenvalue $2$. The trace is $5$
and the determinant $6$ in both representations, as similarity
requires, and in the eigenvector basis the action is plain: scale the
first coordinate by $3$ and the second by $2$. Finding this basis is
the eigenvalue problem of Chapter 10; the change-of-basis algebra here
is the machinery that makes the eigenvalue answer useful.
\end{example}

\subsection{Subspaces of real coordinate space}

Three subspaces recur. The \engterm{zero subspace} $\{\vect{0}\}$ has
dimension zero. A \engterm{line through the origin} is a one-dimensional
subspace; its basis is any one nonzero vector on the line. A
\engterm{plane through the origin} is two-dimensional; its basis is
any two linearly independent vectors in the plane. The phrase
``through the origin'' is load-bearing: a line that does not pass
through the origin is an \emph{affine} subset, not a subspace, because
it is not closed under scalar multiplication.

Affine subsets matter for engineering. The set of solutions of a
linear system $\mat{A}\vect{x} = \vect{b}$ with $\vect{b} \neq
\vect{0}$ is affine, not linear. We return to this distinction in
section 9.3.

\begin{estimation}
\textbf{Question.} A small mechanical truss has $n = 8$ joints, each
with two degrees of freedom in the plane, and four joints constrained
to the ground. How many independent displacement modes does the
truss have, before any member stiffness is computed?

\smallskip
\textit{Estimate, before reading on.}

\smallskip
The unconstrained displacement space has dimension $2n = 16$. Each
ground-constrained joint loses two degrees of freedom; the four
constrained joints remove $8$. The remaining displacement space has
dimension $16 - 8 = 8$. The number is the working dimension the
engineer carries into the stiffness analysis. The constraint
counting is exact only when the constraints themselves are linearly
independent, which is the rank check of section 9.3.
\end{estimation}

\section{Matrices as linear maps}\pathtag{core}

A matrix $\mat{A} \in \R^{m \times n}$ is a rectangular array of
real numbers with $m$ rows and $n$ columns. The matrix-vector
product $\mat{A}\vect{x}$ for $\vect{x} \in \R^{n}$ produces a
vector $\vect{A}\vect{x} \in \R^{m}$ whose $i$-th component is

\[
(\mat{A}\vect{x})_{i} = \sum_{j=1}^{n} A_{ij} x_{j}.
\]

Two interpretations of this product are routinely useful, and an
engineer should be fluent in switching between them.

\subsection{The column-space view}

Write $\mat{A} = [\vect{a}_{1} \mid \vect{a}_{2} \mid \cdots \mid
\vect{a}_{n}]$, where $\vect{a}_{j} \in \R^{m}$ is the $j$-th column.
Then

\[
\mat{A}\vect{x} = x_{1}\vect{a}_{1} + x_{2}\vect{a}_{2} + \cdots
+ x_{n}\vect{a}_{n}.
\]

The product $\mat{A}\vect{x}$ is a linear combination of the
columns of $\mat{A}$, with the entries of $\vect{x}$ as coefficients.
The \engterm{column space} $\col(\mat{A})$ is the span of the
columns. It is a subspace of $\R^{m}$, and it is the set of all
right-hand sides $\vect{b}$ for which $\mat{A}\vect{x} = \vect{b}$
has at least one solution.

This view is the one to default to. When the reader looks at a
matrix, the question \emph{what does this matrix do} is answered by
\emph{which combinations of its columns are reachable}.

\begin{example}[Two readings of one product]
Let $\mat{A} = \begin{psmallmatrix} 2 & 0 & 1 \\ 1 & 3 & 0 \\ 0 & 1
& 4 \end{psmallmatrix}$ and $\vect{x} = (1, 2, -1)^{\transpose}$.
The row reading computes each output component as a dot product:
$(\mat{A}\vect{x})_{1} = 2 - 1 = 1$, $(\mat{A}\vect{x})_{2} = 1 + 6
= 7$, $(\mat{A}\vect{x})_{3} = 2 - 4 = -2$. The column reading
computes the same vector as a combination of columns: $1\cdot(2,1,0)
+ 2\cdot(0,3,1) - 1\cdot(1,0,4) = (2,1,0) + (0,6,2) - (1,0,4) =
(1,7,-2)$. The two readings agree, as they must. The row reading is
the one to use for hand arithmetic; the column reading is the one to
use for reasoning about reachability, because it shows $\mat{A}
\vect{x}$ as a point in the span of the columns.
\end{example}

\subsection{The map view}

A matrix $\mat{A} \in \R^{m \times n}$ defines a function $T:
\R^{n} \to \R^{m}$ by $T(\vect{x}) = \mat{A}\vect{x}$. The function
is \engterm{linear}: $T(\alpha\vect{x} + \beta\vect{y}) = \alpha
T(\vect{x}) + \beta T(\vect{y})$ for any scalars $\alpha, \beta$ and
vectors $\vect{x}, \vect{y}$. Conversely, every linear map between
finite-dimensional real vector spaces is realised by a unique matrix
once bases are fixed. The matrix \emph{is} the linear map, in
coordinates.

Composition of maps is matrix multiplication. If $T(\vect{x}) =
\mat{A}\vect{x}$ takes $\R^{n}$ to $\R^{m}$ and $S(\vect{y}) =
\mat{B}\vect{y}$ takes $\R^{m}$ to $\R^{p}$, then $S \circ T$ takes
$\R^{n}$ to $\R^{p}$ and its matrix is the product $\mat{B}\mat{A}$.
The order matters: $\mat{B}\mat{A}$ applies $\mat{A}$ first. Matrix
multiplication is associative because composition of functions is
associative, and it is non-commutative because applying a rotation
then a shear differs from applying the shear then the rotation. The
entry $(\mat{B}\mat{A})_{ij} = \sum_{k} B_{ik} A_{kj}$ is the
row-times-column dot product, and the cost of multiplying an $m
\times k$ by a $k \times n$ matrix is $2mkn$ flops by the naive
algorithm.

\begin{example}[A rotation as a linear map]
The matrix $\mat{R}_{\theta} = \begin{psmallmatrix} \cos\theta &
-\sin\theta \\ \sin\theta & \cos\theta \end{psmallmatrix}$ rotates
$\R^{2}$ counterclockwise by $\theta$. Its columns are the images of
the standard basis: $\vect{e}_{1} = (1,0)$ maps to $(\cos\theta,
\sin\theta)$ and $\vect{e}_{2} = (0,1)$ maps to $(-\sin\theta,
\cos\theta)$. For $\theta = 90^{\circ}$ the matrix is
$\begin{psmallmatrix} 0 & -1 \\ 1 & 0 \end{psmallmatrix}$, which
sends $(1,0)$ to $(0,1)$ and $(0,1)$ to $(-1,0)$. The composition of
two rotations is a rotation through the sum of the angles:
$\mat{R}_{\alpha}\mat{R}_{\beta} = \mat{R}_{\alpha+\beta}$, and
expanding the product recovers the angle-addition formulas for sine
and cosine (Chapter 3). The rotation matrix is orthogonal
($\mat{R}_{\theta}^{\transpose}\mat{R}_{\theta} = \mat{I}$) and has
determinant $+1$, the algebraic signatures of a length-preserving,
orientation-preserving map.
\end{example}

\subsection{Rank, null space, and the four subspaces}

The \engterm{rank} of $\mat{A}$, written $\rank \mat{A}$, is the
dimension of $\col(\mat{A})$. Equivalently, it is the maximum number
of linearly independent columns, which equals the maximum number of
linearly independent rows.

The \engterm{null space} $\nul(\mat{A}) \subseteq \R^{n}$ is the set
of vectors $\vect{x}$ with $\mat{A}\vect{x} = \vect{0}$. Its
dimension is the \engterm{nullity}. The fundamental relation, the
\emph{rank-nullity theorem}, is

\[
\rank \mat{A} + \dim \nul(\mat{A}) = n.
\]

The theorem is the formal accounting of degrees of freedom: the $n$
input dimensions split into directions that the matrix uses (rank)
and directions it ignores (nullity). We quote the result; the proof
is one page in any text \cite[ch.~3]{text:strang2016}.

Two further subspaces fill out the picture and recur in least
squares (section 9.6). The \engterm{row space} of $\mat{A}$ is
$\col(\mat{A}^{\transpose}) \subseteq \R^{n}$; the \engterm{left
null space} is $\nul(\mat{A}^{\transpose}) \subseteq \R^{m}$. The
four subspaces stand in two orthogonal pairs: in $\R^{n}$, the row
space and the null space are orthogonal complements; in $\R^{m}$,
the column space and the left null space are orthogonal complements.
Figure~\ref{fig:vol02:ch09:four-subspaces} records the picture.

\begin{figure}[ht]
\centering
\begin{tikzpicture}[x=1cm, y=1cm, font=\small]

% Left panel: R^n with row space and null space
\begin{scope}[shift={(0,0)}]
  \node[anchor=south, font=\small\bfseries] at (2.5, 4.6) {Input side $\R^{n}$};
  \draw[thick] (0,0) ellipse (2.6 and 4.2);
  \draw[thick, engblue] (0.6, 1.0) ellipse (1.6 and 2.4);
  \draw[thick, engred] (-0.4, -1.4) ellipse (1.2 and 1.6);
  \node[engblue] at (0.6, 2.8) {row space};
  \node[engblue, font=\scriptsize] at (0.6, 2.3) {$\col(\mat{A}^{\transpose})$};
  \node[engblue, font=\scriptsize] at (0.6, 0.5) {$\dim = r$};
  \node[engred] at (-0.4, -0.5) {null space};
  \node[engred, font=\scriptsize] at (-0.4, -1.0) {$\nul(\mat{A})$};
  \node[engred, font=\scriptsize] at (-0.4, -2.4) {$\dim = n - r$};
\end{scope}

% Center: the map A
\begin{scope}[shift={(5.2, 0)}]
  \draw[->, ultra thick, enggray] (-1.6, 0.3) -- (1.6, 0.3) node[midway, above] {$\mat{A}\vect{x}$};
  \draw[->, ultra thick, enggray] (1.6, -0.3) -- (-1.6, -0.3) node[midway, below] {$\mat{A}^{\transpose}\vect{y}$};
  \node[font=\scriptsize, anchor=south] at (0, 1.0) {rank $r$};
\end{scope}

% Right panel: R^m with column space and left null space
\begin{scope}[shift={(10.4,0)}]
  \node[anchor=south, font=\small\bfseries] at (-2.5, 4.6) {Output side $\R^{m}$};
  \draw[thick] (0,0) ellipse (2.6 and 4.2);
  \draw[thick, engblue] (-0.6, 1.0) ellipse (1.6 and 2.4);
  \draw[thick, engred] (0.4, -1.4) ellipse (1.2 and 1.6);
  \node[engblue] at (-0.6, 2.8) {column space};
  \node[engblue, font=\scriptsize] at (-0.6, 2.3) {$\col(\mat{A})$};
  \node[engblue, font=\scriptsize] at (-0.6, 0.5) {$\dim = r$};
  \node[engred] at (0.4, -0.5) {left null space};
  \node[engred, font=\scriptsize] at (0.4, -1.0) {$\nul(\mat{A}^{\transpose})$};
  \node[engred, font=\scriptsize] at (0.4, -2.4) {$\dim = m - r$};
\end{scope}

\end{tikzpicture}
\caption[The four fundamental subspaces of a matrix]{The four
fundamental subspaces of $\mat{A} \in \R^{m \times n}$. The row
space and null space partition $\R^{n}$ into orthogonal complements;
the column space and left null space partition $\R^{m}$ into
orthogonal complements. The rank $r$ is the dimension shared by the
row space and the column space, and equals the number of independent
input directions the matrix uses.}
\label{fig:vol02:ch09:four-subspaces}
\end{figure}


\begin{example}[The four subspaces of a $3 \times 4$ matrix]
Take
\[
\mat{A} = \begin{pmatrix} 1 & 2 & 0 & 1 \\ 2 & 4 & 1 & 4 \\ 0 & 0 & 1
& 2 \end{pmatrix}.
\]
Row reduction subtracts $2\times$ row 1 from row 2, leaving
$(0,0,1,2)$ in row 2, which equals row 3; subtracting gives a zero
third row. The reduced row-echelon form is
\[
\mat{R} = \begin{pmatrix} 1 & 2 & 0 & 1 \\ 0 & 0 & 1 & 2 \\ 0 & 0 & 0
& 0 \end{pmatrix},
\]
with pivots in columns one and three, so $\rank \mat{A} = 2$. The
four subspaces follow.
\emph{Column space} (in $\R^{3}$, dimension 2): the pivot columns of
the original matrix, $\{(1,2,0)^{\transpose}, (0,1,1)^{\transpose}\}$.
\emph{Row space} (in $\R^{4}$, dimension 2): the nonzero rows of
$\mat{R}$, $\{(1,2,0,1), (0,0,1,2)\}$.
\emph{Null space} (in $\R^{4}$, dimension $4 - 2 = 2$): free
variables $x_{2}$ and $x_{4}$. Setting $x_{2} = 1, x_{4} = 0$ gives
$x_{1} = -2, x_{3} = 0$, so $(-2, 1, 0, 0)^{\transpose}$; setting
$x_{2} = 0, x_{4} = 1$ gives $x_{1} = -1, x_{3} = -2$, so $(-1, 0,
-2, 1)^{\transpose}$.
\emph{Left null space} (in $\R^{3}$, dimension $3 - 2 = 1$): vectors
$\vect{y}$ with $\vect{y}^{\transpose}\mat{A} = \vect{0}$. The third
row reduction recorded $\text{row}_{3} = \text{row}_{2} - 2\,
\text{row}_{1}$ of the working matrix, and tracking the elimination
gives $\vect{y} = (2, -1, 1)^{\transpose}$, which the reader verifies
satisfies $\vect{y}^{\transpose}\mat{A} = \vect{0}$.
The rank-nullity accounting holds in both spaces: $2 + 2 = 4 = n$ in
the domain, $2 + 1 = 3 = m$ in the codomain. Each null-space basis
vector is orthogonal to each row-space basis vector, the orthogonal-%
complement relation Figure~\ref{fig:vol02:ch09:four-subspaces}
records.
\end{example}

\subsection{Special matrices}

Four special structures repay recognition. A matrix is
\engterm{diagonal} when all off-diagonal entries are zero; its action
on $\vect{x}$ is componentwise scaling. It is \engterm{symmetric}
when $\mat{A}^{\transpose} = \mat{A}$; symmetric matrices appear
whenever the underlying problem has a quadratic form with no
preferred orientation, including least squares (section 9.6) and
energy methods in Volume III. It is \engterm{orthogonal} when
$\mat{Q}^{\transpose}\mat{Q} = \mat{I}$; orthogonal matrices preserve
length and inner product, and the $\mat{Q}$ in QR decomposition
(section 9.3) is orthogonal in the rectangular sense
$\mat{Q}^{\transpose}\mat{Q} = \mat{I}$ even when $\mat{Q}\mat{Q}^{\transpose}
\neq \mat{I}$. It is \engterm{positive definite} when $\vect{x}^{\transpose}
\mat{A}\vect{x} > 0$ for all nonzero $\vect{x}$; positive-definite
matrices have positive eigenvalues, are invertible, and admit a
Cholesky factorisation. The Gram matrix $\mat{A}^{\transpose}\mat{A}$
of a full-column-rank $\mat{A}$ is positive definite, a fact the
normal equations of section 9.6 rest on.

\begin{example}[Testing positive definiteness by leading minors]
Checking $\vect{x}^{\transpose}\mat{A}\vect{x} > 0$ for every nonzero
$\vect{x}$ is not something the engineer does directly. Sylvester's
criterion reduces the test to the signs of the leading principal minors,
the determinants of the top-left $1 \times 1$, $2 \times 2$, up to $n
\times n$ submatrices: a symmetric matrix is positive definite exactly
when all of them are positive. Take $\mat{A} = \begin{psmallmatrix} 2 &
-1 & 0 \\ -1 & 2 & -1 \\ 0 & -1 & 2 \end{psmallmatrix}$, the
one-dimensional stiffness matrix. The leading minors are $\det[2] = 2 >
0$, $\det\begin{psmallmatrix} 2 & -1 \\ -1 & 2 \end{psmallmatrix} = 3 >
0$, and $\det\mat{A} = 2(4-1) - (-1)(-2) = 6 - 2 = 4 > 0$. All three
are positive, so $\mat{A}$ is positive definite and the Cholesky
factorisation is guaranteed to run. The minors are not computed
separately in practice; the Cholesky algorithm itself is the efficient
form of the test, since its running diagonal radicands are the ratios of
successive leading minors, and a non-positive radicand is a
non-positive minor. Sylvester's criterion is the hand test at small
size; Cholesky is the machine test at any size. Contrast the indefinite
matrix $\begin{psmallmatrix} 1 & 2 \\ 2 & 1 \end{psmallmatrix}$, whose
second minor is $1 - 4 = -3 < 0$: it fails the criterion, and its
quadratic form $x_{1}^{2} + 4x_{1}x_{2} + x_{2}^{2}$ is negative at
$\vect{x} = (1, -1)^{\transpose}$.
\end{example}

\subsection{Block matrices and the Schur complement}

Large engineering matrices arrive partitioned. A coupled
fluid-structure problem stacks a fluid block, a structure block, and
a coupling block; a network with interior and boundary nodes splits
into four blocks; a saddle-point system from constrained optimisation
(Chapter 15) carries a zero block in one corner. Treating a matrix as
a grid of submatrices, rather than a grid of scalars, is the working
algebra for all of these. A $2 \times 2$ block matrix multiplies and
adds by the same rules as a $2 \times 2$ scalar matrix, provided the
block dimensions are conformable:
\[
\begin{pmatrix} \mat{A} & \mat{B} \\ \mat{C} & \mat{D} \end{pmatrix}
\begin{pmatrix} \vect{x} \\ \vect{y} \end{pmatrix}
= \begin{pmatrix} \mat{A}\vect{x} + \mat{B}\vect{y} \\
\mat{C}\vect{x} + \mat{D}\vect{y} \end{pmatrix}.
\]
The order within each block product is fixed (blocks do not commute),
which is the only departure from scalar arithmetic.

The block view earns its place through the \engterm{Schur
complement}. Eliminating the top block of unknowns from the block
system above, when $\mat{A}$ is invertible, leaves a smaller system
in $\vect{y}$ alone. From the first block row, $\vect{x} =
\mat{A}^{-1}(\vect{b}_{1} - \mat{B}\vect{y})$; substituting into the
second block row gives
\[
(\mat{D} - \mat{C}\mat{A}^{-1}\mat{B})\,\vect{y} = \vect{b}_{2} -
\mat{C}\mat{A}^{-1}\vect{b}_{1}.
\]
The matrix $\mat{S} = \mat{D} - \mat{C}\mat{A}^{-1}\mat{B}$ is the
Schur complement of $\mat{A}$ in the full matrix. It is the exact
algebraic statement of block elimination, and it is the engine of
substructuring: a finite-element model partitioned into a repeated
substructure and a connecting interface eliminates the interior
degrees of freedom once per substructure, leaving a small interface
system, the working principle of static condensation in Volume III.
The block determinant follows at once, $\det \begin{psmallmatrix}
\mat{A} & \mat{B} \\ \mat{C} & \mat{D} \end{psmallmatrix} = \det
\mat{A} \cdot \det \mat{S}$, the block generalisation of the scalar
$2 \times 2$ formula.

\begin{example}[Static condensation by Schur complement]
A finite-element model has four interior degrees of freedom and two
interface degrees of freedom, with the symmetric stiffness matrix
partitioned as $\begin{psmallmatrix} \mat{K}_{ii} & \mat{K}_{ib} \\
\mat{K}_{bi} & \mat{K}_{bb} \end{psmallmatrix}$, where the subscript
$i$ marks interior and $b$ marks boundary (interface). Take the small
illustrative case
\[
\mat{K}_{ii} = \begin{pmatrix} 4 & -1 \\ -1 & 4 \end{pmatrix}, \quad
\mat{K}_{ib} = \begin{pmatrix} -1 & 0 \\ 0 & -1 \end{pmatrix}, \quad
\mat{K}_{bb} = \begin{pmatrix} 3 & -1 \\ -1 & 3 \end{pmatrix},
\]
with $\mat{K}_{bi} = \mat{K}_{ib}^{\transpose}$ by symmetry. The
condensed interface stiffness is the Schur complement $\mat{S} =
\mat{K}_{bb} - \mat{K}_{bi}\mat{K}_{ii}^{-1}\mat{K}_{ib}$. The
interior inverse is $\mat{K}_{ii}^{-1} = \tfrac{1}{15}
\begin{psmallmatrix} 4 & 1 \\ 1 & 4 \end{psmallmatrix}$, so
$\mat{K}_{bi}\mat{K}_{ii}^{-1}\mat{K}_{ib} = \tfrac{1}{15}
\begin{psmallmatrix} 4 & 1 \\ 1 & 4 \end{psmallmatrix}$ and
\[
\mat{S} = \begin{pmatrix} 3 & -1 \\ -1 & 3 \end{pmatrix}
- \frac{1}{15}\begin{pmatrix} 4 & 1 \\ 1 & 4 \end{pmatrix}
= \begin{pmatrix} 2.733 & -1.067 \\ -1.067 & 2.733 \end{pmatrix}.
\]
The interface now carries the interior's elastic response through a
$2 \times 2$ matrix, and the full model never re-touches the interior
unknowns until the interface displacements are known. The condensed
matrix is again symmetric positive definite, the property that lets a
substructured solver chain Schur complements up a tree of
substructures.
\end{example}

\begin{example}[The inverse of a block matrix through the Schur complement]
The Schur complement also assembles the inverse of a block matrix in
closed form, which the engineer meets whenever a covariance matrix or a
saddle-point system must be inverted block by block. With $\mat{S} =
\mat{D} - \mat{C}\mat{A}^{-1}\mat{B}$ the Schur complement of the
invertible top-left block, the inverse is
\[
\begin{pmatrix} \mat{A} & \mat{B} \\ \mat{C} & \mat{D} \end{pmatrix}^{-1}
= \begin{pmatrix}
\mat{A}^{-1} + \mat{A}^{-1}\mat{B}\mat{S}^{-1}\mat{C}\mat{A}^{-1} &
-\mat{A}^{-1}\mat{B}\mat{S}^{-1} \\
-\mat{S}^{-1}\mat{C}\mat{A}^{-1} & \mat{S}^{-1}
\end{pmatrix}.
\]
The formula turns one $(n+m) \times (n+m)$ inversion into an $n \times
n$ inversion (of $\mat{A}$) and an $m \times m$ inversion (of $\mat{S}$),
the block analogue of never inverting more than necessary. Take the
symmetric block matrix with $\mat{A} = \begin{psmallmatrix} 2 & 0 \\ 0 &
2 \end{psmallmatrix}$, $\mat{B} = \mat{C}^{\transpose} =
\begin{psmallmatrix} 1 \\ 0 \end{psmallmatrix}$, and $\mat{D} = [3]$.
Then $\mat{A}^{-1} = \tfrac{1}{2}\mat{I}$, $\mat{C}\mat{A}^{-1}\mat{B} =
(1, 0)\tfrac{1}{2}\mat{I}(1, 0)^{\transpose} = \tfrac{1}{2}$, so
$\mat{S} = 3 - \tfrac{1}{2} = \tfrac{5}{2}$ and $\mat{S}^{-1} =
\tfrac{2}{5}$. The bottom-right block of the inverse is $\tfrac{2}{5}$,
which the reader verifies against a direct $3 \times 3$ inversion. The
same block inverse, applied to a Gaussian covariance matrix, is the
algebra behind conditional distributions in Volume IX: the Schur
complement of a covariance block is the conditional covariance, and the
off-diagonal blocks give the conditional mean's regression coefficients.
\end{example}

\subsection{Kronecker products and structured grids}

A second structure recurs whenever a problem lives on a tensor-product
grid: a rectangular finite-difference mesh, a separable image filter,
a two-factor experimental design. The \engterm{Kronecker product}
$\mat{A} \otimes \mat{B}$ of $\mat{A} \in \R^{m \times n}$ and
$\mat{B} \in \R^{p \times q}$ is the $mp \times nq$ block matrix
whose $(i,j)$ block is $A_{ij}\mat{B}$. Its single most useful
identity rewrites a matrix equation as a matrix-vector product:
\[
(\mat{B}^{\transpose} \otimes \mat{A})\,\operatorname{vec}(\mat{X}) =
\operatorname{vec}(\mat{A}\mat{X}\mat{B}),
\]
where $\operatorname{vec}(\mat{X})$ stacks the columns of $\mat{X}$ into one long
vector. The two-dimensional Laplacian on an $N \times N$ grid is the
canonical use: the discrete operator is $\mat{L}_{2\mathrm{D}} =
\mat{L}_{1\mathrm{D}} \otimes \mat{I} + \mat{I} \otimes
\mat{L}_{1\mathrm{D}}$, built from the one-dimensional second-difference
matrix of the determinant example in section 9.4. The eigenvalues of
the two-dimensional operator are sums of one-dimensional eigenvalues,
$\lambda_{jk} = \mu_{j} + \mu_{k}$, so the conditioning of the
$N^{2} \times N^{2}$ grid Laplacian is read off the $N \times N$
one-dimensional spectrum at no extra cost. The Kronecker form is the
algebraic reason fast Poisson solvers exist: the eigenvectors of a
Kronecker sum factor across the two grid directions, and a separable
problem decouples into independent one-dimensional solves
\cite[ch.~4]{text:v2ch09-golub-vanloan2013}.

\section{Solving linear systems: Gaussian elimination, LU, QR}\pathtag{core}

A square linear system $\mat{A}\vect{x} = \vect{b}$ with $\mat{A}
\in \R^{n \times n}$ has a unique solution when $\mat{A}$ is
invertible, no solutions when $\vect{b} \notin \col(\mat{A})$, and
infinitely many when $\rank \mat{A} < n$ and $\vect{b} \in
\col(\mat{A})$. The classification is exact. The numerical
question, taken up in section 9.7, is whether the unique solution
the algebra promises is the solution the floating-point arithmetic
returns.

\subsection{Gaussian elimination}

\engterm{Gaussian elimination} reduces $\mat{A}$ to upper-triangular
form by adding multiples of one row to another. The standard
algorithm processes columns left to right; in column $j$, the
pivot is $A_{jj}$ (after any previous updates), and rows below row
$j$ are modified to zero out their entries in column $j$. The
operation count is $\tfrac{2}{3}n^{3} + O(n^{2})$ flops for the
elimination phase and $O(n^{2})$ for the back substitution
\cite[ch.~20]{text:trefethen-bau1997}. For $n = 100$ the elimination
takes about $6.7 \times 10^{5}$ flops, a millisecond on any working
machine; for $n = 10^{6}$ the count climbs to $\oom{10^{18}}$ flops,
which is years on a single core. The cubic scaling is the central
fact: doubling the problem size multiplies the work by eight.

When a pivot $A_{jj}$ is zero, or small relative to the entries
below it, the algorithm uses \engterm{partial pivoting}: swap row
$j$ with a row $k > j$ whose entry $A_{kj}$ is largest in absolute
value, then proceed. Partial pivoting is the default in any
production solver. It does not change the complexity but it
controls the growth factor in the elimination, which controls the
floating-point error.

\begin{example}[Why pivoting matters in finite precision]
The system $\begin{psmallmatrix} 10^{-4} & 1 \\ 1 & 1
\end{psmallmatrix}\vect{x} = \begin{psmallmatrix} 1 \\ 2
\end{psmallmatrix}$ has the exact solution $\vect{x} \approx
(1.0001, 0.9999)^{\transpose}$. Without pivoting, the multiplier is
$1/10^{-4} = 10^{4}$, and the $(2,2)$ entry becomes $1 - 10^{4}
\cdot 1 = -9999$, while the right-hand side becomes $2 - 10^{4}
\cdot 1 = -9998$. In three-digit arithmetic, $1 - 10^{4}$ rounds to
$-10^{4}$ and $2 - 10^{4}$ rounds to $-10^{4}$, so back substitution
returns $x_{2} = 1$ and then $x_{1} = (1 - 1)/10^{-4} = 0$, a
catastrophic error in $x_{1}$. With partial pivoting the rows swap
first; the multiplier is now $10^{-4}$, the $(2,2)$ entry becomes
$1 - 10^{-4} \approx 1$, and three-digit arithmetic returns $x_{1}
\approx 1.00$, $x_{2} \approx 1.00$, correct to the precision
carried. The small pivot, not the conditioning of the matrix
($\kappa \approx 2.6$, well conditioned), is what destroyed the
unpivoted solve. Partial pivoting is cheap insurance against this
elementary failure, and it is on by default in every production
solver for exactly this reason.
\end{example}

\subsection{LU factorisation}

The work of Gaussian elimination factors into two parts: a unit
lower-triangular matrix $\mat{L}$ (whose entries record the row
multipliers used) and an upper-triangular matrix $\mat{U}$ (whose
rows are the elimination's final rows). The relation is

\[
\mat{P}\mat{A} = \mat{L}\mat{U},
\]

where $\mat{P}$ is the permutation matrix that records the row
swaps. The factorisation costs the same $\tfrac{2}{3}n^{3}$ flops
as the elimination itself; it pays for itself the second time a
system with the same $\mat{A}$ but a different $\vect{b}$ is
solved.

Once $\mat{P}$, $\mat{L}$, and $\mat{U}$ are stored, solving
$\mat{A}\vect{x} = \vect{b}$ becomes three cheap steps. Permute
$\vect{b}$ to $\mat{P}\vect{b}$. Solve $\mat{L}\vect{y} =
\mat{P}\vect{b}$ by forward substitution ($O(n^{2})$ flops). Solve
$\mat{U}\vect{x} = \vect{y}$ by back substitution ($O(n^{2})$
flops). The reuse pattern is the right discipline whenever the
engineer faces many right-hand sides on the same matrix, as in
finite-element analysis with multiple load cases.

\subsection{QR factorisation}

A second factorisation, $\mat{A} = \mat{Q}\mat{R}$, writes a
full-column-rank matrix $\mat{A} \in \R^{m \times n}$ ($m \geq n$)
as the product of an orthogonal $\mat{Q} \in \R^{m \times n}$
(satisfying $\mat{Q}^{\transpose}\mat{Q} = \mat{I}_{n}$) and an
upper-triangular $\mat{R} \in \R^{n \times n}$. Three algorithms
compute QR. \engterm{Classical Gram-Schmidt} orthogonalises columns
of $\mat{A}$ one at a time; it is unstable in floating point and
should not be used in production. \engterm{Modified Gram-Schmidt}
orthogonalises in a different order and is markedly better, though
not best. \engterm{Householder reflections} reduce $\mat{A}$ to
upper-triangular form by a sequence of orthogonal reflections; the
algorithm is backward stable and is the standard production choice
\cite[ch.~10]{text:trefethen-bau1997}.

The Householder QR costs $2mn^{2} - \tfrac{2}{3}n^{3}$ flops, which
for square matrices ($m = n$) is $\tfrac{4}{3}n^{3}$, twice the cost
of LU. The factor of two is the price for orthogonality, and the
price is well spent in least-squares problems (section 9.6) and in
ill-conditioned systems (section 9.7).

\begin{mastery}
The leading-order flop counts for the standard dense factorisations
of an $n \times n$ matrix, with the structure each one assumes, are
worth carrying as a single table.
\begin{center}
\begin{tabular}{lll}
\toprule
Factorisation & Flops (leading) & Assumed structure \\
\midrule
LU (partial pivot) & $\tfrac{2}{3}n^{3}$ & general square \\
Cholesky & $\tfrac{1}{3}n^{3}$ & symmetric positive definite \\
QR (Householder) & $\tfrac{4}{3}n^{3}$ & general, $m \geq n$ \\
SVD & $\approx 22 n^{3}$ & general, full spectrum \\
Forward/back solve & $2n^{2}$ & triangular \\
\bottomrule
\end{tabular}
\end{center}
Cholesky is half of LU because the symmetry halves the work; it is
the production choice for the symmetric positive-definite systems
that energy methods (Volume III) and least squares (section 9.6)
produce. The SVD constant varies by algorithm and by whether the
singular vectors are requested; the $22 n^{3}$ figure is a working
upper estimate for the full decomposition
\cite[ch.~31]{text:trefethen-bau1997}. The triangular-solve count is
the reason LU reuse is so cheap once the factorisation is in hand.
\end{mastery}

\begin{example}[The economics of LU reuse]
A finite-element model with $n = 2000$ degrees of freedom is solved
under $50$ load cases, all sharing the same stiffness matrix
$\mat{K}$. Solving each case from scratch by Gaussian elimination
costs $50 \times \tfrac{2}{3}(2000)^{3} \approx 2.7 \times 10^{11}$
flops. Factoring $\mat{K} = \mat{L}\mat{U}$ once and reusing the
factors costs $\tfrac{2}{3}(2000)^{3} \approx 5.3 \times 10^{9}$
flops for the factorisation plus $50 \times 2(2000)^{2} = 4 \times
10^{8}$ flops for the fifty triangular solves, a total of $5.7
\times 10^{9}$ flops. The reuse is roughly $47$ times cheaper, and
the saving grows with the number of right-hand sides. The discipline
generalises: whenever the matrix is fixed and only the right-hand
side changes, factor once and back-solve many times.
\end{example}

\subsection{Cholesky factorisation}

When the matrix is symmetric positive definite, the factorisation
specialises. \engterm{Cholesky factorisation} writes such a matrix as
$\mat{A} = \mat{L}\mat{L}^{\transpose}$, with $\mat{L}$ lower
triangular and positive diagonal. The symmetry halves the work to
$\tfrac{1}{3}n^{3}$ flops, no pivoting is needed (positive
definiteness guarantees positive pivots), and the algorithm doubles
as a definiteness test: it runs to completion exactly when the matrix
is positive definite, and it fails at the first non-positive pivot
otherwise. Symmetric positive-definite matrices are the common case
in engineering, arising from every least-squares Gram matrix
(section 9.6), every elastic stiffness matrix (Volume III), and every
grounded network Laplacian (the resistor case above), so Cholesky is
the factorisation an engineer reaches for first whenever symmetry and
positive definiteness are in hand.

\begin{example}[Cholesky by hand]
Factor $\mat{A} = \begin{psmallmatrix} 4 & 2 & 2 \\ 2 & 5 & 3 \\ 2 &
3 & 6 \end{psmallmatrix}$. The algorithm fills $\mat{L}$ column by
column. The first diagonal is $L_{11} = \sqrt{4} = 2$; the rest of
column one is $L_{21} = 2/2 = 1$ and $L_{31} = 2/2 = 1$. The second
diagonal subtracts the squared entries already placed in its row:
$L_{22} = \sqrt{5 - 1^{2}} = 2$, and $L_{32} = (3 - L_{31}L_{21})/
L_{22} = (3 - 1)/2 = 1$. The third diagonal is $L_{33} = \sqrt{6 -
L_{31}^{2} - L_{32}^{2}} = \sqrt{6 - 1 - 1} = 2$. The factor is
\[
\mat{L} = \begin{pmatrix} 2 & 0 & 0 \\ 1 & 2 & 0 \\ 1 & 1 & 2
\end{pmatrix}, \qquad \mat{L}\mat{L}^{\transpose} = \mat{A}.
\]
Every diagonal square root took the root of a positive number, which
is the running certificate that $\mat{A}$ is positive definite. Had
any radicand come out non-positive, the matrix would have been
indefinite and the factorisation would have stopped, the cheapest
positive-definiteness test available. Solving $\mat{A}\vect{x} =
\vect{b}$ then forward-solves $\mat{L}\vect{y} = \vect{b}$ and
back-solves $\mat{L}^{\transpose}\vect{x} = \vect{y}$, two triangular
solves on the one factor.
\end{example}

\begin{example}[Iterative refinement recovers lost digits]
A direct solve in finite precision returns $\hat{\vect{x}}$ with a
residual $\vect{r} = \vect{b} - \mat{A}\hat{\vect{x}}$ that is small
but nonzero. The residual carries information: the error $\vect{e} =
\vect{x} - \hat{\vect{x}}$ satisfies $\mat{A}\vect{e} = \vect{r}$
exactly, the same system with a new right-hand side. \engterm{Iterative
refinement} computes $\vect{r}$ (ideally in higher precision), solves
$\mat{A}\vect{\delta} = \vect{r}$ by reusing the existing LU or
Cholesky factors at $O(n^{2})$ cost, and updates $\hat{\vect{x}}
\leftarrow \hat{\vect{x}} + \vect{\delta}$. For a matrix with
$\kappa(\mat{A}) = 10^{8}$, a first solve in double precision yields
about eight correct digits; one refinement step with the residual
formed in extended precision typically recovers a further several
digits at a cost far below a fresh factorisation. The technique is
the cheap insurance for a moderately ill-conditioned solve, and it is
the reason LU and Cholesky reuse (the previous subsection) pays a
second dividend: the factors that solve the system also drive the
refinement that sharpens it.
\end{example}

\subsection{Worked example}

Consider the system

\[
\mat{A}\vect{x} = \vect{b}, \quad
\mat{A} = \begin{pmatrix} 2 & 1 & 1 \\ 4 & 3 & 3 \\ 8 & 7 & 9
\end{pmatrix}, \quad
\vect{b} = \begin{pmatrix} 4 \\ 10 \\ 30 \end{pmatrix}.
\]

Gaussian elimination subtracts $2 \cdot R_{1}$ from $R_{2}$ and
$4 \cdot R_{1}$ from $R_{3}$, yielding

\[
\begin{pmatrix} 2 & 1 & 1 \\ 0 & 1 & 1 \\ 0 & 3 & 5 \end{pmatrix}
\vect{x} = \begin{pmatrix} 4 \\ 2 \\ 14 \end{pmatrix}.
\]

A second pass subtracts $3 \cdot R_{2}$ from $R_{3}$:

\[
\begin{pmatrix} 2 & 1 & 1 \\ 0 & 1 & 1 \\ 0 & 0 & 2 \end{pmatrix}
\vect{x} = \begin{pmatrix} 4 \\ 2 \\ 8 \end{pmatrix}.
\]

Back substitution gives $x_{3} = 4$, $x_{2} = 2 - 4 = -2$, $x_{1} =
(4 - (-2) - 4)/2 = 1$. The solution is $\vect{x} = (1, -2, 4)$, and
the LU factors recorded by the elimination are

\[
\mat{L} = \begin{pmatrix} 1 & 0 & 0 \\ 2 & 1 & 0 \\ 4 & 3 & 1
\end{pmatrix}, \quad
\mat{U} = \begin{pmatrix} 2 & 1 & 1 \\ 0 & 1 & 1 \\ 0 & 0 & 2
\end{pmatrix}.
\]

The reader who multiplies $\mat{L}\mat{U}$ recovers $\mat{A}$.
Figure~\ref{fig:vol02:ch09:gaussian-elimination} stages the
elimination and records the multipliers.

\begin{figure}[ht]
\centering
\begin{tikzpicture}[x=1cm, y=1cm, font=\small]

% Stage 0: original A
\begin{scope}[shift={(0,0)}]
  \node[anchor=south, font=\small\bfseries] at (1.4, 2.6) {Stage 0};
  \node[anchor=south, font=\scriptsize] at (1.4, 2.2) {original};
  \matrix[matrix of math nodes, left delimiter={(}, right delimiter={)},
          column sep=4pt, row sep=2pt] at (1.4, 1.0) {
    2 & 1 & 1 \\
    4 & 3 & 3 \\
    8 & 7 & 9 \\
  };
\end{scope}

% arrow
\draw[->, thick] (3.2, 1.0) -- (3.8, 1.0) node[midway, above, font=\scriptsize] {pivot $A_{11}$};

% Stage 1
\begin{scope}[shift={(4.0,0)}]
  \node[anchor=south, font=\small\bfseries] at (1.4, 2.6) {Stage 1};
  \node[anchor=south, font=\scriptsize] at (1.4, 2.2) {column 1 zeroed};
  \matrix[matrix of math nodes, left delimiter={(}, right delimiter={)},
          column sep=4pt, row sep=2pt] at (1.4, 1.0) {
    2 & 1 & 1 \\
    |[engblue]| 0 & 1 & 1 \\
    |[engblue]| 0 & 3 & 5 \\
  };
\end{scope}

\draw[->, thick] (7.2, 1.0) -- (7.8, 1.0) node[midway, above, font=\scriptsize] {pivot $A_{22}$};

% Stage 2: U
\begin{scope}[shift={(8.0,0)}]
  \node[anchor=south, font=\small\bfseries] at (1.4, 2.6) {Stage 2};
  \node[anchor=south, font=\scriptsize] at (1.4, 2.2) {$\mat{U}$ ready};
  \matrix[matrix of math nodes, left delimiter={(}, right delimiter={)},
          column sep=4pt, row sep=2pt] at (1.4, 1.0) {
    2 & 1 & 1 \\
    0 & 1 & 1 \\
    0 & |[engblue]| 0 & 2 \\
  };
\end{scope}

% Multipliers stored as L below
\begin{scope}[shift={(2.0, -3.2)}]
  \node[anchor=south, font=\small\bfseries] at (3.0, 1.4) {Multipliers recorded as $\mat{L}$};
  \matrix[matrix of math nodes, left delimiter={(}, right delimiter={)},
          column sep=4pt, row sep=2pt] at (3.0, 0.2) {
    1 & 0 & 0 \\
    |[engred]| 2 & 1 & 0 \\
    |[engred]| 4 & |[engred]| 3 & 1 \\
  };
  \node[anchor=west, font=\scriptsize, text width=4.6cm] at (4.8, 0.2) {
    Row 2: $R_{2} - 2 R_{1}$. \\
    Row 3: $R_{3} - 4 R_{1}$, then $R_{3} - 3 R_{2}$. \\
    Diagonal of $\mat{L}$ is unit; entries below diagonal store the multipliers.
  };
\end{scope}

\end{tikzpicture}
\caption[Two-stage Gaussian elimination on a $3 \times 3$ system]{Gaussian
elimination on the worked example of section 9.3. Stage 1 uses the
pivot $A_{11} = 2$ to zero entries below; stage 2 uses the new pivot
$A_{22} = 1$ to zero the entry below it. The resulting upper-triangular
matrix is $\mat{U}$; the row multipliers used at each step, recorded in
the lower triangle with unit diagonal, give $\mat{L}$, and
$\mat{L}\mat{U} = \mat{A}$. For the back substitution the entries of
$\vect{b}$ are transformed in lockstep.}
\label{fig:vol02:ch09:gaussian-elimination}
\end{figure}


The same example is the input to the LU demonstration code in
\path{code/lu_demo.py}, which solves it by hand-coded Doolittle
elimination and verifies $\mat{L}\mat{U} = \mat{A}$ and
$\mat{A}\vect{x} = \vect{b}$ to machine precision.

\begin{example}[Elimination as a product of elementary matrices]
The elimination of the worked example can be read as left-multiplication
by a sequence of elementary matrices, which is where the $\mat{L}$ of
$\mat{P}\mat{A} = \mat{L}\mat{U}$ comes from. Each row operation
``subtract $m$ times row $j$ from row $i$'' is the matrix $\mat{E}_{ij}(m)
= \mat{I} - m\,\vect{e}_{i}\vect{e}_{j}^{\transpose}$, the identity with
a single $-m$ off the diagonal. The first pass of the worked example
applies $\mat{E}_{21}(2)$ and $\mat{E}_{31}(4)$; the second applies
$\mat{E}_{32}(3)$. Their product acting on $\mat{A}$ produces $\mat{U}$:
\[
\mat{E}_{32}(3)\,\mat{E}_{31}(4)\,\mat{E}_{21}(2)\,\mat{A} = \mat{U}.
\]
The inverse of each elementary matrix simply flips the sign of the
multiplier, $\mat{E}_{ij}(m)^{-1} = \mat{E}_{ij}(-m)$, and the product
of the inverses, taken in reverse order, is $\mat{L}$:
\[
\mat{L} = \mat{E}_{21}(2)^{-1}\mat{E}_{31}(4)^{-1}\mat{E}_{32}(3)^{-1}
= \begin{pmatrix} 1 & 0 & 0 \\ 2 & 1 & 0 \\ 4 & 3 & 1 \end{pmatrix}.
\]
The multipliers land in $\mat{L}$ with their original sign and in their
original positions because the elementary matrices for a single
elimination pass commute in exactly the way that keeps the bookkeeping
trivial. This is the reason $\mat{L}$ can be assembled in place, during
the elimination, from the multipliers already being computed: no
separate matrix products are ever formed. A row swap is a fourth kind
of elementary matrix, a permutation, and collecting the swaps into
$\mat{P}$ is what turns the factorisation into its pivoted form
$\mat{P}\mat{A} = \mat{L}\mat{U}$.
\end{example}

\begin{example}[A factorisation that requires a swap]
The worked example above needed no row swap, so its factorisation was
the unpivoted $\mat{A} = \mat{L}\mat{U}$. A matrix with a zero in the
first pivot position forces the permutation into the open. Take
\[
\mat{A} = \begin{pmatrix} 0 & 2 & 1 \\ 1 & 1 & 0 \\ 2 & 1 & 3
\end{pmatrix}.
\]
The $(1,1)$ entry is zero, so column-one elimination cannot begin;
partial pivoting swaps row $1$ with the row carrying the largest
first-column magnitude, which is row $3$ with entry $2$. The
permutation is $\mat{P} = \begin{psmallmatrix} 0 & 0 & 1 \\ 0 & 1 & 0
\\ 1 & 0 & 0 \end{psmallmatrix}$, and
\[
\mat{P}\mat{A} = \begin{pmatrix} 2 & 1 & 3 \\ 1 & 1 & 0 \\ 0 & 2 & 1
\end{pmatrix}.
\]
The elimination proceeds on the permuted matrix. Subtract
$\tfrac{1}{2}$ row $1$ from row $2$: the
multiplier is $L_{21} = 0.5$ and row $2$ becomes $(0, 0.5, -1.5)$. Row
$3$ already has a zero in column one, so $L_{31} = 0$. For column two,
the pivot is $0.5$ and the entry below it is $2$, so the multiplier is
$L_{32} = 2/0.5 = 4$; subtracting $4$ times the new row $2$ from row
$3$ gives $(0, 0, 1 - 4(-1.5)) = (0, 0, 7)$. The factors are
\[
\mat{L} = \begin{pmatrix} 1 & 0 & 0 \\ 0.5 & 1 & 0 \\ 0 & 4 & 1
\end{pmatrix}, \qquad
\mat{U} = \begin{pmatrix} 2 & 1 & 3 \\ 0 & 0.5 & -1.5 \\ 0 & 0 & 7
\end{pmatrix},
\]
and the reader multiplies $\mat{L}\mat{U}$ to recover $\mat{P}\mat{A}$,
not $\mat{A}$: the factorisation is of the permuted matrix, which is why
solving $\mat{A}\vect{x} = \vect{b}$ through these factors requires
permuting the right-hand side to $\mat{P}\vect{b}$ first. The
determinant follows with a sign for the single swap, $\det\mat{A} =
(-1)^{1}\prod U_{ii} = -(2)(0.5)(7) = -7$, and a direct cofactor
expansion of $\mat{A}$ confirms the value. Figure~\ref{fig:vol02:ch09:pivot-permutation}
stages the swap and the elimination.
\end{example}

\begin{figure}[ht]
\centering
\begin{tikzpicture}[x=1.0cm, y=1.0cm, font=\small, >=Stealth]

% Stage labels along the top
\node[font=\scriptsize] at (1.5, 3.4) {original $\mat{A}$};
\node[font=\scriptsize] at (5.5, 3.4) {after row swap $\mat{P}\mat{A}$};
\node[font=\scriptsize] at (9.5, 3.4) {eliminated $\mat{U}$};

% Stage 1: original, zero (small) pivot flagged
\node at (1.5, 2.0) {$\begin{pmatrix} 0 & 2 & 1 \\ 1 & 1 & 0 \\ 2 & 1 & 3 \end{pmatrix}$};
\node[engred, font=\scriptsize, anchor=north] at (1.5, 0.8) {pivot $=0$: swap needed};

% arrow
\draw[->, thick, enggray] (2.7, 2.0) -- (3.9, 2.0);

% Stage 2: after swapping row 1 with row 3 (largest |entry|)
\node at (5.5, 2.0) {$\begin{pmatrix} 2 & 1 & 3 \\ 1 & 1 & 0 \\ 0 & 2 & 1 \end{pmatrix}$};
\node[engblue, font=\scriptsize, anchor=north] at (5.5, 0.8) {largest entry to pivot};

% arrow
\draw[->, thick, enggray] (7.1, 2.0) -- (8.3, 2.0);

% Stage 3: upper triangular
\node at (9.5, 2.0) {$\begin{pmatrix} 2 & 1 & 3 \\ 0 & 0.5 & -1.5 \\ 0 & 0 & 7 \end{pmatrix}$};
\node[engamber, font=\scriptsize, anchor=north] at (9.5, 0.8) {$\mat{P}\mat{A} = \mat{L}\mat{U}$};

\node[anchor=north, font=\scriptsize, text width=12cm, align=center]
  at (5.5, 0.1) {
  A zero (or small) pivot forces a row exchange before elimination can
  proceed. The permutation matrix $\mat{P}$ records every swap, so the
  factorisation the algorithm actually produces is of the permuted
  matrix, $\mat{P}\mat{A} = \mat{L}\mat{U}$.
};

\end{tikzpicture}
\caption[Partial pivoting produces the permuted factorisation]{Partial
pivoting swaps the row with the largest available entry into the pivot
position before each elimination step, both to avoid a zero pivot and
to control the growth of round-off. The swaps are collected into a
permutation matrix $\mat{P}$, and the factorisation the algorithm
delivers is of the permuted matrix, $\mat{P}\mat{A} = \mat{L}\mat{U}$.
Applying the same permutation to the right-hand side, then solving
$\mat{L}\vect{y} = \mat{P}\vect{b}$ and $\mat{U}\vect{x} = \vect{y}$,
returns the solution of the original system.}
\label{fig:vol02:ch09:pivot-permutation}
\end{figure}


\subsection{The matrix inverse and when not to use it}

The inverse $\mat{A}^{-1}$ of an invertible square matrix satisfies
$\mat{A}\mat{A}^{-1} = \mat{A}^{-1}\mat{A} = \mat{I}$. The inverse
exists, it is unique, and the solution of $\mat{A}\vect{x} =
\vect{b}$ is $\vect{x} = \mat{A}^{-1}\vect{b}$ in exact arithmetic.
The existence of the formula is a frequent source of a bad habit:
solving a linear system by forming $\mat{A}^{-1}$ and multiplying.
For the $2 \times 2$ case the closed form is worth knowing,
\[
\begin{pmatrix} a & b \\ c & d \end{pmatrix}^{-1} = \frac{1}{ad - bc}
\begin{pmatrix} d & -b \\ -c & a \end{pmatrix},
\]
and for hand computation on tiny systems it is convenient. At any
scale, forming the inverse is the wrong computational choice for
three reasons. It costs more: computing $\mat{A}^{-1}$ requires
factoring $\mat{A}$ and then $n$ back-solves to recover the inverse's
columns, roughly three times the work of a single direct solve, and
the subsequent matrix-vector multiply adds $2n^{2}$ more. It is less
accurate: the explicit inverse carries the conditioning of $\mat{A}$
twice, once in its formation and once in the multiply, where a
direct solve carries it once. It destroys structure: the inverse of
a sparse or banded matrix is generally dense, so a $10^{6} \times
10^{6}$ sparse system with a few nonzeros per row produces an
inverse that does not fit in memory. The discipline is the same as
the LU-reuse discipline of the previous subsection: factor once,
solve directly, and never form $\mat{A}^{-1}$ unless the inverse
itself, not the solution of a system, is the required output. The
rare cases that genuinely need the inverse (a variance-covariance
matrix in statistics, a small analytic Jacobian inverse) are
exactly those where $\mat{A}^{-1}$ is the deliverable.

\subsection{Rank-one updates: the Sherman-Morrison formula}

A matrix in an engineering computation is rarely solved once and left
alone. A sensor is recalibrated and one column of a design matrix
shifts; a single spring in a network is stiffened and one entry of the
stiffness matrix changes; an optimiser takes a step and the Hessian
receives a rank-one correction. In each case the new matrix is the old
one plus a rank-one term, $\mat{A} + \vect{u}\vect{v}^{\transpose}$,
and refactoring from scratch discards the work already invested in the
factors of $\mat{A}$. The \engterm{Sherman-Morrison formula} corrects
the inverse instead of recomputing it:
\[
(\mat{A} + \vect{u}\vect{v}^{\transpose})^{-1} = \mat{A}^{-1} -
\frac{\mat{A}^{-1}\vect{u}\,\vect{v}^{\transpose}\mat{A}^{-1}}
{1 + \vect{v}^{\transpose}\mat{A}^{-1}\vect{u}},
\]
valid whenever the scalar denominator $1 + \vect{v}^{\transpose}
\mat{A}^{-1}\vect{u}$ is nonzero, which is exactly the condition that
the updated matrix stays invertible. The formula is never applied by
forming inverses. Its computational content is that the updated solve
$(\mat{A} + \vect{u}\vect{v}^{\transpose})\vect{x} = \vect{b}$ reuses
the existing factors of $\mat{A}$ through two triangular solves, one
for $\vect{y} = \mat{A}^{-1}\vect{b}$ (already in hand if $\mat{A}
\vect{x} = \vect{b}$ was solved) and one for $\vect{z} = \mat{A}^{-1}
\vect{u}$, and then closes with an $\oom{n}$ correction. Each
single-column edit costs $\oom{n^{2}}$ rather than the $\oom{n^{3}}$
of a fresh factorisation, the same reuse discipline the LU section
applied across right-hand sides, now applied across small changes to
the matrix itself. Figure~\ref{fig:vol02:ch09:sherman-morrison}
contrasts the two routes.

\begin{figure}[ht]
\centering
\begin{tikzpicture}[x=1.0cm, y=1.0cm, font=\small,
  box/.style={draw, thick, minimum height=8mm, inner sep=3pt},
  op/.style={font=\scriptsize, enggray}]

% Top path: naive refactor
\node[box, fill=engred!10] (a1) at (0, 2.6) {$\mat{A} + \vect{u}\vect{v}^{\transpose}$};
\node[box, fill=engred!10] (a2) at (4.2, 2.6) {new LU factors};
\node[box, fill=engred!10] (a3) at (8.6, 2.6) {solve};
\draw[->, thick, engred, >=stealth] (a1) -- (a2)
  node[midway, above, op] {$\tfrac{2}{3}n^{3}$ flops};
\draw[->, thick, engred, >=stealth] (a2) -- (a3)
  node[midway, above, op] {$2n^{2}$};
\node[engred, anchor=west, font=\scriptsize] at (10.0, 2.6) {refactor};

% Bottom path: Sherman-Morrison
\node[box, fill=engblue!10] (b1) at (0, 0) {reuse $\mat{A}^{-1}\vect{b}$};
\node[box, fill=engblue!10] (b2) at (4.2, 0) {one extra solve $\mat{A}^{-1}\vect{u}$};
\node[box, fill=engblue!10] (b3) at (8.6, 0) {rank-one update};
\draw[->, thick, engblue, >=stealth] (b1) -- (b2)
  node[midway, above, op] {$2n^{2}$};
\draw[->, thick, engblue, >=stealth] (b2) -- (b3)
  node[midway, above, op] {$\oom{n}$};
\node[engblue, anchor=west, font=\scriptsize] at (10.0, 0) {update};

% brace comparing
\draw[decorate, decoration={brace, amplitude=5pt}, enggray]
  (-1.4, 2.9) -- (-1.4, -0.3);
\node[enggray, font=\scriptsize, rotate=90, anchor=south] at (-1.9, 1.3)
  {same $\mat{A}$, one column changed};

\node[anchor=north, font=\scriptsize, text width=11cm, align=center]
  at (5.0, -1.0) {
  A rank-one change to the matrix does not require a fresh
  factorisation. The Sherman--Morrison identity reuses the existing
  factors and corrects the solution with one auxiliary solve and an
  $\oom{n}$ update, so a sequence of single-column edits costs
  $\oom{n^{2}}$ each rather than $\oom{n^{3}}$.
};

\end{tikzpicture}
\caption[Sherman--Morrison update against refactorisation]{Two routes
to the solution of $(\mat{A} + \vect{u}\vect{v}^{\transpose})\vect{x}
= \vect{b}$ after a rank-one change to a matrix already factored. The
naive route (top) throws away the factors and pays the cubic cost
again. The Sherman--Morrison route (bottom) keeps the factors, spends
one extra triangular solve to form $\mat{A}^{-1}\vect{u}$, and closes
with an $\oom{n}$ correction. The saving is the same discipline as LU
reuse across right-hand sides: never recompute what a cheap update can
correct.}
\label{fig:vol02:ch09:sherman-morrison}
\end{figure}


\begin{example}[A recalibrated sensor by rank-one update]
A weighted least-squares fit has already solved $\mat{A}\vect{x} =
\vect{b}$ with $\mat{A} = \begin{psmallmatrix} 4 & 1 \\ 1 & 3
\end{psmallmatrix}$, whose inverse is $\mat{A}^{-1} = \tfrac{1}{11}
\begin{psmallmatrix} 3 & -1 \\ -1 & 4 \end{psmallmatrix}$. One
measurement is reweighted, adding the rank-one term $\vect{u}
\vect{v}^{\transpose}$ with $\vect{u} = \vect{v} = (1, 0)^{\transpose}$,
so the $(1,1)$ entry rises from $4$ to $5$. Rather than invert the new
matrix, apply the formula. First $\mat{A}^{-1}\vect{u} = \tfrac{1}{11}
(3, -1)^{\transpose}$, and the denominator is $1 + \vect{v}^{\transpose}
\mat{A}^{-1}\vect{u} = 1 + 3/11 = 14/11$. The correction term is
\[
\frac{(\mat{A}^{-1}\vect{u})(\vect{v}^{\transpose}\mat{A}^{-1})}
{14/11} = \frac{11}{14}\cdot\frac{1}{121}
\begin{pmatrix} 3 \\ -1 \end{pmatrix}
\begin{pmatrix} 3 & -1 \end{pmatrix}
= \frac{1}{154}\begin{pmatrix} 9 & -3 \\ -3 & 1 \end{pmatrix},
\]
so the updated inverse is $\mat{A}^{-1}$ minus this, and the reader who
forms $\begin{psmallmatrix} 5 & 1 \\ 1 & 3 \end{psmallmatrix}^{-1} =
\tfrac{1}{14}\begin{psmallmatrix} 3 & -1 \\ -1 & 5 \end{psmallmatrix}$
directly confirms the two agree. At $2 \times 2$ the direct inverse is
as cheap, and the point is the scaling: at $n = 10^{4}$ the update is
one back-solve against a full refactorisation, four orders of magnitude
apart. The denominator also serves as a diagnostic, since a value near
zero warns that the edit is pushing the matrix toward singularity, the
rank-one signature of a measurement about to make the system
ill-posed.
\end{example}

\begin{mastery}
The block generalisation is the \engterm{Woodbury identity}, which
updates the inverse under a rank-$k$ change $\mat{U}\mat{C}
\mat{V}^{\transpose}$ with $\mat{U} \in \R^{n \times k}$, $\mat{C} \in
\R^{k \times k}$, $\mat{V} \in \R^{n \times k}$:
\[
(\mat{A} + \mat{U}\mat{C}\mat{V}^{\transpose})^{-1} = \mat{A}^{-1} -
\mat{A}^{-1}\mat{U}\bigl(\mat{C}^{-1} + \mat{V}^{\transpose}\mat{A}^{-1}
\mat{U}\bigr)^{-1}\mat{V}^{\transpose}\mat{A}^{-1}.
\]
The identity turns an $n \times n$ solve after a rank-$k$ update into
the original solve plus a $k \times k$ solve, a saving that is large
whenever $k \ll n$. Its reach across engineering is wide. The Kalman
filter of Volume IX is the Woodbury identity applied to a covariance
matrix as each measurement of rank equal to the sensor count arrives;
the recursive least-squares estimator updates its normal-equations
inverse the same way as each new data row is appended; and quasi-Newton
optimisers (BFGS and its limited-memory variant) maintain an
approximate inverse Hessian by exactly these low-rank updates, never
factoring the full Hessian at all. Recognising a problem as a sequence
of low-rank updates to a fixed base matrix is the cue to reach for
Woodbury rather than a re-solve, and the small inner matrix $\mat{C}^{-1}
+ \mat{V}^{\transpose}\mat{A}^{-1}\mat{U}$ carries the conditioning of
the update in a $k \times k$ object the engineer can inspect directly.
\end{mastery}

\subsection{The pseudoinverse and the minimum-norm solution}

A square invertible matrix has a clean inverse; a rectangular or
rank-deficient matrix does not, yet engineering still asks for a
``solution'' of $\mat{A}\vect{x} = \vect{b}$. The
\engterm{Moore--Penrose pseudoinverse} $\mat{A}^{+}$ supplies the
answer that two distinct cases require, and it reduces to $\mat{A}^{-1}$
when the matrix is square and invertible. For a full-column-rank
$\mat{A} \in \R^{m \times n}$ with $m > n$ (the overdetermined,
least-squares case), $\mat{A}^{+} = (\mat{A}^{\transpose}\mat{A})^{-1}
\mat{A}^{\transpose}$, and $\mat{A}^{+}\vect{b}$ is the least-squares
solution of section 9.6. For a full-row-rank $\mat{A}$ with $m < n$
(the underdetermined case, more unknowns than equations), $\mat{A}^{+}
= \mat{A}^{\transpose}(\mat{A}\mat{A}^{\transpose})^{-1}$, and
$\mat{A}^{+}\vect{b}$ is the \engterm{minimum-norm solution}: among
the infinitely many $\vect{x}$ that satisfy $\mat{A}\vect{x} =
\vect{b}$ exactly, it is the one of smallest Euclidean norm. The SVD
unifies both cases: with $\mat{A} = \mat{U}\mat{\Sigma}\mat{V}^{\transpose}$,
the pseudoinverse is $\mat{A}^{+} = \mat{V}\mat{\Sigma}^{+}
\mat{U}^{\transpose}$, where $\mat{\Sigma}^{+}$ inverts the nonzero
singular values and transposes the shape, leaving the zero (or
truncated) singular values at zero. This is the form the chapter's
project applies, and it is the only form that handles the
rank-deficient case without an arbitrary choice.

\begin{example}[Minimum-norm solution of an underdetermined system]
A robot arm with three joint angles must place its end effector at a
single target, two scalar conditions on three unknowns. The
linearised constraint is the underdetermined system $\mat{A}\vect{x}
= \vect{b}$ with
\[
\mat{A} = \begin{pmatrix} 1 & 1 & 1 \\ 1 & 2 & 3 \end{pmatrix},
\qquad \vect{b} = \begin{pmatrix} 6 \\ 14 \end{pmatrix}.
\]
The general solution is a line in $\R^{3}$: one particular solution
plus the one-dimensional null space $\spn\{(1, -2, 1)^{\transpose}\}$.
The minimum-norm solution picks the point on that line closest to the
origin, the smallest joint motion that reaches the target. Compute
$\mat{A}\mat{A}^{\transpose} = \begin{psmallmatrix} 3 & 6 \\ 6 & 14
\end{psmallmatrix}$, with inverse $\tfrac{1}{6}\begin{psmallmatrix}
14 & -6 \\ -6 & 3 \end{psmallmatrix}$. Then $(\mat{A}\mat{A}^{
\transpose})^{-1}\vect{b} = \tfrac{1}{6}(14 \cdot 6 - 6 \cdot 14,
-6 \cdot 6 + 3 \cdot 14)^{\transpose} = \tfrac{1}{6}(0, 6)^{\transpose}
= (0, 1)^{\transpose}$, and $\vect{x} = \mat{A}^{\transpose}(0,
1)^{\transpose} = (1, 2, 3)^{\transpose}$. The reader checks
$\mat{A}(1,2,3)^{\transpose} = (6, 14)^{\transpose}$, so the
constraint is met exactly, and that $\vect{x}$ is orthogonal to the
null-space direction $(1, -2, 1)$, the geometric signature of the
minimum-norm point. Any other solution, $\vect{x} + t(1, -2, 1)$,
has strictly larger norm. The redundant degree of freedom is the
robot's freedom of posture; the minimum-norm choice spends it on the
least joint travel.
\end{example}

\begin{example}[The cost of exploiting bandedness]
A one-dimensional finite-difference or finite-element model produces a
\engterm{banded} matrix: nonzeros confined to a few diagonals around
the main diagonal, bandwidth $b$ meaning $A_{ij} = 0$ whenever
$\abs{i - j} > b$. Banded structure changes the economics of the
solve outright. Gaussian elimination on a banded matrix creates no
fill-in outside the band, so the factorisation costs $\oom{n b^{2}}$
flops and the storage is $\oom{n b}$ rather than the dense $\oom{n^{3}}$
and $\oom{n^{2}}$. For a tridiagonal system ($b = 1$), the count
collapses to $\oom{n}$, linear in the size: the Thomas algorithm
solves a tridiagonal system in about $8n$ flops, against the
$\tfrac{2}{3}n^{3}$ of a dense solve. A model with $n = 10^{6}$
unknowns and bandwidth $b = 100$ factorises in $\oom{n b^{2}} =
10^{10}$ flops, about a second on a working core, where the dense
count $\tfrac{2}{3}(10^{6})^{3} \approx 6.7 \times 10^{17}$ is years.
The discipline is to detect the band before the solve and to call a
banded factorisation; the dense routine on a banded matrix wastes the
structure and overflows memory long before it finishes. Pivoting
complicates the picture, since row swaps can widen the band, but a
banded matrix that is diagonally dominant or symmetric positive
definite (the common engineering case) needs no pivoting, and the
band is preserved exactly.
\end{example}

\begin{archetype}[Optimisation]
The least-squares problem (section 9.6) minimises $\norm{\mat{A}
\vect{x} - \vect{b}}_{2}^{2}$ over $\vect{x} \in \R^{n}$. It is the
canonical convex optimisation problem: the objective is a positive
semidefinite quadratic, the feasible set is all of $\R^{n}$, and
the minimum is attained whenever $\mat{A}$ has full column rank.
Every routine engineering fit (Volume I Chapter 8's straight line,
Volume III's truss equilibrium, Volume IX's Kalman update) is a
specialisation of this template. The optimisation archetype, named
in the volume opener, takes its formal home here and recurs in
Chapter 15 with the apparatus of constrained optimisation.
\end{archetype}

\section{Determinants}\pathtag{standard}

The \engterm{determinant} $\det \mat{A}$ of a square matrix is a
scalar that records two pieces of information at once: a sign that
tells whether $\mat{A}$ preserves or reverses orientation, and a
magnitude that tells how $\mat{A}$ scales $n$-dimensional volume.

\subsection{Geometric definition}

For $\mat{A} \in \R^{n \times n}$, $\abs{\det \mat{A}}$ is the
factor by which $\mat{A}$ multiplies $n$-dimensional volume.
Concretely, the unit cube in $\R^{n}$, with edges along
$\vect{e}_{1}, \ldots, \vect{e}_{n}$, maps under $\vect{x} \mapsto
\mat{A}\vect{x}$ to a parallelepiped with edges $\vect{a}_{1},
\ldots, \vect{a}_{n}$. The volume of that parallelepiped is
$\abs{\det \mat{A}}$. The sign of $\det \mat{A}$ is positive when
the columns form a right-handed system (the orientation of
$\vect{e}_{1}, \ldots, \vect{e}_{n}$ is preserved) and negative
when they form a left-handed system.

The geometric definition gives the determinant's properties for
free. A matrix with two equal columns has zero volume in its image,
so $\det \mat{A} = 0$. A matrix with linearly dependent columns
collapses the cube into lower dimension, so its determinant is
zero. Multiplying a column by $c$ scales the volume by $c$, so
$\det$ is linear in each column separately. Swapping two columns
flips orientation, so $\det$ changes sign under a column swap.

\begin{example}[The $2 \times 2$ and $3 \times 3$ formulas]
For $\mat{A} = \begin{psmallmatrix} a & b \\ c & d \end{psmallmatrix}$
the determinant is $ad - bc$, the signed area of the parallelogram
with edges $(a, c)$ and $(b, d)$. For
$\begin{psmallmatrix} 3 & 1 \\ 2 & 4 \end{psmallmatrix}$ the value is
$12 - 2 = 10$. For a $3 \times 3$ matrix the cofactor expansion along
the first row gives
\[
\det \begin{pmatrix} a_{11} & a_{12} & a_{13} \\ a_{21} & a_{22} &
a_{23} \\ a_{31} & a_{32} & a_{33} \end{pmatrix}
= a_{11}(a_{22}a_{33} - a_{23}a_{32}) - a_{12}(a_{21}a_{33} -
a_{23}a_{31}) + a_{13}(a_{21}a_{32} - a_{22}a_{31}).
\]
For $\begin{psmallmatrix} 2 & 0 & 1 \\ 1 & 3 & 2 \\ 0 & 1 & 1
\end{psmallmatrix}$, the expansion is $2(3 - 2) - 0 + 1(1 - 0) = 2 +
1 = 3$. The same value drops out of the product of the pivots in the
LU factorisation, which is the route to take for any matrix larger
than $3 \times 3$.
\end{example}

\begin{example}[A determinant that signals near-singularity]
The matrix $\begin{psmallmatrix} 1 & 1 \\ 1 & 1.0001
\end{psmallmatrix}$ has determinant $1.0001 - 1 = 10^{-4}$, small but
nonzero. The smallness alone does not establish ill-conditioning,
because the determinant scales with the matrix entries: multiplying
the matrix by $10^{-2}$ multiplies the determinant by $10^{-4}$
without changing the conditioning at all. The condition number, a
ratio of singular values, is scale-invariant and is the correct
diagnostic; for this matrix $\kappa \approx 4 \times 10^{4}$
(Exercise 9.7). The determinant answers \emph{is the matrix exactly
singular}; the condition number answers \emph{how close to singular,
in a scale-free sense}. Section 9.7 makes the distinction precise.
\end{example}

\subsection{Algebraic recognition}

The determinant has a sum-over-permutations formula,

\[
\det \mat{A} = \sum_{\sigma \in S_{n}} \sgn(\sigma) \prod_{i=1}^{n}
A_{i, \sigma(i)},
\]

where $S_{n}$ is the set of permutations of $\{1, \ldots, n\}$. The
formula is the bridge between the geometric definition and any
practical computation. We mention it at recognition level: the
reader should know it exists and what it expresses. Computing a
$10 \times 10$ determinant by the permutation formula would require
$10!\approx 3.6 \times 10^{6}$ products and is never the chosen
method.

For practical computation the determinant is read off the LU
factorisation: $\det \mat{A} = \det \mat{P}^{-1} \cdot \det \mat{L}
\cdot \det \mat{U} = (\pm 1) \cdot 1 \cdot \prod_{i=1}^{n} U_{ii}$,
where $\mat{P}^{-1}$'s determinant is $+1$ or $-1$ depending on the
parity of the row swaps. The cost is the LU cost: $\tfrac{2}{3}n^{3}$
flops, the same as solving a system.

\begin{example}[A $4 \times 4$ determinant by elimination]
The cofactor expansion of a $4 \times 4$ determinant costs $24$
products by the permutation formula and is error-prone by hand;
elimination is faster and is the method that scales. Take
\[
\mat{A} = \begin{pmatrix}
2 & 1 & 0 & 0 \\
1 & 2 & 1 & 0 \\
0 & 1 & 2 & 1 \\
0 & 0 & 1 & 2
\end{pmatrix},
\]
the tridiagonal matrix that appears in every one-dimensional
finite-difference scheme. Eliminate column by column, recording each
pivot. Row 2 minus $\tfrac{1}{2}$ row 1 gives a $(2,2)$ pivot of
$2 - 1/2 = 3/2$. Row 3 minus $\tfrac{1}{(3/2)} = \tfrac{2}{3}$ times
the new row 2 gives a $(3,3)$ pivot of $2 - 2/3 = 4/3$. Row 4 minus
$\tfrac{1}{(4/3)} = \tfrac{3}{4}$ times the new row 3 gives a $(4,4)$
pivot of $2 - 3/4 = 5/4$. No row swaps occurred, so $\det \mat{A}$ is
the product of the pivots,
\[
\det \mat{A} = 2 \cdot \tfrac{3}{2} \cdot \tfrac{4}{3} \cdot
\tfrac{5}{4} = 5.
\]
The pivots telescope, $2 \cdot \tfrac{3}{2} \cdot \tfrac{4}{3} \cdot
\tfrac{5}{4} = 5$, which is the general pattern for the $n \times n$
version of this matrix: its determinant is $n + 1$. The positivity of
every pivot is also the certificate that $\mat{A}$ is positive
definite, the property the Cholesky example of section 9.3 exploited
and the property that makes the finite-difference matrix solvable by
the cheap factorisation.
\end{example}

\subsection{Properties used in practice}

Three properties carry their weight in engineering work. First,
$\det(\mat{A}\mat{B}) = \det \mat{A} \cdot \det \mat{B}$, so the
determinant of a product is the product of determinants. Second,
$\det \mat{A}^{\transpose} = \det \mat{A}$, so row and column
operations on the determinant are interchangeable. Third, $\mat{A}$
is invertible if and only if $\det \mat{A} \neq 0$; the converse
provides a clean test of singularity at the level of pencil-and-paper
work, though for finite-precision computation the condition number
of section 9.7 is more informative than the determinant.

\begin{example}[A rank-one update to a determinant]
The \engterm{matrix determinant lemma} is the determinant counterpart of
the Sherman-Morrison inverse update: a rank-one change to a matrix
changes its determinant by a scalar factor,
\[
\det(\mat{A} + \vect{u}\vect{v}^{\transpose}) = \det(\mat{A})\,
\bigl(1 + \vect{v}^{\transpose}\mat{A}^{-1}\vect{u}\bigr).
\]
The same scalar $1 + \vect{v}^{\transpose}\mat{A}^{-1}\vect{u}$ that
appeared in the Sherman-Morrison denominator reappears here, and its
vanishing is again the exact condition that the update makes the matrix
singular. Take $\mat{A} = \begin{psmallmatrix} 4 & 1 \\ 1 & 3
\end{psmallmatrix}$ with $\det\mat{A} = 11$ and $\mat{A}^{-1} =
\tfrac{1}{11}\begin{psmallmatrix} 3 & -1 \\ -1 & 4 \end{psmallmatrix}$,
and add the rank-one term $\vect{u} = \vect{v} = (1, 0)^{\transpose}$
that raised the $(1,1)$ entry to $5$. The scalar is $1 + \vect{v}^{
\transpose}\mat{A}^{-1}\vect{u} = 1 + 3/11 = 14/11$, so the new
determinant is $11 \cdot 14/11 = 14$, which the reader checks against
$\det\begin{psmallmatrix} 5 & 1 \\ 1 & 3 \end{psmallmatrix} = 15 - 1 =
14$. The lemma turns a determinant recomputation into one inner product
against the already-known inverse, and it is the tool behind incremental
determinant tracking in a sequence of low-rank edits, for instance
monitoring how close a stiffness matrix drifts toward singularity as a
structure is progressively modified.
\end{example}

\begin{example}[Tracking a determinant toward singularity]
The single-update lemma becomes a monitoring tool when the edits arrive
in sequence, each one a rank-one change, and the engineer watches the
determinant approach zero. Start from $\mat{A}_{0} = \begin{psmallmatrix}
4 & 1 \\ 1 & 3 \end{psmallmatrix}$ with $\det\mat{A}_{0} = 11$,
representing a two-degree-of-freedom stiffness matrix, and suppose a
progressive weakening removes stiffness from the first degree of
freedom in three steps, each subtracting $1$ from the $(1,1)$ entry,
that is $\vect{u} = \vect{v} = (1, 0)^{\transpose}$ scaled by $-1$ per
step. The scalar factor of the lemma at each step is $1 +
\vect{v}^{\transpose}\mat{A}_{k}^{-1}\vect{u}$ with the sign of the
edit folded in. For the first step, $\mat{A}_{0}^{-1} = \tfrac{1}{11}
\begin{psmallmatrix} 3 & -1 \\ -1 & 4 \end{psmallmatrix}$, so the
$(1,1)$ entry of the inverse is $3/11$ and the factor is $1 - 3/11 =
8/11$; the new determinant is $11 \cdot 8/11 = 8$, matching
$\det\begin{psmallmatrix} 3 & 1 \\ 1 & 3 \end{psmallmatrix} = 8$
directly. For the second step, $\mat{A}_{1} = \begin{psmallmatrix} 3 &
1 \\ 1 & 3 \end{psmallmatrix}$ has inverse $(1,1)$ entry $3/8$, the
factor is $1 - 3/8 = 5/8$, and the determinant drops to $8 \cdot 5/8 =
5$, matching $\det\begin{psmallmatrix} 2 & 1 \\ 1 & 3
\end{psmallmatrix}$. For the third, $\mat{A}_{2} = \begin{psmallmatrix}
2 & 1 \\ 1 & 3 \end{psmallmatrix}$ has inverse $(1,1)$ entry $3/5$, the
factor is $1 - 3/5 = 2/5$, and the determinant falls to $5 \cdot 2/5 =
2$, the value of $\det\begin{psmallmatrix} 1 & 1 \\ 1 & 3
\end{psmallmatrix}$. The sequence $11, 8, 5, 2$ tracks the structure
losing stiffness, and a fourth identical edit would give the factor
$1 - 3/2 = -1/2$ and a determinant of $-1$: the sign flip is the
lemma's warning that the matrix has passed through singularity. Each
update cost one inner product against the running inverse rather than a
fresh $2 \times 2$ determinant, and the same arithmetic at size $n$
costs $\oom{n^{2}}$ per step against the $\oom{n^{3}}$ of recomputing
the determinant from scratch. The vanishing factor, not the determinant
value itself, is the operative alarm, because it is scale free in the
same way the raw determinant is not.
\end{example}

The classroom result \engterm{Cramer's rule} expresses each
component of the solution to $\mat{A}\vect{x} = \vect{b}$ as a
ratio of determinants. Cramer's rule is theoretically clean and
practically useless: solving an $n \times n$ system by Cramer's
rule requires $n + 1$ determinant evaluations, costing
$\oom{n \cdot n^{3}} = \oom{n^{4}}$ flops, against $\oom{n^{3}}$
for elimination. We mention Cramer's rule for recognition; we do
not use it.

\begin{mastery}
Cramer's rule states $x_{i} = \det(\mat{A}_{i})/\det(\mat{A})$,
where $\mat{A}_{i}$ is $\mat{A}$ with its $i$-th column replaced by
$\vect{b}$. For the $2 \times 2$ system $\begin{psmallmatrix} 3 & 1
\\ 2 & 4 \end{psmallmatrix}\vect{x} = \begin{psmallmatrix} 5 \\ 6
\end{psmallmatrix}$, with $\det \mat{A} = 10$, the rule gives $x_{1}
= \det\begin{psmallmatrix} 5 & 1 \\ 6 & 4 \end{psmallmatrix}/10 =
(20 - 6)/10 = 1.4$ and $x_{2} = \det\begin{psmallmatrix} 3 & 5 \\ 2
& 6 \end{psmallmatrix}/10 = (18 - 10)/10 = 0.8$. The rule is
transparent at this size and a liability beyond it. Besides the
$\oom{n^{4}}$ cost, Cramer's rule is numerically poor: it computes
each component through a separate determinant, with no shared
elimination work and no pivoting discipline, so its floating-point
behaviour is worse than a single pivoted solve. The rule earns its
place in derivations (it yields the adjugate formula for the
inverse) and loses its place in computation.
\end{mastery}

\section{Inner products, norms, orthogonality}\pathtag{core}

Linear algebra acquires geometry through the inner product. The
\engterm{Euclidean inner product} (or dot product) on $\R^{n}$ is

\[
\inner*{\vect{u}, \vect{v}} = \vect{u}^{\transpose} \vect{v}
= \sum_{i=1}^{n} u_{i} v_{i}.
\]

Three properties define an inner product abstractly: bilinearity (it
is linear in each argument), symmetry ($\inner*{\vect{u}, \vect{v}}
= \inner*{\vect{v}, \vect{u}}$), and positive definiteness
($\inner*{\vect{v}, \vect{v}} \geq 0$, with equality if and only if
$\vect{v} = \vect{0}$). The Euclidean inner product satisfies all
three. Other inner products, weighted by a positive-definite matrix
$\mat{W}$ as $\inner*{\vect{u}, \vect{v}}_{\mat{W}} =
\vect{u}^{\transpose}\mat{W}\vect{v}$, recur in weighted
least-squares (Volume I Chapter 4 measurement uncertainty meets
section 9.6 here) and in numerical solvers for partial differential
equations (Chapter 16).

\subsection{Norms}

A \engterm{norm} on $\R^{n}$ is a function $\norm{\cdot}: \R^{n} \to
[0, \infty)$ satisfying $\norm{\vect{v}} = 0 \iff \vect{v} =
\vect{0}$, $\norm{c\vect{v}} = \abs{c}\norm{\vect{v}}$, and the
triangle inequality $\norm{\vect{u} + \vect{v}} \leq \norm{\vect{u}}
+ \norm{\vect{v}}$. The Euclidean norm $\norm{\vect{v}}_{2} =
\sqrt{\inner*{\vect{v}, \vect{v}}}$ is the one any inner product
defines. Two other norms are routinely useful. The 1-norm
$\norm{\vect{v}}_{1} = \sum_{i} \abs{v_{i}}$ measures total
absolute deviation; it appears in robust regression and in
compressed-sensing problems. The infinity-norm $\norm{\vect{v}}_{\infty}
= \max_{i} \abs{v_{i}}$ measures the worst component; it is the
engineering tolerance norm.

The three norms agree on $\vect{v} = \vect{0}$ and disagree
elsewhere. They are all equivalent on $\R^{n}$ in the sense that
each bounds the others up to a dimension-dependent constant; the
engineering reading of equivalence is that the choice of norm
matters when $n$ is large and the constants matter for the
inequality the engineer is using.

\begin{example}[The three norms disagree by the dimension]
For $\vect{v} = (3, -4, 12)^{\transpose}$, the norms are
$\norm{\vect{v}}_{1} = 3 + 4 + 12 = 19$, $\norm{\vect{v}}_{2} =
\sqrt{9 + 16 + 144} = \sqrt{169} = 13$, and $\norm{\vect{v}}_{\infty}
= 12$. The chain $\norm{\vect{v}}_{\infty} \leq \norm{\vect{v}}_{2}
\leq \norm{\vect{v}}_{1}$ holds, and the equivalence constants are
visible: $\norm{\vect{v}}_{1} \leq \sqrt{n}\,\norm{\vect{v}}_{2}$ and
$\norm{\vect{v}}_{2} \leq \sqrt{n}\,\norm{\vect{v}}_{\infty}$ for
$n = 3$, so $\sqrt{3}$ is the worst-case gap. The choice is not
cosmetic. A tolerance specified as ``no component exceeds
$0.1\,\si{\milli\meter}$'' is an infinity-norm bound; a tolerance
specified as ``the total deviation is below $0.1\,\si{\milli\meter}$''
is a 1-norm bound; a least-squares residual is a 2-norm. Reading a
specification means reading which norm it imposes.
\end{example}

A matrix norm $\norm{\mat{A}}$ measures how much $\mat{A}$ can
amplify a vector: $\norm{\mat{A}} = \max_{\vect{x} \neq \vect{0}}
\norm{\mat{A}\vect{x}}/\norm{\vect{x}}$. The induced 2-norm equals
the largest singular value $\sigma_{1}(\mat{A})$, a fact we use in
section 9.7. The induced 1-norm and infinity-norm have closed forms
that need no maximisation: the induced 1-norm is the largest
absolute column sum, and the induced infinity-norm is the largest
absolute row sum. These two are computable by inspection, which makes
them the norms to reach for when a quick bound on amplification is
wanted and the singular values are not in hand.

\begin{example}[The three induced matrix norms by hand]
Take $\mat{A} = \begin{psmallmatrix} 1 & -2 \\ 3 & 4
\end{psmallmatrix}$. The induced 1-norm is the largest column sum of
absolute values: column one sums to $\abs{1} + \abs{3} = 4$, column
two to $\abs{-2} + \abs{4} = 6$, so $\norm{\mat{A}}_{1} = 6$. The
induced infinity-norm is the largest row sum of absolute values: row
one gives $\abs{1} + \abs{-2} = 3$, row two gives $\abs{3} + \abs{4}
= 7$, so $\norm{\mat{A}}_{\infty} = 7$. The induced 2-norm is the
largest singular value, which requires the eigenvalues of
$\mat{A}^{\transpose}\mat{A} = \begin{psmallmatrix} 10 & 10 \\ 10 &
20 \end{psmallmatrix}$: the characteristic equation $(10 - \lambda)
(20 - \lambda) - 100 = 0$ gives $\lambda^{2} - 30\lambda + 100 = 0$,
so $\lambda = 15 \pm \sqrt{125}$, and $\sigma_{1} = \sqrt{15 +
5\sqrt{5}} \approx 5.12$. The three norms disagree, as they must,
and each bounds the true amplification in its own vector norm. The
1-norm and infinity-norm bracketed the 2-norm here without any
eigenvalue work, which is their value: an $\oom{n^{2}}$ scan of the
entries yields two bounds on the amplification, where the 2-norm
demands the $\oom{n^{3}}$ singular-value computation. The Frobenius
norm $\norm{\mat{A}}_{F} = \sqrt{\sum_{ij} A_{ij}^{2}} = \sqrt{1 + 4
+ 9 + 16} = \sqrt{30} \approx 5.48$ is a fourth measure, not induced
by any vector norm but cheap and often used as a stand-in for the
2-norm since $\norm{\mat{A}}_{2} \leq \norm{\mat{A}}_{F} \leq
\sqrt{r}\,\norm{\mat{A}}_{2}$ with $r$ the rank.
\end{example}

\subsection{Orthogonality}

Two vectors $\vect{u}$ and $\vect{v}$ are \engterm{orthogonal} when
$\inner*{\vect{u}, \vect{v}} = 0$. The Pythagorean identity holds
for orthogonal pairs: $\norm{\vect{u} + \vect{v}}_{2}^{2} =
\norm{\vect{u}}_{2}^{2} + \norm{\vect{v}}_{2}^{2}$. A set
$\{\vect{q}_{1}, \ldots, \vect{q}_{k}\}$ is \engterm{orthonormal}
when its members are mutually orthogonal and each has unit norm.
An orthonormal basis is a basis of orthonormal vectors.

Orthonormal bases simplify computation. If $\{\vect{q}_{1}, \ldots,
\vect{q}_{n}\}$ is an orthonormal basis of a subspace $W$ and
$\vect{w} \in W$, then the coordinates of $\vect{w}$ are inner
products with the basis vectors: $\vect{w} = \sum_{i} \inner*{\vect{w},
\vect{q}_{i}}\vect{q}_{i}$. The expansion is the engineer's working
form of the Fourier expansion the reader will meet in
Chapter 16.

The \engterm{Gram-Schmidt procedure} manufactures an orthonormal
basis from any independent set. Given $\vect{a}_{1}, \ldots,
\vect{a}_{k}$, it sets $\vect{q}_{1} = \vect{a}_{1}/\norm{\vect{a}_{1}}$
and then, for each subsequent vector, subtracts off its components
along the orthonormal vectors already built and normalises the
remainder:
\[
\vect{w}_{j} = \vect{a}_{j} - \sum_{i < j} \inner*{\vect{a}_{j},
\vect{q}_{i}}\vect{q}_{i}, \qquad \vect{q}_{j} = \vect{w}_{j}/
\norm{\vect{w}_{j}}.
\]
The orthogonalisation and the QR factorisation of section 9.3 are the
same arithmetic: the coefficients $\inner*{\vect{a}_{j}, \vect{q}_{i}}$
and the norms $\norm{\vect{w}_{j}}$ are exactly the entries of
$\mat{R}$, and the $\vect{q}_{j}$ are the columns of $\mat{Q}$.

\begin{example}[Gram-Schmidt on three vectors]
Orthonormalise $\vect{a}_{1} = (1, 1, 0)^{\transpose}$, $\vect{a}_{2}
= (1, 0, 1)^{\transpose}$, $\vect{a}_{3} = (0, 1, 1)^{\transpose}$.
First, $\norm{\vect{a}_{1}} = \sqrt{2}$, so $\vect{q}_{1} = (1, 1,
0)/\sqrt{2}$. Second, $\inner*{\vect{a}_{2}, \vect{q}_{1}} =
1/\sqrt{2}$, so $\vect{w}_{2} = (1, 0, 1) - (1/2, 1/2, 0) = (1/2,
-1/2, 1)$, with $\norm{\vect{w}_{2}} = \sqrt{1/4 + 1/4 + 1} =
\sqrt{3/2}$, giving $\vect{q}_{2} = (1, -1, 2)/\sqrt{6}$. Third,
$\inner*{\vect{a}_{3}, \vect{q}_{1}} = 1/\sqrt{2}$ and
$\inner*{\vect{a}_{3}, \vect{q}_{2}} = (0 - 1 + 2)/\sqrt{6} =
1/\sqrt{6}$, so $\vect{w}_{3} = (0, 1, 1) - (1/2)(1, 1, 0) -
(1/6)(1, -1, 2) = (-2/3, 2/3, 2/3)$, with $\norm{\vect{w}_{3}} =
\sqrt{4/9 \cdot 3} = 2/\sqrt{3}$, giving $\vect{q}_{3} = (-1, 1,
1)/\sqrt{3}$. The reader verifies $\inner*{\vect{q}_{i},
\vect{q}_{j}} = \delta_{ij}$ for all pairs. The three $\vect{q}$
vectors are an orthonormal basis of $\R^{3}$.
\end{example}

\begin{mastery}
A weighted inner product $\inner*{\vect{u}, \vect{v}}_{\mat{W}} =
\vect{u}^{\transpose}\mat{W}\vect{v}$, with $\mat{W}$ symmetric
positive definite, induces a weighted norm and a weighted notion of
orthogonality. The engineering use is in weighted least squares: when
measurements have unequal variances $\sigma_{i}^{2}$ (Volume I
Chapter 4), the right objective weights each residual by its inverse
variance, $\mat{W} = \diag(1/\sigma_{1}^{2}, \ldots, 1/\sigma_{m}^{2})$,
and the projection that minimises the weighted residual is the
maximum-likelihood estimate under Gaussian noise. The weighted
normal equations are $\mat{A}^{\transpose}\mat{W}\mat{A}\hat{\vect{x}}
= \mat{A}^{\transpose}\mat{W}\vect{b}$, and the conditioning analysis
of section 9.7 applies to $\mat{A}^{\transpose}\mat{W}\mat{A}$ in
place of the unweighted Gram matrix. The weighting is the formal
channel by which measurement uncertainty enters the fit.
\end{mastery}

\subsection{Cauchy-Schwarz and the angle}

The Cauchy-Schwarz inequality $\abs{\inner*{\vect{u}, \vect{v}}}
\leq \norm{\vect{u}}_{2} \norm{\vect{v}}_{2}$ holds for any inner
product. It justifies defining the angle $\theta$ between two
nonzero vectors by

\[
\cos \theta = \frac{\inner*{\vect{u}, \vect{v}}}
{\norm{\vect{u}}_{2} \norm{\vect{v}}_{2}}.
\]

The cosine is in $[-1, 1]$ by Cauchy-Schwarz, and the formula
recovers the elementary-geometry definition of angle in $\R^{2}$
and $\R^{3}$. The angle is a working notion of \emph{how far} two
directions disagree, and it underwrites the projection formulas of
section 9.6.

\begin{example}[Orthogonality in a function space]
The inner-product machinery is identical when the vectors are
functions. On the interval $[-1, 1]$, take the $L^{2}$ inner product
$\inner*{f, g} = \int_{-1}^{1} f(x) g(x)\,\dd x$. The monomials
$1$, $x$, $x^{2}$ are independent but not orthogonal: $\inner*{1, x^{2}}
= \int_{-1}^{1} x^{2}\,\dd x = 2/3 \neq 0$. Apply Gram-Schmidt. The
first vector is $p_{0} = 1$, with $\norm{p_{0}}^{2} = \int_{-1}^{1}
\dd x = 2$. The second: $\inner*{x, 1} = \int_{-1}^{1} x\,\dd x = 0$,
so $x$ is already orthogonal to $1$, and $p_{1} = x$, with
$\norm{p_{1}}^{2} = \int_{-1}^{1} x^{2}\,\dd x = 2/3$. The third:
subtract the components of $x^{2}$ along $p_{0}$ and $p_{1}$. The
$p_{1}$ component vanishes by symmetry ($\int_{-1}^{1} x^{3}\,\dd x =
0$); the $p_{0}$ component is $\inner*{x^{2}, 1}/\norm{p_{0}}^{2} =
(2/3)/2 = 1/3$, so $p_{2} = x^{2} - 1/3$. These are the Legendre
polynomials up to normalisation. The orthogonality the construction
forces, $\inner*{p_{i}, p_{j}} = 0$ for $i \neq j$, is exactly the
property that makes the Legendre basis well conditioned for
polynomial fitting where the monomial basis fails (section 9.7). The
finite-dimensional Gram-Schmidt of the previous example and this
infinite-dimensional one are the same arithmetic; the sum over
components becomes an integral, and nothing else changes. The
spectral methods of Volume IV and the orthogonal-function expansions
of Chapter 16 run on this identity.
\end{example}

\section{Projection, least squares}\pathtag{core}

A linear system $\mat{A}\vect{x} = \vect{b}$ may have no exact
solution, either because $\vect{b} \notin \col(\mat{A})$ or because
the system is overdetermined ($m > n$ with no redundancy in
$\vect{b}$). The least-squares problem replaces \emph{solve} with
\emph{come closest}: find $\hat{\vect{x}}$ minimising
$\norm{\mat{A}\vect{x} - \vect{b}}_{2}^{2}$.

\subsection{Orthogonal projection}

Let $W$ be a subspace of $\R^{m}$ and $\vect{b} \in \R^{m}$. The
\engterm{orthogonal projection} of $\vect{b}$ onto $W$ is the
unique vector $\proj_{W}\vect{b} \in W$ such that $\vect{b} -
\proj_{W}\vect{b}$ is orthogonal to every vector in $W$. The
projection minimises $\norm{\vect{b} - \vect{w}}_{2}$ over
$\vect{w} \in W$.

When $W = \col(\mat{A})$ for full-column-rank $\mat{A}$, the
projection is

\[
\proj_{\col(\mat{A})}\vect{b} = \mat{A}(\mat{A}^{\transpose}\mat{A})^{-1}
\mat{A}^{\transpose}\vect{b}.
\]

The matrix $\mat{P} = \mat{A}(\mat{A}^{\transpose}\mat{A})^{-1}
\mat{A}^{\transpose}$ is the projection matrix onto $\col(\mat{A})$;
it satisfies $\mat{P}^{2} = \mat{P}$ (idempotent) and
$\mat{P}^{\transpose} = \mat{P}$ (symmetric), the algebraic
fingerprint of an orthogonal projection.
Figure~\ref{fig:vol02:ch09:projection} draws the geometry.

\begin{figure}[ht]
\centering
\begin{tikzpicture}[x=1cm, y=1cm, font=\small]

% Plane W in 3D-ish perspective
\draw[fill=enggray!10, draw=enggray, thick]
  (-3.5, -1.6) -- (3.5, -1.0) -- (4.2, 1.6) -- (-2.8, 1.0) -- cycle;
\node[enggray, anchor=north west] at (-3.5, -1.6) {$W = \col(\mat{A})$};

% Origin
\fill (0, 0) circle (1.5pt);
\node[anchor=north east, font=\scriptsize] at (0, 0) {$\vect{0}$};

% Basis vectors a_1, a_2 spanning W
\draw[->, thick, engblue] (0, 0) -- (2.4, -0.4) node[anchor=west] {$\vect{a}_{1}$};
\draw[->, thick, engblue] (0, 0) -- (1.0, 1.2) node[anchor=south] {$\vect{a}_{2}$};

% Vector b
\draw[->, very thick, engred] (0, 0) -- (1.6, 3.0) node[anchor=south] {$\vect{b}$};

% Projection P b onto W
\draw[->, very thick, engblue] (0, 0) -- (1.6, 0.6) node[anchor=north west, pos=0.6] {$\proj_{W}\vect{b}$};

% Residual b - P b
\draw[->, thick, engamber, dashed] (1.6, 0.6) -- (1.6, 3.0)
  node[midway, anchor=west] {$\vect{b} - \proj_{W}\vect{b}$};

% Right-angle marker
\draw[thick] (1.6, 0.6) -- ++(0.18, -0.08) -- ++(0.08, 0.18) -- ++(-0.18, 0.08);

\node[anchor=north, font=\scriptsize, text width=10cm, align=center] at (0.5, -2.4) {
  The residual $\vect{b} - \proj_{W}\vect{b}$ is orthogonal to every vector in $W$;
  $\norm{\vect{b} - \vect{w}}_{2}$ is minimised over $\vect{w} \in W$ at $\vect{w} = \proj_{W}\vect{b}$.
};

\end{tikzpicture}
\caption[Orthogonal projection onto a column space]{Orthogonal
projection of $\vect{b} \in \R^{m}$ onto $W = \col(\mat{A})$. The
fitted vector $\hat{\vect{b}} = \mat{A}\hat{\vect{x}}$ is the
projection; the residual $\vect{r} = \vect{b} - \hat{\vect{b}}$ is
the component of $\vect{b}$ that cannot be reached as any linear
combination of the columns of $\mat{A}$. Least squares chooses
$\hat{\vect{x}}$ so the residual is orthogonal to every column of
$\mat{A}$, which is the geometric statement of the normal
equations.}
\label{fig:vol02:ch09:projection}
\end{figure}


\begin{example}[Distance from a point to a plane by projection]
The projection gives the shortest distance from a point to a subspace
directly, without calculus. The plane $W = \spn\{(1, 0, 1)^{\transpose},
(0, 1, 1)^{\transpose}\}$ through the origin is the column space of
$\mat{A} = \begin{psmallmatrix} 1 & 0 \\ 0 & 1 \\ 1 & 1
\end{psmallmatrix}$, and we want the distance from the point $\vect{b}
= (1, 1, 4)^{\transpose}$ to it. Form the Gram matrix
$\mat{A}^{\transpose}\mat{A} = \begin{psmallmatrix} 2 & 1 \\ 1 & 2
\end{psmallmatrix}$, invert it to $\tfrac{1}{3}\begin{psmallmatrix} 2 &
-1 \\ -1 & 2 \end{psmallmatrix}$, and compute $\mat{A}^{\transpose}
\vect{b} = (1 + 4, 1 + 4)^{\transpose} = (5, 5)^{\transpose}$. The
least-squares coefficients are $\hat{\vect{x}} = \tfrac{1}{3}
\begin{psmallmatrix} 2 & -1 \\ -1 & 2 \end{psmallmatrix}(5, 5)^{
\transpose} = \tfrac{1}{3}(5, 5)^{\transpose} = (5/3, 5/3)^{\transpose}$,
so the projection is $\proj_{W}\vect{b} = \mat{A}\hat{\vect{x}} =
(5/3, 5/3, 10/3)^{\transpose}$. The residual is $\vect{b} -
\proj_{W}\vect{b} = (1 - 5/3, 1 - 5/3, 4 - 10/3)^{\transpose} =
(-2/3, -2/3, 2/3)^{\transpose}$, and its length is the distance,
$\norm{\vect{b} - \proj_{W}\vect{b}}_{2} = \tfrac{2}{3}\sqrt{3} =
\tfrac{2}{\sqrt{3}} \approx 1.155$. The residual is orthogonal to both
columns of $\mat{A}$, which the reader checks: $(-2/3, -2/3, 2/3) \cdot
(1, 0, 1) = 0$ and $(-2/3, -2/3, 2/3) \cdot (0, 1, 1) = 0$. The plane
here is $\{(x, y, z) : x + y = z\}$, whose unit normal is $(1, 1, -1)/
\sqrt{3}$, and the elementary point-to-plane formula $\abs{(1)(1) +
(1)(1) + (-1)(4)}/\sqrt{3} = 2/\sqrt{3}$ reproduces the same distance,
confirming that the projection computes the geometric distance the
normal-vector formula gives. Figure~\ref{fig:vol02:ch09:point-plane-distance}
draws the perpendicular. The projection route is the one that
generalises: it needs no normal vector and works in any dimension and
for a subspace of any codimension, where the normal-vector formula
applies only to a hyperplane.
\end{example}

\begin{figure}[ht]
\centering
\begin{tikzpicture}[x=1.0cm, y=1.0cm, font=\small, >=Stealth]

% A plane drawn as a parallelogram (the subspace W = col(A))
\fill[engblue!10] (-0.5, 0) -- (5.0, 0) -- (6.5, 1.6) -- (1.0, 1.6) -- cycle;
\draw[engblue, thick] (-0.5, 0) -- (5.0, 0) -- (6.5, 1.6) -- (1.0, 1.6) -- cycle;
\node[engblue, anchor=west, font=\scriptsize] at (5.2, 0.3) {$W = \col(\mat{A})$};

% Origin
\fill[engblack] (1.5, 0.55) circle (1.4pt);
\node[anchor=north east, font=\scriptsize] at (1.5, 0.55) {$\vect{0}$};

% Point b above the plane
\fill[engred] (3.7, 3.0) circle (1.6pt);
\node[engred, anchor=west, font=\scriptsize] at (3.85, 3.05) {$\vect{b}$};

% Foot of the perpendicular (projection)
\fill[engamber] (3.7, 0.85) circle (1.6pt);
\node[engamber, anchor=north, font=\scriptsize] at (3.7, 0.72) {$\proj_{W}\vect{b}$};

% Residual (perpendicular) segment
\draw[engred, thick, ->] (3.7, 0.85) -- (3.7, 2.92);
\node[engred, anchor=west, font=\scriptsize] at (3.78, 1.9) {$\vect{b} - \proj_{W}\vect{b}$};

% Right-angle mark at the foot
\draw[engamber] (3.7, 0.85) ++(-0.28, 0) -- ++(0, 0.28) -- ++(0.28, 0);

% Vector from origin to b and to the projection
\draw[enggray, dashed] (1.5, 0.55) -- (3.7, 3.0);
\draw[enggray, dashed] (1.5, 0.55) -- (3.7, 0.85);

\node[anchor=north, font=\scriptsize, text width=11cm, align=center]
  at (2.8, -0.6) {
  The projection is the closest point of the subspace to $\vect{b}$, and
  the residual meets the subspace at a right angle. The length of the
  residual is the point-to-subspace distance the least-squares solve
  reports.
};

\end{tikzpicture}
\caption[Distance from a point to a subspace by projection]{The
orthogonal projection $\proj_{W}\vect{b}$ is the foot of the
perpendicular dropped from $\vect{b}$ to the subspace $W$, and the
residual $\vect{b} - \proj_{W}\vect{b}$ is orthogonal to every vector
in $W$. Among all points of $W$, the projection is the unique nearest
one to $\vect{b}$, so its distance $\norm{\vect{b} - \proj_{W}\vect{b}}_{2}$
is the least-squares residual. Reading the least-squares solution as a
distance-minimising projection is the geometric content behind the
normal equations.}
\label{fig:vol02:ch09:point-plane-distance}
\end{figure}


\subsection{The normal equations}

Setting the gradient of $\norm{\mat{A}\vect{x} - \vect{b}}_{2}^{2}$
to zero gives the \engterm{normal equations}:

\[
\mat{A}^{\transpose}\mat{A}\hat{\vect{x}} = \mat{A}^{\transpose}\vect{b}.
\]

When $\mat{A}$ has full column rank, $\mat{A}^{\transpose}\mat{A}$
is positive definite and invertible, and

\[
\hat{\vect{x}} = (\mat{A}^{\transpose}\mat{A})^{-1}
\mat{A}^{\transpose}\vect{b}.
\]

This is the formula Volume I Chapter 8 used for the slope and
intercept of a linear fit. The chapter is now in a position to read
it as projection: the fitted values $\mat{A}\hat{\vect{x}}$ are the
orthogonal projection of the data $\vect{b}$ onto the column space
of the design matrix.

The normal equations are the canonical least-squares formula. They
are also numerically dangerous when $\mat{A}$ is ill-conditioned, a
fact section 9.7 makes precise. The condition number of
$\mat{A}^{\transpose}\mat{A}$ is the square of the condition number
of $\mat{A}$. Forming $\mat{A}^{\transpose}\mat{A}$ explicitly
amplifies whatever conditioning trouble $\mat{A}$ already has.

\subsection{The QR alternative}

If $\mat{A} = \mat{Q}\mat{R}$ with $\mat{Q} \in \R^{m \times n}$
orthogonal in the rectangular sense and $\mat{R} \in \R^{n \times n}$
upper triangular, then $\mat{A}^{\transpose}\mat{A} =
\mat{R}^{\transpose}\mat{Q}^{\transpose}\mat{Q}\mat{R} =
\mat{R}^{\transpose}\mat{R}$. The normal equations become
$\mat{R}^{\transpose}\mat{R}\hat{\vect{x}} = \mat{R}^{\transpose}
\mat{Q}^{\transpose}\vect{b}$, which simplifies to

\[
\mat{R}\hat{\vect{x}} = \mat{Q}^{\transpose}\vect{b}.
\]

The system is upper triangular and back-solves in $O(n^{2})$ flops
once $\mat{Q}$ and $\mat{R}$ are in hand. The QR route avoids ever
forming $\mat{A}^{\transpose}\mat{A}$, and it works in a regime
where the normal equations would be numerically dangerous. Its
backward error is governed by the conditioning of $\mat{A}$, not
the squared conditioning of $\mat{A}^{\transpose}\mat{A}$
\cite[ch.~19]{text:trefethen-bau1997}.

\subsection{Worked example: a straight-line fit}

Volume I Chapter 8's straight-line fit had the design matrix

\[
\mat{A} = \begin{pmatrix} 1 & x_{1} \\ 1 & x_{2} \\ \vdots & \vdots
\\ 1 & x_{m} \end{pmatrix}, \quad
\vect{b} = \begin{pmatrix} y_{1} \\ y_{2} \\ \vdots \\ y_{m}
\end{pmatrix}.
\]

The first column carries the intercept and the second the slope.
The Gram matrix is

\[
\mat{A}^{\transpose}\mat{A} = \begin{pmatrix} m & \sum x_{i} \\
\sum x_{i} & \sum x_{i}^{2} \end{pmatrix}.
\]

The normal equations recover the closed-form expressions Volume I
Chapter 8 stated for the intercept $\hat{\beta}_{0}$ and slope
$\hat{\beta}_{1}$:

\[
\hat{\beta}_{1} = \frac{m\sum x_{i} y_{i} - \sum x_{i} \sum y_{i}}
{m\sum x_{i}^{2} - (\sum x_{i})^{2}}, \quad
\hat{\beta}_{0} = \bar{y} - \hat{\beta}_{1}\bar{x}.
\]

The reader who works through the algebra sees that the $(1,1)$
entry of $(\mat{A}^{\transpose}\mat{A})^{-1}$ multiplied by the
appropriate component of $\mat{A}^{\transpose}\vect{b}$ produces
$\hat{\beta}_{0}$, and similarly for the slope. The Volume I
formula is the matrix formula in two-dimensional coordinates.

\subsection{Higher-degree polynomial fits}

Replace the design matrix with the \engterm{Vandermonde matrix}

\[
\mat{V} = \begin{pmatrix} 1 & x_{1} & x_{1}^{2} & \cdots & x_{1}^{p}
\\ 1 & x_{2} & x_{2}^{2} & \cdots & x_{2}^{p} \\ \vdots & \vdots &
\vdots & \ddots & \vdots \\ 1 & x_{m} & x_{m}^{2} & \cdots &
x_{m}^{p} \end{pmatrix},
\]

and the same least-squares formulas fit a polynomial of degree $p$.
The Vandermonde matrix is the canonical example of an
ill-conditioned design matrix as $p$ grows: its columns become
nearly linearly dependent (the higher powers of $x$ correlate
strongly with each other on any interval that does not include
zero), and the condition number of $\mat{V}^{\transpose}\mat{V}$
explodes. Section 9.7 quantifies the trouble; the project of this
chapter walks the reader through it, and
Figure~\ref{fig:vol02:ch09:vandermonde-conditioning} previews the
result for $m = 100$ equispaced points in three bases.

\begin{figure}[p]
\centering
\begin{tikzpicture}[x=1.0cm, y=0.55cm, font=\small]

% Title
\node[font=\small\bfseries, anchor=south] at (5.5, 11.2) {Vandermonde conditioning on $m = 100$ equispaced points};

% Axes
\draw[->, thick] (0, 0) -- (11.2, 0) node[anchor=west] {polynomial degree $p$};
\draw[->, thick] (0, 0) -- (0, 10.8) node[anchor=south] {$\log_{10}\kappa$};

% x ticks: degree 1..10
\foreach \p in {1,2,3,4,5,6,7,8,9,10} {
  \draw (\p, 0.1) -- (\p, -0.1) node[anchor=north, font=\scriptsize] {\p};
}
% y ticks
\foreach \y in {0,2,4,6,8,10} {
  \draw (0.1, \y) -- (-0.1, \y) node[anchor=east, font=\scriptsize] {\y};
}

% Curve A: x in [0,1] (high conditioning)
% log10 kappa approx: 1->1.5, 2->3, 3->4.5, 4->6.2, 5->7.8, 6->9.3, 7->10.6 (clip), 8->11.8, 9->12.8, 10->13.6
\draw[very thick, engred]
  (1, 1.5) -- (2, 3.0) -- (3, 4.5) -- (4, 6.2) -- (5, 7.8) -- (6, 9.3) -- (7, 10.6);
\node[engred, anchor=west] at (3.2, 4.8) {monomial basis on $[0,1]$};

% Curve B: x in [-1,1]
% log10 kappa approx: 1->0.5, 2->1.0, 3->1.6, 4->2.2, 5->2.9, 6->3.6, 7->4.4, 8->5.2, 9->6.0, 10->6.8
\draw[very thick, engblue]
  (1, 0.5) -- (2, 1.0) -- (3, 1.6) -- (4, 2.2) -- (5, 2.9) -- (6, 3.6)
  -- (7, 4.4) -- (8, 5.2) -- (9, 6.0) -- (10, 6.8);
\node[engblue, anchor=west] at (5.6, 3.3) {monomial basis on $[-1,1]$};

% Curve C: Chebyshev basis
\draw[very thick, enggray, dashed]
  (1, 0.3) -- (2, 0.5) -- (3, 0.7) -- (4, 0.9) -- (5, 1.1) -- (6, 1.3)
  -- (7, 1.5) -- (8, 1.7) -- (9, 1.9) -- (10, 2.1);
\node[enggray, anchor=west] at (6.4, 1.0) {Chebyshev basis};

% Double-precision normal-equations refusal line: log10 kappa = 8 means Gram cond 16; we plot Gram refuse at sqrt
\draw[dashed, engamber] (0, 8) -- (10.8, 8);
\node[engamber, anchor=west, font=\scriptsize] at (7.4, 8.3) {refuse normal equations (double)};

% double precision break line
\draw[dashed, enggray] (0, 16);  % off-axis but show note instead

% Caption-overlay note
\node[font=\scriptsize, text width=10cm, align=center, anchor=north] at (5.5, -1.1) {
  Condition numbers computed for a $100 \times (p+1)$ Vandermonde matrix on
  equispaced points; values rounded to one decimal in $\log_{10}\kappa$.
};

\end{tikzpicture}
\caption[Vandermonde conditioning as polynomial degree grows]{Order-of-magnitude
condition number of the design matrix for a polynomial fit of degree
$p$ on $100$ equispaced points, in three formulations. The monomial
basis on $[0,1]$ is the canonical pathology: each added degree
adds roughly $1.5$ orders of magnitude to $\log_{10}\kappa$, and the
degree-six fit already lies past the normal-equations refusal line
in double precision. Rescaling $x$ to $[-1,1]$ recovers about three
orders of magnitude. Switching to the Chebyshev basis on $[-1,1]$
keeps the conditioning near unity through degree ten. The figure
encodes the central project result: the conditioning of a fit is a
choice, not a fact of the data, and that choice is the engineer's
to make.}
\label{fig:vol02:ch09:vandermonde-conditioning}
\end{figure}


\subsection{Worked example: a quadratic drag fit}

A wind-tunnel test measures the drag force $F$ on a model at five
airspeeds $v$. The recorded pairs $(v\,[\si{\meter\per\second}],
F\,[\si{\newton}])$ are $(10, 2.1)$, $(20, 8.0)$, $(30, 18.2)$,
$(40, 31.9)$, $(50, 50.5)$. Aerodynamic drag scales as $F = c v^{2}$
to leading order, but the fit is taken as $F = c_{0} + c_{1} v +
c_{2} v^{2}$ to let the data report any linear or constant offset.
The design matrix and right-hand side are
\[
\mat{V} = \begin{pmatrix} 1 & 10 & 100 \\ 1 & 20 & 400 \\ 1 & 30 &
900 \\ 1 & 40 & 1600 \\ 1 & 50 & 2500 \end{pmatrix}, \qquad
\vect{F} = \begin{pmatrix} 2.1 \\ 8.0 \\ 18.2 \\ 31.9 \\ 50.5
\end{pmatrix}.
\]
The Gram matrix and the right-hand side of the normal equations,
computed from the column sums $\sum 1 = 5$, $\sum v = 150$, $\sum
v^{2} = 5500$, $\sum v^{3} = 2.25 \times 10^{5}$, $\sum v^{4} =
9.79 \times 10^{6}$, are
\[
\mat{V}^{\transpose}\mat{V} = \begin{pmatrix} 5 & 150 & 5500 \\ 150
& 5500 & 225000 \\ 5500 & 225000 & 9790000 \end{pmatrix}, \qquad
\mat{V}^{\transpose}\vect{F} = \begin{pmatrix} 110.7 \\ 4528 \\
197080 \end{pmatrix}.
\]
Solving the $3 \times 3$ system (by the elimination of section 9.3)
returns $\hat{\vect{c}} \approx (0.38, -0.032, 0.0206)^{\transpose}$.
The quadratic coefficient $c_{2} \approx 0.0206\,\si{\newton}/
(\si{\meter\per\second})^{2}$ dominates, as the drag law predicts;
the small constant and linear terms absorb the measurement scatter.
The fitted curve evaluated at $v = 30$ gives $0.38 - 0.95 + 18.6
\approx 18.0\,\si{\newton}$, against the measured $18.2$, a residual
of about $0.2\,\si{\newton}$, consistent with the scatter in the
other points. The condition number of this $\mat{V}$ on the raw
airspeeds is about $6.7 \times 10^{3}$, modest for a degree-two fit,
and the Gram matrix squares it to about $4.5 \times 10^{7}$, which
still leaves about nine digits in double precision. The script
\path{code/least_squares_three_ways.py} reproduces the fit and
reports the three condition numbers; the dataset is recorded in
\path{data/drag-fit.csv}. Push the same data to a degree-four fit on
raw airspeeds and the Gram conditioning passes $10^{15}$, the regime
section 9.7 tells the engineer to refuse.

\begin{example}[The same drag fit through QR, digit for digit]
The drag fit above formed $\mat{V}^{\transpose}\mat{V}$ and solved the
normal equations. The QR route never forms the Gram matrix, and on a
mildly conditioned problem the two must agree; carrying the QR solve on
the same data shows where the agreement comes from and where it would
break. Factor $\mat{V} = \mat{Q}\mat{R}$ with $\mat{Q} \in \R^{5 \times
3}$ and $\mat{R} \in \R^{3 \times 3}$ upper triangular. The
least-squares solution then back-solves $\mat{R}\hat{\vect{c}} =
\mat{Q}^{\transpose}\vect{F}$. The diagonal of $\mat{R}$ carries the
conditioning: for this $\mat{V}$ the entries fall off as roughly
$\abs{R_{11}} \approx 2.24$, $\abs{R_{22}} \approx 31.6$, $\abs{R_{33}}
\approx 1.1 \times 10^{3}$ once the columns are taken in raw units, and
the ratio of largest to smallest diagonal magnitude, about $10^{3}$,
matches $\kappa(\mat{V}) \approx 6.7 \times 10^{3}$ to within the usual
factor. The back-solve returns $\hat{\vect{c}} \approx (0.38, -0.032,
0.0206)^{\transpose}$, identical to the normal-equations answer to the
digits reported, because at $\kappa(\mat{V}) \approx 6.7 \times 10^{3}$
the squared Gram conditioning of $4.5 \times 10^{7}$ still leaves nine
double-precision digits and the two routes cannot yet diverge. The
lesson is the contrast waiting one degree higher: at the degree-four
fit whose Gram conditioning passes $10^{15}$, the QR diagonal ratio
would sit near $10^{7.5}$ while the normal-equations conditioning would
sit near $10^{15}$, and the two routes would part company by several
digits per coefficient. The QR route is insurance whose premium (twice
the flops of the normal equations) buys back the squared half of the
conditioning exactly when the problem needs it.
\end{example}

\subsection{Worked case study: adjusting a spirit-levelling network}

The drag fit had one independent variable and a single response. A
larger and more representative least-squares problem has a network
structure: many measurements constrain a few unknowns through an
incidence pattern, the measurements disagree, and the engineer must
both estimate the unknowns and decide which measurements to trust.
Surveying supplies the cleanest example, and the apparatus that
solves it is the same apparatus a structural analysis, a sensor
fusion, or a circuit solve uses. We carry one network from field
sheet to adjusted heights to a flagged blunder.

\paragraph{The measurements.}
A survey crew runs a spirit level between four benchmarks. Benchmark
$A$ is a national datum mark whose height is known and fixed,
$H_{A} = 100.000\,\si{\meter}$. The heights of $B$, $C$, and $D$ are
unknown. Each levelling run measures a height \emph{difference}
between two benchmarks, not an absolute height; a run from $P$ to $Q$
returns $\ell = H_{Q} - H_{P}$ plus measurement noise. The crew
records five runs, listed in Figure~\ref{fig:vol02:ch09:leveling-network}:
\[
\begin{array}{llr}
A \to B: & \ell_{1} = +5.10 & \text{(1 km)} \\
A \to C: & \ell_{2} = +2.34 & \text{(1 km)} \\
B \to C: & \ell_{3} = -2.81 & \text{(2 km)} \\
C \to D: & \ell_{4} = +0.92 & \text{(1 km)} \\
B \to D: & \ell_{5} = -2.95 & \text{(4 km)}
\end{array}
\]
Three unknowns, five observations: the network is overdetermined by
two, and the two redundant measurements are the loop closures.

\begin{figure}[ht]
\centering
\begin{tikzpicture}[x=1.0cm, y=1.0cm, font=\small,
  bm/.style={circle, draw, thick, fill=engblue!12, minimum size=8mm,
    inner sep=1pt},
  fixed/.style={circle, draw, very thick, fill=engamber!25,
    minimum size=8mm, inner sep=1pt},
  obs/.style={->, thick, engred, >=stealth}]

% Benchmarks (nodes)
\node[fixed] (A) at (0, 0) {$A$};
\node[bm]    (B) at (4, 0.8) {$B$};
\node[bm]    (C) at (3.2, -2.4) {$C$};
\node[bm]    (D) at (7, -1.0) {$D$};

% Fixed-height label
\node[anchor=north, font=\scriptsize, engamber] at (0, -0.55)
  {fixed, $H_A = 100.000$};

% Observed height differences (directed leg = measured rise from tail to head)
\draw[obs] (A) -- (B) node[midway, above, font=\scriptsize]
  {$\ell_1 = +5.10$};
\draw[obs] (A) -- (C) node[midway, below left, font=\scriptsize]
  {$\ell_2 = +2.34$};
\draw[obs] (B) -- (C) node[midway, right=2pt, font=\scriptsize]
  {$\ell_3 = -2.81$};
\draw[obs] (C) -- (D) node[midway, below, font=\scriptsize]
  {$\ell_4 = +0.92$};
\draw[obs] (B) -- (D) node[midway, above right, font=\scriptsize]
  {$\ell_5 = -2.95$};

\node[anchor=north, font=\scriptsize, text width=10.5cm, align=center]
  at (3.5, -3.4) {
  Five spirit-levelling runs (red, in metres) connect one fixed
  benchmark $A$ to three unknown benchmarks $B$, $C$, $D$. Each run
  measures a height \emph{difference}; the network has more
  observations ($m = 5$) than unknown heights ($n = 3$), so the
  redundant loops let least squares both estimate the heights and
  expose measurement blunders.
};

\end{tikzpicture}
\caption[A redundant spirit-levelling network]{A spirit-levelling
network with one fixed benchmark $A$ and three unknown benchmarks
$B$, $C$, $D$. Each directed leg records a measured height
difference. Three unknowns and five observations leave two redundant
measurements, the loop closures $A\!\to\!B\!\to\!C\!\to\!A$ and
$B\!\to\!D\!\to\!C\!\to\!B$, which would each sum to zero with
perfect instruments. The non-zero loop misclosures are the raw
material of the least-squares adjustment in the worked case of this
section.}
\label{fig:vol02:ch09:leveling-network}
\end{figure}


\paragraph{The misclosures, read before any algebra.}
The estimation habit of this book asks for the loop misclosures
before the normal equations. The loop $A \to B \to C$ should return
to the same height the direct run $A \to C$ reports:
$\ell_{1} + \ell_{3}$ against $\ell_{2}$ gives $5.10 - 2.81 - 2.34
= -0.05\,\si{\meter}$, a $5\,\si{\centi\meter}$ misclosure over
$4\,\si{\kilo\meter}$ of levelling, within the tolerance of ordinary
work. The loop $B \to D \to C$ closes through $\ell_{5}$, the reverse
of $\ell_{4}$, and the reverse of $\ell_{3}$: $\ell_{5} - \ell_{4}
- \ell_{3} = -2.95 - 0.92 + 2.81 = -1.06\,\si{\meter}$. A
$1.06\,\si{\meter}$ misclosure is a blunder, not noise. One of the
three runs in that loop carries a gross error, and the longest, the
$4\,\si{\kilo\meter}$ run $B \to D$, is the first suspect. The
adjustment will assign the residual; the misclosure already tells the
crew where to look.

\paragraph{The linear model.}
With $A$ fixed, the unknown vector is $\vect{x} = (H_{B}, H_{C},
H_{D})^{\transpose}$. Each observation $\ell_{k} = H_{Q} - H_{P}$ is
linear in the unknowns, so the design matrix has a $+1$ in the
column of the head benchmark, a $-1$ in the column of the tail, and
the fixed datum moves to the right-hand side. The system $\mat{A}
\vect{x} = \vect{b}$ is
\[
\mat{A} = \begin{pmatrix}
1 & 0 & 0 \\
0 & 1 & 0 \\
-1 & 1 & 0 \\
0 & -1 & 1 \\
-1 & 0 & 1
\end{pmatrix}, \qquad
\vect{b} = \begin{pmatrix}
\ell_{1} + H_{A} \\ \ell_{2} + H_{A} \\ \ell_{3} \\ \ell_{4} \\ \ell_{5}
\end{pmatrix}
= \begin{pmatrix}
105.10 \\ 102.34 \\ -2.81 \\ 0.92 \\ -2.95
\end{pmatrix}.
\]
The first column is the incidence of $B$ across the five runs, and so
on. The matrix has the structure of a graph incidence matrix, which
is the recurring algebraic shape of every conservation law on a
network: levelling, the resistor circuit of section 9.6's companion
example, a pipe flow balance, a traffic assignment. The rank is three
(the three columns are independent because $A$ anchors the network),
so the unknowns are determined.

\paragraph{The weights.}
Levelling error grows with distance, because the random errors of
successive setups accumulate along the run. Standard practice models
the variance of a run as proportional to its length, so the weight is
the inverse length: $w_{k} = 1/L_{k}$. The weight matrix is
$\mat{W} = \diag(1, 1, 1/2, 1, 1/4)$. The $4\,\si{\kilo\meter}$ run
$B \to D$ carries one-quarter the weight of a $1\,\si{\kilo\meter}$
run, the formal statement that the engineer trusts it least, which is
the weighted-least-squares connection to measurement uncertainty the
mastery box of section 9.5 set up.

\paragraph{The solve.}
The weighted normal equations are $\mat{A}^{\transpose}\mat{W}\mat{A}
\hat{\vect{x}} = \mat{A}^{\transpose}\mat{W}\vect{b}$. The weighted
Gram matrix is
\[
\mat{A}^{\transpose}\mat{W}\mat{A} = \begin{pmatrix}
1.75 & -0.50 & -0.25 \\
-0.50 & 1.50 & -1.00 \\
-0.25 & -1.00 & 1.25
\end{pmatrix},
\]
a symmetric, diagonally weighted graph Laplacian; its condition
number is $6.05$, so the system is benign and either the normal
equations or QR returns the same answer. Solving (by the elimination
of section 9.3, or by the script
\path{code/leveling_adjustment.py}) gives the adjusted heights
\[
\hat{H}_{B} = 105.203\,\si{\meter}, \quad
\hat{H}_{C} = 102.237\,\si{\meter}, \quad
\hat{H}_{D} = 102.976\,\si{\meter}.
\]
The covariance of the estimate is $\hat{\sigma}_{0}^{2}
(\mat{A}^{\transpose}\mat{W}\mat{A})^{-1}$, with the reference
variance $\hat{\sigma}_{0}^{2}$ read off the weighted residual sum of
squares divided by the redundancy $m - n = 2$. The reference standard
deviation comes out at $\hat{\sigma}_{0} = 0.31\,\si{\meter}$, and the
per-benchmark standard errors are $0.26\,\si{\meter}$ for $B$ and $C$
and $0.37\,\si{\meter}$ for $D$. $D$ is least certain, as it must be:
it sits at the far end of the long, lightly weighted run.

\paragraph{The residuals localise the blunder.}
The fitted residuals $\vect{r} = \mat{A}\hat{\vect{x}} - \vect{b}$
are, run by run,
\[
\vect{r} = (+0.10,\; -0.10,\; -0.16,\; -0.18,\; +0.72)^{\transpose}\,
\si{\meter}.
\]
Four residuals are at or below $0.18\,\si{\meter}$, in line with the
$0.31\,\si{\meter}$ reference deviation. The fifth, the $B \to D$
run, carries $+0.72\,\si{\meter}$, more than twice the reference
deviation, exactly where the loop misclosure pointed. The least
squares did not hide the bad run; it distributed most of the
discrepancy onto the lightly weighted measurement and left a residual
large enough to flag. The crew would re-run $B \to D$, and an
adjustment without that observation would tighten $\hat{H}_{D}$
toward the value the $A \to C \to D$ chain implies.

\paragraph{What the case teaches.}
The levelling adjustment is the least-squares template in working
dress. The design matrix is the incidence matrix of a measurement
graph; the weights encode measurement uncertainty as an inverse
variance; the normal-equations Gram matrix is a weighted graph
Laplacian, symmetric positive definite, so Cholesky is the natural
factorisation; the redundancy $m - n$ is the number of independent
checks the network supplies; and the residual vector, not the
point estimate alone, is the deliverable, because the residuals are
where blunders surface. Every structural-analysis adjustment, every
GPS network solution, and every sensor-fusion update in Volume IX
carries this shape. The dataset is recorded in
\path{data/leveling-network.csv} and the full adjustment, including
the covariance propagation, is in \path{code/leveling_adjustment.py}.

\subsection{Companion case: a resistor network by nodal analysis}

The same incidence-matrix algebra solves a linear circuit, with a
sign that makes the system symmetric positive definite outright. The
network of Figure~\ref{fig:vol02:ch09:resistor-network} has four
nodes, node $0$ grounded, a $1\,\si{\ampere}$ current injected at
node $1$, and five resistors. Nodal analysis writes Kirchhoff's
current law at each non-ground node: the sum of currents leaving the
node through its resistors equals the injected current. With
conductances $g_{ij} = 1/R_{ij}$, the unknown node voltages
$\vect{v} = (v_{1}, v_{2}, v_{3})^{\transpose}$ satisfy $\mat{G}
\vect{v} = \vect{i}$, where the diagonal entry $G_{kk}$ is the sum of
conductances touching node $k$ and the off-diagonal entry $G_{kj}$ is
the negative of the conductance directly between nodes $k$ and $j$.

\begin{figure}[ht]
\centering
\begin{tikzpicture}[x=1.0cm, y=1.0cm, font=\small,
  node/.style={circle, draw, thick, fill=engblue!12, minimum size=8mm,
    inner sep=1pt}]

% Nodes
\node[node] (n1) at (0, 0) {$1$};
\node[node] (n2) at (4, 1.6) {$2$};
\node[node] (n3) at (4, -1.6) {$3$};
\node[node] (g)  at (8, 0) {$0$};
\node[anchor=west, font=\scriptsize] at (8.5, 0) {ground};

% Current source into node 1
\draw[->, very thick, engred, >=stealth] (-2.4, 0) -- (n1)
  node[midway, above, font=\scriptsize] {$I = 1\,\text{A}$};

% Resistors as labelled edges
\draw[thick] (n1) -- (n2)
  node[midway, above left=-1pt, font=\scriptsize] {$R_{12} = 2\,\Omega$};
\draw[thick] (n1) -- (n3)
  node[midway, below left=-1pt, font=\scriptsize] {$R_{13} = 4\,\Omega$};
\draw[thick] (n2) -- (n3)
  node[midway, right=1pt, font=\scriptsize] {$R_{23} = 1\,\Omega$};
\draw[thick] (n2) -- (g)
  node[midway, above right=-1pt, font=\scriptsize] {$R_{20} = 3\,\Omega$};
\draw[thick] (n3) -- (g)
  node[midway, below right=-1pt, font=\scriptsize] {$R_{30} = 5\,\Omega$};

\node[anchor=north, font=\scriptsize, text width=10.5cm, align=center]
  at (3.0, -2.6) {
  A linear resistor network with a $1\,\text{A}$ current injected at
  node 1 and node 0 held at ground. Nodal analysis writes one
  Kirchhoff-current-law equation per unknown node, producing a
  symmetric positive-definite conductance matrix whose solve returns
  the node voltages.
};

\end{tikzpicture}
\caption[A resistor network as a symmetric positive-definite
system]{A four-node resistor network. Kirchhoff's current law at the
three non-ground nodes yields a $3 \times 3$ linear system
$\mat{G}\vect{v} = \vect{i}$, where $\mat{G}$ is the conductance
matrix, $\vect{v}$ the unknown node voltages, and $\vect{i}$ the
injected currents. The conductance matrix is symmetric (reciprocity
of resistors) and positive definite (the network dissipates positive
power for any non-trivial voltage), so Cholesky is the production
factorisation, half the cost of LU.}
\label{fig:vol02:ch09:resistor-network}
\end{figure}


\begin{example}[Solving the conductance system]
With $R_{12} = 2$, $R_{13} = 4$, $R_{23} = 1$, $R_{20} = 3$, $R_{30}
= 5$ ohms, the conductances are $g_{12} = 0.5$, $g_{13} = 0.25$,
$g_{23} = 1$, $g_{20} = 1/3$, $g_{30} = 0.2$ siemens. The conductance
matrix and injection are
\[
\mat{G} = \begin{pmatrix}
0.75 & -0.5 & -0.25 \\
-0.5 & 1.833 & -1.0 \\
-0.25 & -1.0 & 1.45
\end{pmatrix}, \qquad
\vect{i} = \begin{pmatrix} 1 \\ 0 \\ 0 \end{pmatrix}.
\]
The matrix is symmetric, because a resistor between two nodes
contributes equally to both rows, the algebraic face of reciprocity.
It is positive definite, because the power dissipated for any
non-trivial voltage pattern is $\vect{v}^{\transpose}\mat{G}\vect{v}
> 0$. The grounding of node $0$ is what makes it positive definite
rather than merely positive semidefinite: without a reference node
the Laplacian would have the constant vector in its null space, the
floating-ground singularity. Cholesky factorises $\mat{G} =
\mat{L}\mat{L}^{\transpose}$ at half the cost of LU, and forward and
back substitution return $\vect{v} \approx (3.21, 1.89,
1.85)^{\transpose}\,\si{\volt}$. The voltage at the driven node is
highest, as the injected current must climb the network's effective
resistance to reach ground. The same matrix shape, a grounded graph
Laplacian, governs heat conduction on a discretised bar, steady
diffusion, and the stiffness assembly of Volume III; recognising the
shape tells the engineer to reach for Cholesky before LU.
\end{example}

\subsection{Worked case study: a truss by the direct-stiffness method}\pathtag{standard}

The levelling and resistor networks were measurement and conservation
problems: an incidence matrix, a weighting, a solve. A structural
analysis assembles its matrix differently, by adding up the
contributions of physical elements, and the assembly is worth one full
worked pass because it is the algebraic skeleton of the entire
finite-element method that Volume III builds on. We carry the
three-bar truss of Figure~\ref{fig:vol02:ch09:truss-stiffness} from
geometry to displacements to member forces.

\begin{figure}[ht]
\centering
\begin{tikzpicture}[x=1.0cm, y=1.0cm, font=\small,
  joint/.style={circle, draw, thick, fill=engblue!12, minimum size=7mm,
    inner sep=1pt}]

% Geometry: a three-bar planar truss.
% Node 1 at origin (pinned), node 2 at (4,0) (roller), node 3 at (2,3) (loaded).
\coordinate (p1) at (0, 0);
\coordinate (p2) at (4, 0);
\coordinate (p3) at (2, 3);

% Members
\draw[very thick] (p1) -- (p2) node[midway, below=2pt, font=\scriptsize] {member $c$ ($1\!-\!2$)};
\draw[very thick] (p1) -- (p3) node[midway, above left=-1pt, font=\scriptsize] {member $a$ ($1\!-\!3$)};
\draw[very thick] (p2) -- (p3) node[midway, above right=-1pt, font=\scriptsize] {member $b$ ($2\!-\!3$)};

% Joints
\node[joint] at (p1) {$1$};
\node[joint] at (p2) {$2$};
\node[joint] at (p3) {$3$};

% Pin support at node 1
\draw[thick] (-0.45, -0.45) -- (0.45, -0.45);
\draw[thick] (0, 0) -- (-0.35, -0.45);
\draw[thick] (0, 0) -- (0.0, -0.45);
\draw[thick] (0, 0) -- (0.35, -0.45);
\foreach \dx in {-0.4,-0.2,0.0,0.2,0.4}
  {\draw[thick] ({\dx-0.12}, -0.6) -- ({\dx+0.05}, -0.45);}
\node[font=\scriptsize, anchor=north] at (0, -0.75) {pin};

% Roller support at node 2
\draw[thick] (3.7, -0.5) circle (0.12);
\draw[thick] (4.3, -0.5) circle (0.12);
\draw[thick] (3.5, -0.65) -- (4.5, -0.65);
\node[font=\scriptsize, anchor=north] at (4, -0.8) {roller};

% Applied load at node 3
\draw[->, very thick, engred, >=stealth] (p3) -- ++(0.9, -1.5)
  node[midway, right=1pt, font=\scriptsize] {$\vect{P} = (3,-4)\,\text{kN}$};

\node[anchor=north, font=\scriptsize, text width=10.5cm, align=center]
  at (2.0, -1.4) {
  A planar three-bar truss. Node 1 is pinned (two reactions), node 2
  is a roller (one vertical reaction), and a load is applied at node 3.
  Each bar contributes a $4 \times 4$ element stiffness matrix; the
  direct-stiffness method assembles these into a global matrix indexed
  by the six nodal degrees of freedom.};

\end{tikzpicture}
\caption[A three-bar truss and the direct-stiffness assembly]{A
planar three-bar truss solved by the direct-stiffness method. Each
member contributes an element stiffness matrix in global coordinates;
the assembly scatters these into a $6 \times 6$ global matrix, the
support conditions delete the constrained rows and columns, and the
reduced system returns the free displacements. Member axial forces
and support reactions follow by back-substitution. The global
stiffness matrix is symmetric positive definite once the rigid-body
modes are removed by the supports, so Cholesky is the natural
factorisation.}
\label{fig:vol02:ch09:truss-stiffness}
\end{figure}


\paragraph{The element stiffness matrix.}
A pin-jointed bar carries only axial force. Along its own axis a bar
of cross-sectional area $A$, modulus $E$, and length $L$ behaves as a
spring of stiffness $k = EA/L$: an axial extension $\delta$ produces
an axial force $k\delta$. In the global $(x, y)$ frame, a bar inclined
at angle $\theta$ to the horizontal maps its two nodal displacements
to two nodal forces through the $4 \times 4$ element stiffness matrix
\[
\mat{k}_{\mathrm{e}} = \frac{EA}{L}
\begin{pmatrix}
c^{2} & cs & -c^{2} & -cs \\
cs & s^{2} & -cs & -s^{2} \\
-c^{2} & -cs & c^{2} & cs \\
-cs & -s^{2} & cs & s^{2}
\end{pmatrix},
\qquad c = \cos\theta,\; s = \sin\theta,
\]
where the four rows and columns index the $x$ and $y$ displacements of
the bar's two end nodes. The matrix is symmetric (the bar stores
elastic energy, a quadratic form) and singular on its own (a bar with
no supports has rigid-body modes, the source of the zero determinant).
The block structure $\begin{psmallmatrix} \mat{K}_{0} & -\mat{K}_{0}
\\ -\mat{K}_{0} & \mat{K}_{0} \end{psmallmatrix}$ with $\mat{K}_{0} =
\tfrac{EA}{L}\begin{psmallmatrix} c^{2} & cs \\ cs & s^{2}
\end{psmallmatrix}$ is the block algebra of the previous section in
physical dress: action and reaction across the bar.

\paragraph{Geometry and element matrices.}
Take $EA = 1$ in consistent units for every bar, so the stiffness of
each is $1/L$. Node 1 sits at $(0, 0)$, node 2 at $(4, 0)$, node 3 at
$(2, 3)$. The three bars have:
\[
\begin{array}{lllll}
\text{member } c\ (1\!-\!2): & L = 4, & \theta = 0^{\circ}, & c = 1, & s = 0; \\
\text{member } a\ (1\!-\!3): & L = \sqrt{13}, & \cos\theta = 2/\sqrt{13}, & & \sin\theta = 3/\sqrt{13}; \\
\text{member } b\ (2\!-\!3): & L = \sqrt{13}, & \cos\theta = -2/\sqrt{13}, & & \sin\theta = 3/\sqrt{13}.
\end{array}
\]
Member $c$ is horizontal, so its only nonzero stiffness is along
$x$: $k = 1/4 = 0.25$. For members $a$ and $b$, $1/L = 1/\sqrt{13}
\approx 0.2774$, and the direction cosines give $c^{2} = 4/13$, $s^{2}
= 9/13$, $\abs{cs} = 6/13$, so the per-bar block scales these by
$0.2774$: $\tfrac{EA}{L}c^{2} \approx 0.0854$, $\tfrac{EA}{L}s^{2}
\approx 0.1921$, $\tfrac{EA}{L}\abs{cs} \approx 0.1280$.

\paragraph{Assembly.}
The global stiffness matrix $\mat{K}$ is $6 \times 6$, indexed by the
six degrees of freedom $(u_{1}, v_{1}, u_{2}, v_{2}, u_{3}, v_{3})$.
Each element matrix is \emph{scattered} into the global matrix: the
bar's four local entries land in the global rows and columns of the
two nodes it connects, and overlapping contributions add. Node 3, the
apex, receives contributions from both inclined bars, so its diagonal
block sums them. Assembling the three element matrices gives the free
part of $\mat{K}$ for the two unconstrained degrees of freedom at node
3, $(u_{3}, v_{3})$:
\[
\mat{K}_{\mathrm{ff}} = \begin{pmatrix}
\tfrac{EA}{L_{a}}c_{a}^{2} + \tfrac{EA}{L_{b}}c_{b}^{2} & \tfrac{EA}{L_{a}}c_{a}s_{a} + \tfrac{EA}{L_{b}}c_{b}s_{b} \\[2pt]
\tfrac{EA}{L_{a}}c_{a}s_{a} + \tfrac{EA}{L_{b}}c_{b}s_{b} & \tfrac{EA}{L_{a}}s_{a}^{2} + \tfrac{EA}{L_{b}}s_{b}^{2}
\end{pmatrix}
= \begin{pmatrix} 0.1708 & 0 \\ 0 & 0.3842 \end{pmatrix}.
\]
The off-diagonal coupling vanishes because the two inclined bars are
mirror images: $c_{a}s_{a} = -c_{b}s_{b}$, so the $cs$ terms cancel
exactly. The symmetry of the structure has decoupled the horizontal
and vertical response at the apex, a result the assembly hands the
engineer for free.

\paragraph{The constrained solve.}
Node 1 is pinned ($u_{1} = v_{1} = 0$) and node 2 is on a roller
($v_{2} = 0$, with $u_{2}$ free). Imposing the supports deletes the
constrained rows and columns. For the load $\vect{P} = (3, -4)\,
\text{kN}$ applied at node 3, and because the apex block decoupled,
the displacements at node 3 follow directly:
\[
u_{3} = \frac{3}{0.1708} \approx 17.6, \qquad
v_{3} = \frac{-4}{0.3842} \approx -10.4,
\]
in the consistent displacement units. The horizontal displacement at
node 2 couples in through member $c$ and resolves alongside; the full
$3 \times 3$ reduced system (the two apex freedoms plus $u_{2}$) is
symmetric positive definite, so Cholesky returns all three free
displacements in one factorisation. The positive definiteness is the
payoff of removing the rigid-body modes: the unconstrained $6 \times
6$ matrix is singular (it has three planar rigid-body modes, two
translations and one rotation), and the supports remove exactly those
three, leaving an invertible reduced matrix.

\paragraph{Member forces and reactions.}
The deliverable of a structural analysis is not the displacements but
the member forces and support reactions. Each bar's axial force is
$N = k\,(\text{axial extension}) = \tfrac{EA}{L}(c, s) \cdot
(\Delta u, \Delta v)$, the projection of the relative end displacement
onto the bar axis times the bar stiffness. The support reactions
recover from the deleted rows: multiplying the full $\mat{K}$ by the
now-known displacement vector returns the forces at the constrained
degrees of freedom, which are the reactions. The reactions must
balance the applied load, $\sum \vect{R} + \vect{P} = \vect{0}$, and
checking that balance is the engineer's arithmetic verification that
the assembly and solve are correct.

\paragraph{What the case teaches.}
The direct-stiffness method is the least-squares and network algebra
turned inside out. The global matrix is not measured and inverted; it
is \emph{assembled} from element contributions, each a small dense
block scattered into a large sparse whole. The supports are not data
to be fitted but constraints that remove the singularity. The matrix
is symmetric positive definite once constrained, so Cholesky is the
factorisation, and the sparsity (each node couples only to its bar
neighbours) is the structure a production code exploits to reach
$10^{6}$ degrees of freedom. Volume III develops the method in full
for frames, plates, and solids; the algebra is the algebra of this
chapter. The assembly pattern, scatter-and-add of element blocks, is
the same whether the element is a bar, a beam, a triangle, or a
tetrahedron.

\subsection{Applied vignette: positioning by trilateration}\pathtag{standard}

The levelling network estimated heights from height differences on a
graph. A satellite navigation receiver, an indoor beacon system, and an
acoustic-tracking rig all solve a cousin of that problem: they estimate
a position from measured distances to known anchors, and the quality of
the estimate depends less on the precision of any single range than on
the geometry of the anchors. That dependence is a condition number, and
naming it is the point of the vignette.

\paragraph{The measurement model.}
An anchor at known location $\vect{p}_{k} \in \R^{3}$ reports its
distance $d_{k}$ to the unknown receiver position $\vect{x}$, so
$\norm{\vect{x} - \vect{p}_{k}}_{2} = d_{k}$. The equation is nonlinear
in $\vect{x}$ because of the square root inside the norm, so the working
method linearises about a current guess $\vect{x}_{0}$. The gradient of
the range with respect to position is the unit vector from the anchor to
the receiver, $\vect{e}_{k} = (\vect{x}_{0} - \vect{p}_{k})/
\norm{\vect{x}_{0} - \vect{p}_{k}}_{2}$, so a small position correction
$\Delta\vect{x}$ changes the modelled range by $\vect{e}_{k}^{\transpose}
\Delta\vect{x}$. Stacking the $m$ anchors gives the overdetermined
linear least-squares system $\mat{A}\,\Delta\vect{x} = \Delta\vect{d}$,
where row $k$ of $\mat{A}$ is the line-of-sight unit vector
$\vect{e}_{k}^{\transpose}$ and $\Delta\vect{d}$ is the vector of
measured-minus-modelled ranges. The receiver refines $\vect{x}_{0}
\leftarrow \vect{x}_{0} + \Delta\vect{x}$ and repeats, the Gauss--Newton
iteration of Chapter 11 running on the least-squares apparatus of this
section.

\paragraph{The design matrix is a matrix of directions.}
Every row of $\mat{A}$ is a unit vector pointing from an anchor to the
receiver. The conditioning of the fit is therefore set entirely by how
those directions are spread. When the anchors surround the receiver,
the rows point in well-separated directions, $\mat{A}^{\transpose}
\mat{A}$ is close to a scaled identity, and the condition number is near
one. When the anchors cluster in one part of the sky, the rows are
nearly parallel, $\mat{A}^{\transpose}\mat{A}$ is nearly rank
deficient, and the condition number is large. The navigation community
has a name for the square root of the relevant diagonal of
$(\mat{A}^{\transpose}\mat{A})^{-1}$, the \engterm{dilution of
precision}, and it is the condition number of this section wearing a
domain hat: it multiplies the ranging error into the position error
exactly as the sensitivity bound of section 9.7 says a condition number
must. Figure~\ref{fig:vol02:ch09:trilateration-gdop} contrasts a
well-spread constellation with a clustered one.

\paragraph{A worked estimate.}
Take four anchors and a receiver, with the line-of-sight unit vectors
(after linearising) coming out as $\vect{e}_{1} = (1, 0, 0)$,
$\vect{e}_{2} = (0, 1, 0)$, $\vect{e}_{3} = (0, 0, 1)$, and
$\vect{e}_{4} = (1, 1, 1)/\sqrt{3}$. The design matrix stacks these
four rows, and the Gram matrix is
\[
\mat{A}^{\transpose}\mat{A} = \begin{pmatrix}
4/3 & 1/3 & 1/3 \\
1/3 & 4/3 & 1/3 \\
1/3 & 1/3 & 4/3
\end{pmatrix},
\]
whose eigenvalues are $2$ (once) and $1$ (twice), so $\kappa(\mat{A}^{
\transpose}\mat{A}) = 2$ and $\kappa(\mat{A}) = \sqrt{2}$: a
well-conditioned fix, and a ranging error of $\SI{1}{\meter}$ maps to a
position error of about a metre. Now cluster the anchors so that the
four unit vectors all lie within a narrow cone, say $\vect{e}_{k}
\approx (0, 0, 1) + \vect{\delta}_{k}$ with the $\vect{\delta}_{k}$
small and mostly horizontal. The vertical direction is now measured by
every row and the two horizontal directions by almost none, so the
smallest eigenvalue of $\mat{A}^{\transpose}\mat{A}$ collapses toward
zero while the largest stays near $m$, and $\kappa$ climbs past
$10^{3}$. The same $\SI{1}{\meter}$ ranging error now maps to a
horizontal position error of tens of metres, with the receiver pinned
in altitude and adrift in the plane, the elongated error region of the
figure's right panel.

\paragraph{What the vignette teaches.}
The precision of a position fix is a property of geometry, not of the
ranging hardware alone. A receiver reporting a position without its
dilution of precision is reporting a least-squares estimate without its
condition number, the incompleteness section 9.8 refuses. The remedy is
also geometric: add an anchor that observes the poorly conditioned
direction, exactly the row an engineer would add to a rank-deficient
design matrix to restore full column rank. The trilateration fit, the
levelling adjustment, and the resistor solve are three faces of the one
least-squares template, and the condition number is the single number
that tells the engineer whether the geometry supports the answer.

\begin{figure}[ht]
\centering
\begin{tikzpicture}[x=1.0cm, y=1.0cm, font=\small, >=Stealth]

% Left panel: strong geometry (well spread beacons)
\begin{scope}[shift={(0,0)}]
  \node[font=\scriptsize, anchor=south] at (2.0, 3.3) {well-spread beacons};
  % receiver
  \fill[engblack] (2.0, 1.5) circle (1.4pt);
  \node[anchor=north, font=\scriptsize] at (2.0, 1.4) {$\vect{x}$};
  % three beacons around
  \foreach \bx/\by in {0.2/0.4, 3.8/0.5, 2.0/3.0}{
    \fill[engblue] (\bx, \by) circle (1.4pt);
    \draw[engblue, thick] (2.0, 1.5) -- (\bx, \by);
  }
  % small well-conditioned intersection region
  \fill[englime!40] (2.0, 1.5) ellipse (0.18 and 0.16);
\end{scope}

% Right panel: poor geometry (clustered beacons)
\begin{scope}[shift={(6.5,0)}]
  \node[font=\scriptsize, anchor=south] at (2.0, 3.3) {clustered beacons};
  % receiver
  \fill[engblack] (2.0, 1.5) circle (1.4pt);
  \node[anchor=north, font=\scriptsize] at (2.0, 1.4) {$\vect{x}$};
  % three beacons nearly in a line, one direction
  \foreach \bx/\by in {1.4/3.0, 2.0/3.0, 2.6/3.0}{
    \fill[engred] (\bx, \by) circle (1.4pt);
    \draw[engred, thick] (2.0, 1.5) -- (\bx, \by);
  }
  % elongated ill-conditioned region
  \fill[engamber!35] (2.0, 1.5) ellipse (0.7 and 0.16);
\end{scope}

\node[anchor=north, font=\scriptsize, text width=12cm, align=center]
  at (4.5, -0.5) {
  The same range measurements pin the receiver tightly when the
  beacons surround it (left) and loosely when they cluster in one
  direction (right). The stretched error region on the right is the
  geometric face of a large condition number: the dilution of
  precision.
};

\end{tikzpicture}
\caption[Geometry sets the conditioning of trilateration]{The
conditioning of a trilateration fit is set by the geometry of the
beacons, not by the precision of the individual ranges. Beacons that
surround the receiver (left) give a design matrix of rows pointing in
well-separated directions, a small condition number, and a compact
error region. Beacons clustered in one direction (right) give nearly
parallel rows, a large condition number, and an error region stretched
along the poorly observed direction. This ratio is what navigation
engineers call the geometric dilution of precision, and it is the
condition number of the linearised design matrix in disguise.}
\label{fig:vol02:ch09:trilateration-gdop}
\end{figure}


\subsection{Worked case study: a Leontief input-output economy}\pathtag{standard}

A linear system need not come from geometry or measurement. The
Leontief input-output model, for which Wassily Leontief received the
1973 economics Nobel, treats an economy as a set of sectors that
consume one another's output to produce their own, and asks how much
each sector must produce to meet a given final demand. The model is a
square linear system whose matrix carries an economic interpretation
the engineer can read directly, and it introduces the Neumann series,
the linear-algebra tool behind iterative solvers and behind the
analysis of feedback in Volume IX.

\paragraph{The technical-coefficient matrix.}
Number the sectors $1, \ldots, n$. The \engterm{technical coefficient}
$a_{ij}$ is the amount of sector $i$'s output consumed in producing
one unit of sector $j$'s output. The matrix $\mat{A} = (a_{ij})$ is
the technology of the economy. If sector $j$ produces $x_{j}$ units,
it consumes $a_{ij}x_{j}$ units of sector $i$, so the total
intermediate consumption of sector $i$ is $\sum_{j} a_{ij}x_{j} =
(\mat{A}\vect{x})_{i}$. The accounting identity is that each sector's
output equals what other sectors consume of it plus the final demand
$d_{i}$ (households, exports, investment):
\[
\vect{x} = \mat{A}\vect{x} + \vect{d}, \qquad\text{so}\qquad
(\mat{I} - \mat{A})\vect{x} = \vect{d}.
\]
The matrix $\mat{I} - \mat{A}$ is the \engterm{Leontief matrix}, and
the gross output that meets demand $\vect{d}$ is $\vect{x} = (\mat{I}
- \mat{A})^{-1}\vect{d}$, with the inverse computed never explicitly
but through a single LU solve.

\paragraph{A three-sector economy.}
Take the economy of Figure~\ref{fig:vol02:ch09:leontief} with sectors
agriculture, manufacturing, and energy, and technical coefficients
\[
\mat{A} = \begin{pmatrix}
0.2 & 0.2 & 0.3 \\
0.1 & 0.0 & 0.2 \\
0.0 & 0.1 & 0.1
\end{pmatrix}.
\]
Column $j$ reads as the recipe for one unit of sector $j$: producing
one unit of manufacturing (column 2) consumes $0.2$ of agriculture and
$0.1$ of energy. The column sums, $0.3$, $0.3$, $0.6$, are all below
one, the sign that each sector delivers net output rather than
consuming more than it makes. The Leontief matrix is
\[
\mat{I} - \mat{A} = \begin{pmatrix}
0.8 & -0.2 & -0.3 \\
-0.1 & 1.0 & -0.2 \\
0.0 & -0.1 & 0.9
\end{pmatrix}.
\]

\paragraph{The solve.}
For a final demand $\vect{d} = (100, 50, 30)^{\transpose}$ (in
monetary units), Gaussian elimination on $(\mat{I} - \mat{A})\vect{x}
= \vect{d}$ returns the gross outputs $\vect{x} \approx (164.6, 88.6,
43.2)^{\transpose}$. Each sector must produce more than its final
demand because part of every sector's output is swallowed as
intermediate input to the others: agriculture produces about $165$ to
deliver $100$ to final demand, the extra $65$ feeding manufacturing
and energy. The off-diagonal couplings are the economy's supply chain,
and the inverse $(\mat{I} - \mat{A})^{-1}$ records the total (direct
plus indirect) requirement of each sector per unit of each final
demand.

\paragraph{Productivity and the Neumann series.}
The model has a non-negative, finite solution for every non-negative
demand exactly when the economy is \engterm{productive}, and the
algebraic condition is that the spectral radius of $\mat{A}$ (its
largest eigenvalue magnitude) is below one. When it is, the inverse
expands as the convergent \engterm{Neumann series}
\[
(\mat{I} - \mat{A})^{-1} = \mat{I} + \mat{A} + \mat{A}^{2} +
\mat{A}^{3} + \cdots,
\]
the matrix analogue of the geometric series $1/(1 - a) = 1 + a +
a^{2} + \cdots$ for $\abs{a} < 1$. The series has a direct economic
reading: $\mat{I}$ is the demand itself, $\mat{A}\vect{d}$ is the
first round of intermediate inputs to make that demand, $\mat{A}^{2}
\vect{d}$ is the inputs to make those inputs, and so on. The
convergence of the series is the statement that the supply chain
terminates: each round of indirect demand is smaller than the last.
For the matrix above, the spectral radius is about $0.45$, comfortably
below one, so the series converges geometrically and the first few
terms already approximate the gross output. The same Neumann series is
how iterative solvers (the stationary methods of section 9.7) generate
their iterates, and how a feedback loop's closed-loop response is
expanded in Volume IX; the economic model and the control loop are the
same algebra.

\begin{example}[Low-rank approximation as image compression]
The truncated SVD of section 9.7 has a direct application that makes
the rank concept concrete: compression. A grayscale image is a matrix
$\mat{M} \in \R^{m \times n}$ of pixel intensities. Its SVD $\mat{M}
= \sum_{i} \sigma_{i}\vect{u}_{i}\vect{v}_{i}^{\transpose}$ writes the
image as a sum of rank-one layers, ordered by importance through the
singular values. Keeping the first $r$ layers gives the rank-$r$
approximation $\mat{M}_{r} = \sum_{i=1}^{r} \sigma_{i}\vect{u}_{i}
\vect{v}_{i}^{\transpose}$, and the \engterm{Eckart--Young theorem}
certifies that no other rank-$r$ matrix is closer to $\mat{M}$ in
either the 2-norm or the Frobenius norm
\cite{paper:v2ch10-eckart-young1936}. The error of the truncation is
exactly the first discarded singular value, $\norm{\mat{M} -
\mat{M}_{r}}_{2} = \sigma_{r+1}$. The storage drops from $mn$ numbers
to $r(m + n + 1)$ numbers, the $r$ singular values plus their left and
right vectors. For a $1000 \times 1000$ image whose singular values
decay quickly (Figure~\ref{fig:vol02:ch09:low-rank}), a rank-$50$
truncation stores $50 \times 2001 \approx 10^{5}$ numbers against the
original $10^{6}$, a tenfold compression, while the discarded tail
$\sigma_{51}, \sigma_{52}, \ldots$ is small enough that the eye cannot
distinguish the result. The same truncation is principal-component
analysis when $\mat{M}$ is a data matrix, model-order reduction when
$\mat{M}$ is a state-transition operator, and noise filtering when the
discarded tail is measurement noise. The rank of a matrix, an abstract
count of independent directions, becomes a compression ratio.
\end{example}

\begin{figure}[ht]
\centering
\begin{tikzpicture}[x=1.0cm, y=1.0cm, font=\small,
  sector/.style={rectangle, rounded corners=2pt, draw, thick,
    fill=engblue!12, minimum width=22mm, minimum height=10mm,
    align=center, inner sep=2pt}]

\node[sector] (agr) at (0, 0) {Agriculture};
\node[sector] (man) at (5, 0) {Manufacturing};
\node[sector] (ene) at (2.5, -3) {Energy};

% Inter-industry flows (a_ij = input from i needed per unit output of j)
\draw[->, thick, >=stealth] (agr) to[bend left=12]
  node[midway, above, font=\scriptsize] {$0.1$} (man);
\draw[->, thick, >=stealth] (man) to[bend left=12]
  node[midway, below, font=\scriptsize] {$0.2$} (agr);
\draw[->, thick, >=stealth] (ene) to[bend left=10]
  node[midway, left=1pt, font=\scriptsize] {$0.3$} (agr);
\draw[->, thick, >=stealth] (ene) to[bend right=10]
  node[midway, right=1pt, font=\scriptsize] {$0.2$} (man);
\draw[->, thick, >=stealth] (man) to[bend left=10]
  node[midway, right=1pt, font=\scriptsize] {$0.1$} (ene);
\draw[->, thick, engred, >=stealth] (agr) edge[loop above]
  node[font=\scriptsize] {$0.2$} (agr);

% Final demand arrows out
\draw[->, very thick, enggray, >=stealth] (agr) -- ++(-2.0, 0.0)
  node[left, font=\scriptsize] {demand $d_1$};
\draw[->, very thick, enggray, >=stealth] (man) -- ++(2.0, 0.0)
  node[right, font=\scriptsize] {demand $d_2$};
\draw[->, very thick, enggray, >=stealth] (ene) -- ++(0.0, -1.6)
  node[below, font=\scriptsize] {demand $d_3$};

\node[anchor=north, font=\scriptsize, text width=10.5cm, align=center]
  at (2.5, -5.2) {
  A three-sector Leontief economy. An arrow labelled $a_{ij}$ from
  sector $i$ to sector $j$ means producing one unit of sector $j$'s
  output consumes $a_{ij}$ units of sector $i$'s output. The
  technical-coefficient matrix $\mat{A}$ collects these labels; the
  total output $\vect{x}$ that meets a final demand $\vect{d}$ solves
  $(\mat{I} - \mat{A})\vect{x} = \vect{d}$.};

\end{tikzpicture}
\caption[A Leontief input-output economy as a linear
system]{A three-sector input-output economy. Each sector's output is
consumed partly by the other sectors as intermediate input and partly
by final demand. The technical-coefficient matrix $\mat{A}$ records
the inter-industry requirements, and the gross output that meets a
final demand $\vect{d}$ solves the Leontief system $(\mat{I} -
\mat{A})\vect{x} = \vect{d}$. The economy is productive (the system
has a non-negative solution for every non-negative demand) exactly
when the spectral radius of $\mat{A}$ is below one, which is the
convergence condition for the Neumann series $(\mat{I} - \mat{A})^{-1}
= \sum_{k \geq 0} \mat{A}^{k}$.}
\label{fig:vol02:ch09:leontief}
\end{figure}


\begin{figure}[ht]
\centering
\begin{tikzpicture}[x=1.0cm, y=1.0cm, font=\small]

% Decaying singular-value spectrum on a log scale (illustrative).
\def\axw{8.2}
\def\axh{4.2}
\draw[->, thick] (0,0) -- (\axw, 0) node[right, font=\scriptsize] {rank index $k$};
\draw[->, thick] (0,0) -- (0, \axh) node[above, font=\scriptsize] {$\log_{10}\sigma_k$};

% gridline ticks on y
\foreach \y/\lab in {0/-3, 1/-2, 2/-1, 3/0, 4/1}
  {\draw[enggray!50] (0, \y) -- (\axw, \y);
   \node[anchor=east, font=\scriptsize] at (-0.1, \y) {\lab};}

% A decaying spectrum: sigma_k roughly 10^{1 - 0.18 k}, levelling into noise.
\draw[very thick, engblue]
  plot[smooth] coordinates {
    (0.3, 4.0) (1.0, 3.55) (1.8, 3.05) (2.6, 2.55) (3.4, 2.05)
    (4.2, 1.55) (5.0, 1.15) (5.8, 0.85) (6.6, 0.6) (7.4, 0.5) (8.0, 0.48)};

% Truncation cutoff
\draw[dashed, very thick, engred] (4.0, 0) -- (4.0, \axh);
\node[engred, anchor=south, font=\scriptsize, align=center]
  at (4.0, \axh) {truncation rank $r$};
\node[anchor=west, font=\scriptsize, text width=3.5cm] at (4.3, 2.9)
  {keep $\sigma_1,\ldots,\sigma_r$:\\ the signal subspace};
\node[anchor=west, font=\scriptsize, text width=3.2cm] at (4.3, 0.85)
  {discard the tail:\\ near-noise directions};

\node[anchor=north, font=\scriptsize, text width=10.5cm, align=center]
  at (4.0, -0.7) {
  A typical singular-value spectrum of a structured matrix (a natural
  image, a covariance matrix, a discretised smooth operator) decays
  rapidly, then levels into a noise floor. Truncating at rank $r$
  keeps the dominant directions and discards the rest. The
  Eckart--Young theorem certifies that this truncation is the best
  rank-$r$ approximation in both the 2-norm and the Frobenius norm.};

\end{tikzpicture}
\caption[Singular-value decay and low-rank truncation]{The
singular-value spectrum of a structured matrix. Rapid decay means the
matrix is well approximated by a low-rank truncation: keeping the
first $r$ singular triplets reproduces the matrix to an error equal
to the first discarded singular value $\sigma_{r+1}$ (in the 2-norm).
The Eckart--Young theorem makes the truncated SVD the optimal
rank-$r$ approximation, which is the algebraic content of image
compression, principal-component analysis, and model-order reduction.}
\label{fig:vol02:ch09:low-rank}
\end{figure}


\section{Conditioning and numerical sensitivity}\pathtag{core}

A linear system $\mat{A}\vect{x} = \vect{b}$ has a unique solution
when $\mat{A}$ is invertible. Whether floating-point arithmetic
returns that solution accurately depends on the conditioning of
$\mat{A}$, a quantity that measures how much small perturbations
in $\mat{A}$ or $\vect{b}$ amplify into changes in $\vect{x}$.

\subsection{The condition number}

For an invertible square matrix $\mat{A}$, the \engterm{condition
number} (in the 2-norm) is

\[
\kappa(\mat{A}) = \norm{\mat{A}}_{2} \cdot \norm{\mat{A}^{-1}}_{2}
= \frac{\sigma_{1}(\mat{A})}{\sigma_{n}(\mat{A})},
\]

where $\sigma_{1} \geq \cdots \geq \sigma_{n} > 0$ are the singular
values of $\mat{A}$ \cite[ch.~12]{text:trefethen-bau1997}. The
ratio of largest to smallest singular value answers the question
\emph{how much can $\mat{A}$ stretch some directions while it
shrinks others}.

The standard sensitivity bound, valid for any consistent norm, is

\[
\frac{\norm{\Delta\vect{x}}}{\norm{\vect{x}}} \leq \kappa(\mat{A})
\left(\frac{\norm{\Delta\mat{A}}}{\norm{\mat{A}}}
+ \frac{\norm{\Delta\vect{b}}}{\norm{\vect{b}}}\right) + O(\epsilon^{2}),
\]

where $\Delta\mat{A}$ and $\Delta\vect{b}$ are perturbations and
$\Delta\vect{x}$ is the resulting change in the solution. The
inequality reads as a budget. A relative perturbation of size
$\epsilon$ in the data can produce a relative change of size up to
$\kappa(\mat{A}) \cdot \epsilon$ in the solution. When
$\kappa(\mat{A}) = 10^{6}$ and the data are accurate to $10^{-8}$
(double precision near the machine epsilon $2.2 \times 10^{-16}$
times typical magnitude), the solution is accurate to about
$10^{-2}$.

\begin{example}[Computing a condition number by hand]
Take $\mat{A} = \begin{psmallmatrix} 2 & 0 \\ 0 & 0.01
\end{psmallmatrix}$, already diagonal so its singular values are the
absolute diagonal entries, $\sigma_{1} = 2$ and $\sigma_{2} = 0.01$.
The condition number is $\kappa(\mat{A}) = 2/0.01 = 200$. The
geometric reading: $\mat{A}$ stretches the first axis by $2$ and the
second by $0.01$, an aspect-ratio distortion of $200$ to $1$. A
right-hand side aligned with the second axis produces a solution
$200$ times more sensitive to perturbation than one aligned with the
first. For a non-diagonal matrix the singular values come from the
eigenvalues of $\mat{A}^{\transpose}\mat{A}$: for $\mat{A} =
\begin{psmallmatrix} 1 & 2 \\ 0 & 1 \end{psmallmatrix}$, $\mat{A}^{
\transpose}\mat{A} = \begin{psmallmatrix} 1 & 2 \\ 2 & 5
\end{psmallmatrix}$ has eigenvalues $3 \pm 2\sqrt{2}$, so $\sigma_{1}
= \sqrt{3 + 2\sqrt{2}} \approx 2.414$, $\sigma_{2} = \sqrt{3 -
2\sqrt{2}} \approx 0.414$, and $\kappa \approx 5.83$. The shear,
mild as it looks, already carries a condition number near six.
\end{example}

\begin{example}[A small residual is not a small error]
The most dangerous misreading of a computed solution is to trust it
because its residual is small. Take the ill-conditioned system
$\mat{A} = \begin{psmallmatrix} 1 & 1 \\ 1 & 1.0001
\end{psmallmatrix}$, $\vect{b} = (2, 2.0001)^{\transpose}$, whose
exact solution is $\vect{x} = (1, 1)^{\transpose}$. Consider the
candidate $\hat{\vect{x}} = (2, 0)^{\transpose}$, wrong in both
components. Its residual is $\vect{r} = \vect{b} - \mat{A}
\hat{\vect{x}} = (2, 2.0001) - (2, 2) = (0, 0.0001)^{\transpose}$,
with norm $10^{-4}$, reassuringly tiny against $\norm{\vect{b}}
\approx 2.83$. Yet the error $\hat{\vect{x}} - \vect{x} = (1,
-1)^{\transpose}$ has norm $\sqrt{2}$, a $100\%$ relative error. The
residual and the error are related through the inverse, $\vect{e} =
\mat{A}^{-1}\vect{r}$, so the error can be as large as
$\norm{\mat{A}^{-1}}$ times the residual, and for this matrix
$\norm{\mat{A}^{-1}}_{2} \approx 2 \times 10^{4}$. The conditioning
is the bridge: $\norm{\vect{e}}/\norm{\vect{x}} \leq \kappa(\mat{A})
\cdot \norm{\vect{r}}/\norm{\vect{b}}$, and with $\kappa \approx 4
\times 10^{4}$ a residual of $10^{-4}/2.83 \approx 3.5 \times
10^{-5}$ permits a relative error up to $\approx 1.4$, consistent
with the $100\%$ error observed. The lesson is operative: on an
ill-conditioned system, verify the solution by an error estimate
(an inverse-norm bound, a condition estimate, or iterative
refinement), never by the residual alone.
\end{example}

\begin{example}[An innocuous-looking triangular matrix that is ill-conditioned]
A matrix can carry a large condition number with no small entry, no
zero pivot, and a determinant of exactly one, which is why the
determinant is the wrong diagnostic. Take the upper-triangular matrix
with ones on the diagonal and $-1$ everywhere above it,
\[
\mat{A} = \begin{pmatrix} 1 & -1 & -1 & -1 \\ 0 & 1 & -1 & -1 \\ 0 & 0
& 1 & -1 \\ 0 & 0 & 0 & 1 \end{pmatrix}.
\]
The determinant is the product of the diagonal, $\det\mat{A} = 1$, and
every entry is $\pm 1$, so nothing on the face of the matrix looks
dangerous. The inverse tells the real story: back-substitution
propagates the $-1$ entries into powers of two, and
\[
\mat{A}^{-1} = \begin{pmatrix} 1 & 1 & 2 & 4 \\ 0 & 1 & 1 & 2 \\ 0 & 0
& 1 & 1 \\ 0 & 0 & 0 & 1 \end{pmatrix},
\]
whose largest entry is $2^{n-2} = 4$ at $n = 4$. The infinity-norm of
$\mat{A}$ is the largest absolute row sum, $\norm{\mat{A}}_{\infty} =
4$, and of the inverse, $\norm{\mat{A}^{-1}}_{\infty} = 8$, so
$\kappa_{\infty}(\mat{A}) = 32$ already at $n = 4$. Grow the matrix to
$n = 30$ and the top-right entry of the inverse is $2^{28} \approx 2.7
\times 10^{8}$, so $\kappa$ passes $10^{9}$ and a solve loses nine
digits, all from a matrix of $\pm 1$ entries with unit determinant.
This is the standard warning against reading conditioning off the
determinant or off the entry magnitudes: the growth lives in the
inverse, where back-substitution accumulates it, and only the condition
number sees it. The example also explains why partial pivoting, which
bounds the multipliers by one, still cannot bound the growth factor in
the worst case: this matrix is exactly the pattern on which Gaussian
elimination's growth factor reaches its $2^{n-1}$ theoretical maximum.
\end{example}

\subsection{The rule of thumb}

The working rule of thumb: a matrix with condition number
$\kappa$ loses about $\log_{10}\kappa$ decimal digits of accuracy
in the solution relative to the data. Double-precision floating
point carries about 16 decimal digits; a condition number of
$10^{8}$ leaves about 8 digits in the solution, a condition number
of $10^{12}$ leaves about 4, and a condition number above $10^{16}$
leaves nothing.

The rule is loose. Backward stability of the algorithm, structure
in the right-hand side, and the specific norm in use all change
the constant. The rule's value is as a triage: the engineer who
sees $\kappa = 10^{14}$ knows, before any computation, that the
problem in this formulation has at most one or two reliable
digits. Figure~\ref{fig:vol02:ch09:condition-number} encodes the
budget across the three working IEEE precisions.

\begin{figure}[ht]
\centering
\begin{tikzpicture}[x=1.0cm, y=0.8cm, font=\small]

% Axes
\draw[->, thick] (-0.4, 0) -- (10.4, 0) node[anchor=west] {$\log_{10}\kappa(\mat{A})$};
\draw[->, thick] (0, -0.4) -- (0, 8.4) node[anchor=south] {decimal digits retained};

% Ticks x
\foreach \x in {0,2,4,6,8,10} {
  \draw (\x, 0.1) -- (\x, -0.1) node[anchor=north, font=\scriptsize] {\x};
}
% Ticks y
\foreach \y in {0,2,4,6,8} {
  \draw (0.1, \y) -- (-0.1, \y) node[anchor=east, font=\scriptsize] {\y};
}

% Double precision: y = 16 - x (clipped at 0)
\draw[very thick, engblue] (0, 8) -- (8, 0) -- (10, 0);
\node[engblue, anchor=south west] at (3.6, 4.2) {double ($16 - \log_{10}\kappa$)};

% Single precision: y = 7 - x
\draw[very thick, engred] (0, 7) -- (7, 0) -- (10, 0);
\node[engred, anchor=south west] at (0.2, 6.3) {single ($7 - \log_{10}\kappa$)};

% danger zones
\draw[dashed, enggray] (8, 0) -- (8, 7.6);
\node[enggray, font=\scriptsize, anchor=south, align=center] at (8, 7.7) {refuse normal\\ equations};

\draw[dashed, enggray] (15, 0) -- (15, 7.6);

% Half-precision
\draw[very thick, engamber] (0, 3.3) -- (3.3, 0) -- (10, 0);
\node[engamber, anchor=south west, font=\scriptsize] at (0.2, 2.7) {half ($3.3 - \log_{10}\kappa$)};

% Annotation
\node[font=\scriptsize, text width=8.5cm, align=center, anchor=north] at (5, -1.4) {
  Rule of thumb: a condition number $\kappa$ costs $\log_{10}\kappa$ decimal digits
  in the solution. Floor at zero retained digits is the algebraic answer the algorithm
  must refuse to report.
};

\end{tikzpicture}
\caption[Decimal-digit budget against condition number]{Working
rule-of-thumb chart for condition-number budgeting. Each precision
gives a fixed total digit budget (about 16 in double, 7 in single,
3.3 in half); each factor of $10$ in $\kappa(\mat{A})$ removes one
digit from that budget. Above $\kappa \approx 10^{8}$ the normal
equations formulation is to be refused in double precision (the
Gram matrix doubles the loss); above $\kappa \approx 10^{15}$ even
QR has lost almost everything and the engineer must reformulate. The
figure is a graphical encoding of the budget inequality of section
9.7.}
\label{fig:vol02:ch09:condition-number}
\end{figure}


\begin{mastery}
The condition number is a property of the \emph{problem}; backward
stability is a property of the \emph{algorithm}. The two combine in
the central error estimate of numerical linear algebra. A backward
stable algorithm returns the exact solution of a nearby problem,
$(\mat{A} + \Delta\mat{A})\hat{\vect{x}} = \vect{b} + \Delta\vect{b}$
with $\norm{\Delta\mat{A}}/\norm{\mat{A}}$ and $\norm{\Delta\vect{b}}/
\norm{\vect{b}}$ both of order the machine epsilon $\epsilon_{\text{mach}}$.
The forward error in $\hat{\vect{x}}$ is then bounded by the product
of the two,
\[
\frac{\norm{\hat{\vect{x}} - \vect{x}}}{\norm{\vect{x}}} \lesssim
\kappa(\mat{A})\,\epsilon_{\text{mach}},
\]
the relation Trefethen and Bau call ``forward error $\lesssim$
condition number $\times$ backward error''
\cite[ch.~15]{text:trefethen-bau1997}. Gaussian elimination with
partial pivoting, Householder QR, and Cholesky are all backward
stable; classical Gram-Schmidt is not, which is why its output can
lose orthogonality even on a well-conditioned matrix. The engineer's
control is split accordingly: choose a backward stable algorithm to
keep the backward error at $\epsilon_{\text{mach}}$, and reformulate
the problem (rescale, change basis, regularise) to keep
$\kappa(\mat{A})$ small. Neither alone suffices.
\end{mastery}

\begin{estimation}
\textbf{Question.} A daily-mean temperature record at a single
weather station, $m = 365$ points over one year, is to be fitted by
a polynomial of degree $p$ in the monomial basis on the day index
$x \in \{1, \ldots, 365\}$, in double precision. At what degree $p$
does the rule of thumb predict that the normal equations have lost
all reliable digits in the fitted coefficients, assuming the data
are accurate to about three significant figures?

\smallskip
\textit{Estimate, before reading on.}

\smallskip
Vandermonde conditioning on equispaced points scales roughly as
$\log_{10}\kappa(\mat{V}) \approx 1.5 p$ in the monomial basis on
an interval that does not include the origin. Squaring for the Gram
matrix gives $\log_{10}\kappa(\mat{V}^{\transpose}\mat{V}) \approx
3 p$. The data carry three digits; double precision carries sixteen;
the working budget is about $13$ digits. The rule of thumb says
$3p \approx 13$, i.e.\ $p \approx 4$ to $5$, exhausts the budget.
The script \path{code/condition_visualiser.py} confirms the
estimate to within a factor of three: at $p = 5$ the Gram condition
number is already above $10^{14}$ on $[0, 1]$, leaving at most one
or two coefficient digits. Rescaling $x$ to $[-1, 1]$ buys back
about three orders of magnitude per degree; the Chebyshev basis
buys back the rest. The estimate is the triage that tells the
engineer to rescale or change basis before the fit, not after.
\end{estimation}

\subsection{The condition number of the Gram matrix}

For $\mat{A} \in \R^{m \times n}$ with full column rank, the Gram
matrix $\mat{A}^{\transpose}\mat{A}$ has condition number
$\kappa(\mat{A}^{\transpose}\mat{A}) = \kappa(\mat{A})^{2}$. The
squaring is the central numerical fact about the normal equations.
A design matrix with $\kappa(\mat{A}) = 10^{4}$ produces a Gram
matrix with $\kappa(\mat{A}^{\transpose}\mat{A}) = 10^{8}$. Solving
the normal equations directly carries the squared conditioning;
solving via QR (section 9.6) carries the original conditioning.

This is why a polynomial fit of moderate degree on closely-spaced
data is the canonical example of when not to use the normal
equations. The Vandermonde columns correlate strongly, the matrix
is ill-conditioned, and the squared conditioning of the Gram matrix
makes the normal-equations solution numerically meaningless before
the engineer notices.

\subsection{Cheap bounds: Gershgorin discs and diagonal dominance}

The condition number is the exact diagnostic, and it costs a
singular-value decomposition to compute. Before paying that price, the
engineer can bound where the spectrum lies from the matrix entries
alone. The \engterm{Gershgorin disc theorem} states that every
eigenvalue of $\mat{A}$ lies inside at least one of the $n$ discs in
the complex plane, the $i$-th disc centred on the diagonal entry
$a_{ii}$ with radius the sum of the off-diagonal magnitudes in that row,
$R_{i} = \sum_{j \neq i} \abs{a_{ij}}$. The proof is one line: if
$\mat{A}\vect{x} = \lambda\vect{x}$ and $x_{i}$ is the largest component
in magnitude, the $i$-th equation rearranges to $(\lambda - a_{ii})x_{i}
= \sum_{j \neq i} a_{ij}x_{j}$, and dividing by $x_{i}$ and taking
magnitudes gives $\abs{\lambda - a_{ii}} \leq R_{i}$.
Figure~\ref{fig:vol02:ch09:gershgorin} draws the discs for a small
matrix.

The theorem earns its keep as an invertibility test that costs only a
scan of the rows. When no disc reaches the origin, no eigenvalue can be
zero, so the matrix is nonsingular. The condition ``no disc contains the
origin'' is exactly $\abs{a_{ii}} > R_{i}$ for every row, which is
\engterm{strict diagonal dominance}. A strictly diagonally dominant
matrix is invertible, needs no pivoting in Gaussian elimination, and, if
symmetric with positive diagonal, is positive definite. The stiffness
and conductance matrices of the network examples are diagonally
dominant by construction, which is why the earlier examples could reach
for Cholesky without a pivoting discussion. Diagonal dominance is the
structural property an engineer checks first, because it settles
solvability and stability together at the cost of a single pass over the
matrix.

\begin{figure}[ht]
\centering
\begin{tikzpicture}[x=1.0cm, y=1.0cm, font=\small]

% Axes (complex plane)
\draw[->, thick] (-1.0, 0) -- (9.5, 0) node[anchor=west, font=\scriptsize] {real};
\draw[->, thick] (4.0, -3.0) -- (4.0, 3.0) node[anchor=south, font=\scriptsize] {imag};

% Disc centre 2 (shifted to x=2 on plot from origin at x=4 => plot x=2? we place plot origin at 4)
% Map: plot_x = 4 + value. Centres at 2, 5, 8 -> plot 6, 9? keep small. Use centres 2,5 and radii 1, 1.5
% disc 1: centre 2 (plot 4+? ) simpler: put real origin at plot x=4, unit=1cm
% centre a11 = 2 -> plot (6,0), radius 1
\draw[thick, engblue, fill=engblue!12, fill opacity=0.5] (6, 0) circle (1);
\node[engblue, font=\scriptsize] at (6, 1.3) {$D_{1}$};
\fill[engblue] (6, 0) circle (1.4pt);
\node[engblue, font=\scriptsize, anchor=north] at (6, -0.15) {$2$};

% centre a22 = -1 -> plot (3,0), radius 0.5
\draw[thick, engred, fill=engred!12, fill opacity=0.5] (3, 0) circle (0.5);
\node[engred, font=\scriptsize] at (3, 0.85) {$D_{2}$};
\fill[engred] (3, 0) circle (1.4pt);
\node[engred, font=\scriptsize, anchor=north] at (3, -0.15) {$-1$};

% centre a33 = 5 -> plot (9,0) too far; use 4.5 -> plot (8.5,0) radius 0.8
\draw[thick, engamber, fill=engamber!12, fill opacity=0.5] (8.5, 0) circle (0.8);
\node[engamber, font=\scriptsize] at (8.5, 1.1) {$D_{3}$};
\fill[engamber] (8.5, 0) circle (1.4pt);
\node[engamber, font=\scriptsize, anchor=north] at (8.5, -0.15) {$4.5$};

% Real axis origin tick
\draw (4, 0.08) -- (4, -0.08) node[anchor=north east, font=\scriptsize] {$0$};

\node[anchor=north, font=\scriptsize, text width=11cm, align=center]
  at (4.0, -3.3) {
  Each Gershgorin disc is centred on a diagonal entry with radius the
  sum of the off-diagonal magnitudes in that row. Every eigenvalue
  lies in the union of the discs. No disc here contains the origin, so
  the matrix is nonsingular and its eigenvalues stay clear of zero, a
  read-off of invertibility that costs only a scan of the rows.
};

\end{tikzpicture}
\caption[Gershgorin discs localise the spectrum]{Gershgorin discs for
a $3 \times 3$ matrix, one disc per row, centred on the diagonal entry
$a_{ii}$ with radius $\sum_{j \neq i} \abs{a_{ij}}$. The theorem
guarantees every eigenvalue sits inside the union. When no disc
reaches the origin, as here, the matrix cannot have a zero eigenvalue
and is therefore invertible; when the discs are small and well
separated from zero, the matrix is well away from singularity. The
disc scan is the cheapest available bound on where a matrix's spectrum
can lie, and a strictly diagonally dominant matrix is exactly one
whose discs all miss the origin.}
\label{fig:vol02:ch09:gershgorin}
\end{figure}


\begin{example}[Gershgorin invertibility and a dominance check]
Take $\mat{A} = \begin{psmallmatrix} 5 & 1 & -1 \\ 0 & 4 & 1 \\ 1 & -1
& 6 \end{psmallmatrix}$. The row radii are $R_{1} = \abs{1} + \abs{-1}
= 2$, $R_{2} = \abs{0} + \abs{1} = 1$, $R_{3} = \abs{1} + \abs{-1} =
2$. The three discs are centred at $5$, $4$, $6$ with radii $2$, $1$,
$2$, so they span the real intervals $[3, 7]$, $[3, 5]$, $[4, 8]$. None
reaches the origin, so $\mat{A}$ is invertible and its eigenvalues all
lie between $3$ and $8$; without computing a single eigenvalue the
engineer knows the matrix is well clear of singularity. The dominance
check confirms it directly: $\abs{5} > 2$, $\abs{4} > 1$, $\abs{6} >
2$, so $\mat{A}$ is strictly diagonally dominant, Gaussian elimination
needs no pivoting, and the smallest possible eigenvalue magnitude is at
least $\min_{i}(\abs{a_{ii}} - R_{i}) = 3$. That lower bound feeds a
crude condition-number ceiling, since $\norm{\mat{A}}_{\infty} = 8$
gives $\kappa_{\infty} \lesssim 8/3 \approx 2.7$ once the inverse's
norm is bounded through the dominance gap. The discs converted a scan
of nine entries into a guarantee that no factorisation would surprise
the engineer.
\end{example}

\subsection{Scaling and equilibration}

Not every large condition number reflects a hard problem. A matrix
whose columns carry different physical units, a displacement in metres
beside a rotation in radians, or a stiffness beside a compliance, can
show a condition number of $10^{8}$ that collapses to single digits the
moment the columns are rescaled to a common magnitude. The
ill-conditioning is an artefact of the units, not a property of the map,
and \engterm{equilibration} removes it. The operation divides each
column of $\mat{A}$ by its norm, $\mat{A} = \tilde{\mat{A}}\mat{D}$ with
$\mat{D}$ diagonal, solves the better-conditioned $\tilde{\mat{A}}$
system, and rescales the solution by $\mat{D}$ at the end. The solution
of the original problem is recovered exactly; only the conditioning
changed. Row scaling does the same for equations that carry mismatched
units on the right-hand side.
Figure~\ref{fig:vol02:ch09:scaling-conditioning} shows the condition
number sweeping through a basin as the column scales are brought into
balance.

\begin{figure}[ht]
\centering
\begin{tikzpicture}[x=1.0cm, y=1.0cm, font=\small]

% Axes
\draw[->, thick] (0, 0) -- (10.6, 0)
  node[anchor=west, font=\scriptsize] {column-scale ratio $\log_{10}(s_{2}/s_{1})$};
\draw[->, thick] (0, 0) -- (0, 6.4)
  node[anchor=south, font=\scriptsize] {$\log_{10}\kappa$};

% x ticks
\foreach \x/\lab in {1/-4, 3/-2, 5/0, 7/2, 9/4} {
  \draw (\x, 0.08) -- (\x, -0.08) node[anchor=north, font=\scriptsize] {\lab};
}
% y ticks
\foreach \y in {0,2,4,6} {
  \draw (0.08, \y) -- (-0.08, \y) node[anchor=east, font=\scriptsize] {\y};
}

% U-shaped conditioning curve: minimum at balanced scaling (x=5), rising both ways
\draw[very thick, engblue]
  (1, 5.6) .. controls (3, 3.0) and (4.2, 1.1) ..
  (5, 0.7) .. controls (5.8, 1.1) and (7, 3.0) .. (9, 5.6);

% minimum marker
\fill[engred] (5, 0.7) circle (2pt);
\node[engred, anchor=south, font=\scriptsize] at (5, 0.85) {equilibrated: $\kappa$ minimal};

% unscaled point off to one side
\fill[engamber] (8, 4.4) circle (2pt);
\node[engamber, anchor=south west, font=\scriptsize] at (7.2, 4.5) {as-supplied units};
\draw[->, engamber, >=stealth, dashed] (7.9, 4.2) -- (5.3, 0.95);
\node[engamber, font=\scriptsize, anchor=west] at (5.6, 2.6) {rescale};

\node[anchor=north, font=\scriptsize, text width=11cm, align=center]
  at (5.3, -0.9) {
  Condition number against the relative scale of two matrix columns.
  A matrix whose columns carry mismatched physical units (a stiffness
  in $\si{\newton\per\meter}$ beside one in
  $\si{\newton\meter\per\radian}$) sits far up one arm; dividing each
  column by its norm moves it toward the basin, recovering digits the
  units had thrown away.
};

\end{tikzpicture}
\caption[Column scaling and the condition number]{Condition number of a
two-column matrix as a function of the ratio between the columns'
scales. Mismatched column magnitudes, common when the columns carry
different physical units, push the matrix up either arm of the curve and
inflate $\kappa$ for a reason that has nothing to do with the problem's
true sensitivity. Column equilibration, dividing each column by its norm
before the solve, moves the matrix toward the basin where $\kappa$ is
smallest, and the corresponding row or column scaling of the right-hand
side leaves the solution unchanged. Scaling is the cheapest defence
against artificial ill-conditioning, and a solver that reports $\kappa$
should report it after equilibration, not before.}
\label{fig:vol02:ch09:scaling-conditioning}
\end{figure}


\begin{example}[Equilibration rescues a mixed-unit system]
A two-variable calibration relates a temperature reading and a pressure
reading through
\[
\mat{A} = \begin{pmatrix} 1 & 10^{4} \\ 1 & -10^{4} \end{pmatrix},
\qquad \vect{b} = \begin{pmatrix} 10^{4} \\ -10^{4} + 2 \end{pmatrix},
\]
where the second column carries a variable scaled in different units.
The singular values are near $\sqrt{2}\cdot 10^{4}$ and $\sqrt{2}$, so
$\kappa(\mat{A}) \approx 10^{4}$, and a naive solve loses four digits.
Equilibrate by dividing the second column by $10^{4}$, giving
$\tilde{\mat{A}} = \begin{psmallmatrix} 1 & 1 \\ 1 & -1
\end{psmallmatrix}$ with $\kappa(\tilde{\mat{A}}) = 1$, a perfectly
conditioned system. Solve $\tilde{\mat{A}}\tilde{\vect{x}} = \vect{b}$
to get $\tilde{x}_{1} = 1$, $\tilde{x}_{2} = 10^{4} - 1$, then rescale
the second unknown by dividing by $10^{4}$ to recover $x_{2} = 1 -
10^{-4}$ in the original units. The four digits the raw conditioning
threatened are recovered in full, because the difficulty was never in
the problem, only in the coordinates it was posed in. A production
solver equilibrates by default and reports the condition number of the
equilibrated matrix, since the pre-equilibration number would flag a
difficulty that a change of units dissolves.
\end{example}

\subsection{Regularisation: solving what conditioning would refuse}

The refusal discipline of section 9.8 says when to decline an
ill-conditioned solve. Regularisation is what the engineer does instead
when a solve is genuinely required and the raw problem is too
ill-conditioned to trust. \engterm{Tikhonov regularisation} (called
ridge regression in statistics) adds a penalty on the solution norm to
the least-squares objective,
\[
\min_{\vect{x}} \; \norm{\mat{A}\vect{x} - \vect{b}}_{2}^{2} +
\lambda\norm{\vect{x}}_{2}^{2},
\]
whose minimiser solves the modified normal equations $(\mat{A}^{
\transpose}\mat{A} + \lambda\mat{I})\vect{x} = \mat{A}^{\transpose}
\vect{b}$. The added $\lambda\mat{I}$ lifts every eigenvalue of the Gram
matrix by $\lambda$, so the smallest eigenvalue can no longer approach
zero and the effective condition number is capped at roughly
$(\sigma_{1}^{2} + \lambda)/(\sigma_{n}^{2} + \lambda)$, which for a
tiny $\sigma_{n}$ is near $\sigma_{1}^{2}/\lambda$ rather than
$\sigma_{1}^{2}/\sigma_{n}^{2}$. The parameter $\lambda$ is the
engineer's dial between fidelity to the data (small $\lambda$, low bias,
high variance) and stability of the solution (large $\lambda$, high
bias, low variance). Choosing it is the L-curve problem of
Figure~\ref{fig:vol02:ch09:tikhonov-lcurve}: sweep $\lambda$, plot the
solution norm against the residual norm, and take the corner.

\begin{figure}[ht]
\centering
\begin{tikzpicture}[x=1.0cm, y=1.0cm, font=\small]

% Axes
\draw[->, thick] (0, 0) -- (10.4, 0)
  node[anchor=west, font=\scriptsize] {residual norm $\norm{\mat{A}\vect{x}_{\lambda} - \vect{b}}_{2}$};
\draw[->, thick] (0, 0) -- (0, 6.4)
  node[anchor=south, font=\scriptsize] {solution norm $\norm{\vect{x}_{\lambda}}_{2}$};

% The L-curve: large lambda -> small solution norm, large residual (bottom right);
% small lambda -> large solution norm (noise amplified), small residual (top left)
\draw[very thick, engblue]
  (0.6, 6.0) .. controls (0.9, 3.2) and (1.6, 1.6) ..
  (2.4, 1.15) .. controls (4.5, 0.75) and (7.5, 0.55) .. (9.6, 0.45);

% corner marker (optimal lambda)
\fill[engred] (2.4, 1.15) circle (2pt);
\node[engred, anchor=south west, font=\scriptsize] at (2.5, 1.25) {corner: balanced $\lambda$};

% Over-regularised end
\node[engamber, anchor=west, font=\scriptsize] at (7.0, 0.95)
  {$\lambda$ large: over-smoothed};
\draw[->, engamber, >=stealth] (8.4, 0.85) -- (9.0, 0.55);

% Under-regularised end
\node[enggray, anchor=west, font=\scriptsize] at (1.0, 5.3)
  {$\lambda \to 0$: noise amplified};
\draw[->, enggray, >=stealth] (1.0, 5.6) -- (0.7, 5.9);

\node[anchor=north, font=\scriptsize, text width=11cm, align=center]
  at (5.2, -0.9) {
  Tracing the ridge solution $\vect{x}_{\lambda} =
  (\mat{A}^{\transpose}\mat{A} + \lambda\mat{I})^{-1}
  \mat{A}^{\transpose}\vect{b}$ as $\lambda$ sweeps from large to small
  draws an L. The vertical arm is the ill-conditioned regime where a
  small drop in residual buys a large, noise-driven growth in the
  solution; the horizontal arm is over-smoothing. The corner is the
  defensible choice.
};

\end{tikzpicture}
\caption[The L-curve of Tikhonov regularisation]{The L-curve for a
Tikhonov (ridge) regularised least-squares problem, plotting the
solution norm against the residual norm as the regularisation
parameter $\lambda$ varies. Small $\lambda$ leaves the problem
ill-conditioned and lets measurement noise inflate the solution norm
(the steep vertical arm); large $\lambda$ over-smooths and inflates the
residual (the flat horizontal arm). The corner of the L marks the
parameter that trades the two against each other, and locating it is a
standard, defensible way to choose $\lambda$ without a separate
noise-variance estimate. The regularisation caps the effective
condition number at roughly $\sigma_{1}^{2}/\lambda$, converting a
refusal into a controlled solve.}
\label{fig:vol02:ch09:tikhonov-lcurve}
\end{figure}


\begin{example}[Ridge regularisation of a near-singular fit]
The near-singular matrix $\mat{A} = \begin{psmallmatrix} 1 & 1 \\ 1 &
1.0001 \end{psmallmatrix}$, with $\kappa \approx 4 \times 10^{4}$, is to
be used in a least-squares solve for $\vect{b} = (2, 2.0001)^{\transpose}$
where the data carry noise at the $10^{-3}$ level. The unregularised
solution $\vect{x} = (1, 1)^{\transpose}$ is exact for clean data but
swings wildly under the noise the conditioning amplifies. Add
$\lambda = 10^{-3}$. The Gram matrix $\mat{A}^{\transpose}\mat{A} =
\begin{psmallmatrix} 2 & 2.0001 \\ 2.0001 & 2.0002
\end{psmallmatrix}$ becomes $\begin{psmallmatrix} 2.001 & 2.0001 \\
2.0001 & 2.0012 \end{psmallmatrix}$, whose smallest eigenvalue has been
lifted from about $2.5 \times 10^{-9}$ to about $10^{-3}$, dropping the
effective condition number from $\approx 1.6 \times 10^{9}$ to $\approx
4 \times 10^{3}$. The regularised solution is pulled slightly toward
the origin, trading a small, controlled bias for a solution that no
longer amplifies the measurement noise by four orders of magnitude. The
choice $\lambda \approx (\text{noise level})^{2}$ is the working default,
and stating $\lambda$ with the fit is the regularisation counterpart of
stating the SVD cutoff: it declares, on the record, which directions the
data were allowed to determine.
\end{example}

\subsection{Applied vignette: deblurring an image as an inverse problem}

A camera with an out-of-focus lens or a moving subject records a blurred
image. Recovering the sharp image is a linear inverse problem, and it is
the cleanest engineering demonstration of why an invertible matrix can
still be useless to invert. The blur is a linear operation: each output
pixel is a weighted average of nearby input pixels, so if the sharp
image is the vector $\vect{x}$ (its pixels stacked into a column) and the
blurred image is $\vect{b}$, then $\vect{b} = \mat{A}\vect{x} +
\vect{n}$, where $\mat{A}$ is the blur matrix and $\vect{n}$ is sensor
noise. The blur matrix is square and, for most physical blurs,
technically invertible. The naive recovery $\vect{x} = \mat{A}^{-1}
\vect{b}$ is a disaster.

\paragraph{Why the direct inverse fails.}
A blur is a smoothing operation, and smoothing suppresses high spatial
frequencies. In the singular-value picture, the small singular values of
$\mat{A}$ correspond to the fine-detail directions the blur nearly
destroys. Inverting the matrix divides by those singular values, so
$\mat{A}^{-1}$ multiplies the fine-detail directions by huge factors.
The sensor noise $\vect{n}$, spread across all directions including the
fine-detail ones, is amplified by $1/\sigma_{n}$, and for a blur of any
severity $\sigma_{n}$ is small enough that the amplified noise buries the
recovered image in speckle. The condition number is the exact statement
of the trouble: $\kappa(\mat{A}) = \sigma_{1}/\sigma_{n}$ is large
precisely because the blur has a wide dynamic range between the
frequencies it preserves and the frequencies it suppresses, and the
sensitivity bound of this section says the relative error in $\vect{x}$
can be $\kappa(\mat{A})$ times the relative noise in $\vect{b}$. A blur
with $\kappa = 10^{6}$ turns a clean $0.1\%$ sensor noise into a $10^{5}
\times$ error, which is why the direct inverse returns noise.

\paragraph{The truncated-SVD recovery.}
The SVD names the fix. Write $\vect{x} = \sum_{i} (\vect{u}_{i}^{
\transpose}\vect{b}/\sigma_{i})\vect{v}_{i}$, the SVD form of
$\mat{A}^{-1}\vect{b}$. The terms with small $\sigma_{i}$ are the ones
that explode, and they are also the ones the blur has already stripped
of reliable signal, so the noise dominates them. Truncating the sum at
the singular values that stand above the noise floor, keeping only the
$r$ terms with $\sigma_{i}$ above a cutoff, discards the noise-dominated
directions and keeps the signal-dominated ones. The recovered image is
sharper than the blurred input and free of the speckle the direct
inverse produced. The cutoff is the engineering choice: too high and the
image stays blurred (over-truncation), too low and the speckle returns
(under-truncation). The truncation is the low-rank approximation of the
earlier example applied to the operator rather than the image, and the
cutoff plays the role the SVD cutoff plays in the chapter's project.

\paragraph{The regularised recovery.}
Tikhonov regularisation reaches the same end by a smooth dial instead of
a hard cutoff. The ridge solution $\vect{x}_{\lambda} = (\mat{A}^{
\transpose}\mat{A} + \lambda\mat{I})^{-1}\mat{A}^{\transpose}\vect{b}$
weights the $i$-th SVD term by the \engterm{filter factor}
$\sigma_{i}^{2}/(\sigma_{i}^{2} + \lambda)$, which is near one for large
singular values (kept) and near zero for small ones (suppressed). The
filter is a soft version of the truncation, and the L-curve of
Figure~\ref{fig:vol02:ch09:tikhonov-lcurve} chooses $\lambda$. In
practice the regularisation term is often not $\norm{\vect{x}}_{2}^{2}$
but a roughness penalty $\norm{\mat{L}\vect{x}}_{2}^{2}$ with $\mat{L}$
a discrete derivative, which penalises jagged solutions and prefers
smooth ones, the prior knowledge that a real image is mostly smooth. The
same machinery, with the same L-curve, recovers a deblurred image, a
denoised measurement, or a regularised regression coefficient; the
inverse problem, the conditioning failure, and the regularised fix are
one pattern across imaging, geophysics, and tomography.

\paragraph{What the vignette teaches.}
Invertibility is not usability. The blur matrix is invertible in exact
arithmetic and worthless to invert in finite precision with noisy data,
because the conditioning routes the noise straight into the answer. The
engineer's move is to declare, through a truncation cutoff or a
regularisation parameter, which directions the data can support, and to
solve the reduced problem those directions define. This is the refusal
discipline of section 9.8 in constructive form: rather than refuse the
solve outright, refuse the noise-dominated directions and solve the
rest.

\subsection{Iterative solvers and sparse systems}

The cubic cost of direct factorisation sets a ceiling on the dense
systems an engineer can solve, around $n \approx 10^{4}$ on a single
core within a working session. Real engineering systems run far
larger, a finite-element mesh of a bridge or an aircraft has $10^{6}$
to $10^{8}$ unknowns, and they are saved by structure. A
\engterm{sparse} matrix has almost all entries zero; a mesh stiffness
matrix has a handful of nonzeros per row, because each node couples
only to its mesh neighbours. Storing only the nonzeros, in a
compressed-row or coordinate format, turns an impossible $\oom{n^{2}}$
storage into a manageable $\oom{n}$, and a sparse-direct
factorisation with a good ordering keeps the fill-in (the new
nonzeros the elimination creates) under control. The ordering is not a
detail. The same sparse matrix, eliminated in two different orders, can
produce a factor that is nearly as sparse as the original or one that is
almost dense, and Figure~\ref{fig:vol02:ch09:fillin-reordering} shows
the two extremes on an arrow-shaped matrix. Finding a good ordering is a
combinatorial problem the sparse-direct libraries solve with
nested-dissection and minimum-degree heuristics before any arithmetic
begins.

\begin{figure}[ht]
\centering
\begin{tikzpicture}[x=1.0cm, y=1.0cm, font=\small]

% Left panel: arrow matrix, original ordering (heavy fill-in).
% Cell (row r, col c) plots at x = (c-0.5)*0.5, y = (6-r+0.5)*0.5.
\begin{scope}[shift={(0,0)}]
\draw[thick] (0,0) rectangle (3, 3);
% diagonal
\foreach \r in {1,...,6} {
  \pgfmathsetmacro{\cx}{(\r-0.5)*0.5}
  \pgfmathsetmacro{\cy}{(6-\r+0.5)*0.5}
  \fill[engblue] ({\cx},{\cy}) circle (2.4pt);
}
% first row (columns 2..6)
\foreach \c in {2,...,6} {
  \pgfmathsetmacro{\cx}{(\c-0.5)*0.5}
  \pgfmathsetmacro{\cy}{(6-1+0.5)*0.5}
  \fill[engblue] ({\cx},{\cy}) circle (2.4pt);
}
% first column (rows 2..6)
\foreach \r in {2,...,6} {
  \pgfmathsetmacro{\cx}{(1-0.5)*0.5}
  \pgfmathsetmacro{\cy}{(6-\r+0.5)*0.5}
  \fill[engblue] ({\cx},{\cy}) circle (2.4pt);
}
% fill-in in the trailing 5x5 block (rows 2..6, cols 2..6)
\foreach \r in {2,...,6} {
  \foreach \c in {2,...,6} {
    \pgfmathsetmacro{\cx}{(\c-0.5)*0.5}
    \pgfmathsetmacro{\cy}{(6-\r+0.5)*0.5}
    \node[engred, font=\tiny] at ({\cx},{\cy}) {$\times$};
  }
}
\node[anchor=north, font=\scriptsize] at (1.5, -0.2) {pivot on hub first};
\node[anchor=south, font=\scriptsize] at (1.5, 3.15) {dense fill-in ($\times$)};
\end{scope}

% Right panel: reordered, minimal fill-in.
\begin{scope}[shift={(6.2,0)}]
\draw[thick] (0,0) rectangle (3, 3);
% diagonal
\foreach \r in {1,...,6} {
  \pgfmathsetmacro{\cx}{(\r-0.5)*0.5}
  \pgfmathsetmacro{\cy}{(6-\r+0.5)*0.5}
  \fill[engblue] ({\cx},{\cy}) circle (2.4pt);
}
% last row (columns 1..5)
\foreach \c in {1,...,5} {
  \pgfmathsetmacro{\cx}{(\c-0.5)*0.5}
  \pgfmathsetmacro{\cy}{(6-6+0.5)*0.5}
  \fill[engblue] ({\cx},{\cy}) circle (2.4pt);
}
% last column (rows 1..5)
\foreach \r in {1,...,5} {
  \pgfmathsetmacro{\cx}{(6-0.5)*0.5}
  \pgfmathsetmacro{\cy}{(6-\r+0.5)*0.5}
  \fill[engblue] ({\cx},{\cy}) circle (2.4pt);
}
\node[anchor=north, font=\scriptsize] at (1.5, -0.2) {pivot on hub last};
\node[anchor=south, font=\scriptsize] at (1.5, 3.15) {no fill-in};
\end{scope}

\node[anchor=north, font=\scriptsize, text width=11.5cm, align=center]
  at (4.6, -0.9) {
  The same arrow matrix under two orderings. Eliminating the densely
  connected hub node first (left) couples every remaining unknown and
  fills the trailing block; eliminating it last (right) leaves the
  pattern untouched. Ordering is the free lever that separates a
  tractable sparse factorisation from a dense one.
};

\end{tikzpicture}
\caption[Fill-in depends on elimination order]{Nonzero patterns of an
arrow-shaped sparse matrix under two elimination orderings. Filled
circles are original nonzeros; red crosses are fill-in, the new nonzeros
the elimination creates. Eliminating the hub node first (left panel)
couples every remaining unknown to every other and fills the trailing
block completely, so the factorisation costs and stores as if the matrix
were dense. Eliminating the hub last (right panel) creates no fill-in at
all. The reordering is a symmetric permutation
$\mat{P}\mat{A}\mat{P}^{\transpose}$ that leaves the solution unchanged
while changing the factorisation cost by orders of magnitude; a sparse
direct solver's first job is to find a good ordering, and
nested-dissection and minimum-degree heuristics are the standard tools.}
\label{fig:vol02:ch09:fillin-reordering}
\end{figure}


\subsection{Applied vignette: the discrete Poisson problem on a grid}

The single most common large linear system in engineering is the
discretised Poisson equation, $-\nabla^{2}u = f$, which governs steady
heat conduction, electrostatic potential, incompressible-flow pressure,
and the elastic membrane. Discretising it on a grid produces a linear
system whose structure makes every idea in this section concrete at
once: sparsity, bandedness, the Kronecker product, positive
definiteness, conditioning that grows with resolution, and the choice
between direct and iterative solvers.

\paragraph{From the equation to the matrix.}
On a one-dimensional grid of $N$ interior points with spacing $h$, the
second derivative at point $i$ is approximated by the central difference
$u''(x_{i}) \approx (u_{i-1} - 2u_{i} + u_{i+1})/h^{2}$, so the operator
$-u''$ becomes the tridiagonal matrix $\mat{T} = h^{-2}\,
\mathrm{tridiag}(-1, 2, -1)$, the same matrix whose determinant the
$4 \times 4$ example of section 9.4 computed as $n + 1$. On a
two-dimensional $N \times N$ grid, each interior node couples to its
four neighbours, the five-point stencil weighting the node by $+4$ and
each neighbour by $-1$ (before the $h^{-2}$ scaling), drawn in the left
panel of Figure~\ref{fig:vol02:ch09:poisson-stencil}. The full operator
is the Kronecker sum $\mat{L}_{2\mathrm{D}} = \mat{T} \otimes \mat{I} +
\mat{I} \otimes \mat{T}$ of section 9.2, an $N^{2} \times N^{2}$ matrix
that is sparse (five nonzeros per row), symmetric, positive definite,
and banded with bandwidth $N$.

\begin{figure}[ht]
\centering
\begin{tikzpicture}[x=1.0cm, y=1.0cm, font=\small]

% Left: the grid with five-point stencil.
\begin{scope}[shift={(0,0)}]
% 5x5 grid of nodes
\foreach \i in {0,...,4} {
  \foreach \j in {0,...,4} {
    \fill[enggray!55] (\i*0.75, \j*0.75) circle (1.3pt);
  }
}
% highlight centre node and its four neighbours
\def\cx{1.5} \def\cy{1.5}
\draw[thick, engblue] (\cx,\cy) -- ({\cx+0.75},\cy);
\draw[thick, engblue] (\cx,\cy) -- ({\cx-0.75},\cy);
\draw[thick, engblue] (\cx,\cy) -- (\cx,{\cy+0.75});
\draw[thick, engblue] (\cx,\cy) -- (\cx,{\cy-0.75});
\fill[engred] (\cx,\cy) circle (2.6pt);
\node[engred, font=\scriptsize, anchor=north east] at (\cx,\cy) {$4$};
\foreach \dx/\dy in {0.75/0, -0.75/0, 0/0.75, 0/-0.75} {
  \fill[engblue] ({\cx+\dx},{\cy+\dy}) circle (2.3pt);
  \node[engblue, font=\scriptsize] at ({\cx+\dx*1.35},{\cy+\dy*1.35}) {$-1$};
}
\node[anchor=north, font=\scriptsize] at (1.5, -0.4) {five-point stencil};
\end{scope}

% Right: Kronecker structure of the operator.
\begin{scope}[shift={(5.6,0.2)}]
\node[font=\small] at (2.6, 2.6)
  {$\mat{L}_{2\mathrm{D}} = \mat{T} \otimes \mat{I} + \mat{I} \otimes \mat{T}$};
\node[font=\scriptsize, text width=5.4cm, anchor=north] at (2.6, 2.2) {
  with $\mat{T}$ the $1$D second-difference matrix
  $\mathrm{tridiag}(-1, 2, -1)$. The block matrix is banded and
  symmetric positive definite; its eigenvalues are the pairwise sums
  $\lambda_{jk} = \mu_{j} + \mu_{k}$ of the $1$D spectrum, so its
  condition number is read off the $1$D problem at no extra cost.
};
% sketch of block-banded pattern
\begin{scope}[shift={(0.4,-1.9)}]
\draw[thick] (0,0) rectangle (4.2,1.4);
\foreach \k in {0,1,2} {
  \fill[engblue!35] ({\k*1.4},{1.4-0.47*(\k+1)}) rectangle ({\k*1.4+1.4},{1.4-0.47*\k});
}
% off-diagonal identity blocks (super- and sub-diagonal)
\foreach \k in {0,1} {
  \fill[engamber!30] ({\k*1.4+1.4},{1.4-0.47*(\k+1)}) rectangle ({\k*1.4+1.4+1.4},{1.4-0.47*\k});
  \fill[engamber!30] ({\k*1.4},{1.4-0.47*(\k+2)}) rectangle ({\k*1.4+1.4},{1.4-0.47*(\k+1)});
}
\node[font=\scriptsize, anchor=north] at (2.1,-0.15) {block-tridiagonal};
\end{scope}
\end{scope}

\node[anchor=north, font=\scriptsize, text width=11.5cm, align=center]
  at (4.5, -2.5) {
  The discrete Laplacian on a square grid couples each interior node to
  its four neighbours (left). Written with the Kronecker product
  (right), the two-dimensional operator is a Kronecker sum of the
  one-dimensional second-difference matrix, which makes its bandedness,
  its symmetry, and its whole spectrum inheritable from the cheap
  one-dimensional problem.
};

\end{tikzpicture}
\caption[The discrete Poisson operator and its Kronecker structure]{The
five-point finite-difference Laplacian on a structured grid and its
Kronecker-sum form. Each interior grid node contributes one row that
weights the node by $+4$ and its four neighbours by $-1$ (left),
producing a large sparse matrix that is block-tridiagonal in the natural
node ordering. The Kronecker identity $\mat{L}_{2\mathrm{D}} = \mat{T}
\otimes \mat{I} + \mat{I} \otimes \mat{T}$ (right) exposes the structure
the solver exploits: the operator is symmetric positive definite, its
bandwidth is the grid width, and its eigenvalues are sums of the
one-dimensional eigenvalues, so the conditioning of an $N^{2} \times
N^{2}$ grid problem follows from an $N \times N$ analysis. This is the
algebraic reason fast Poisson solvers and multigrid methods exist.}
\label{fig:vol02:ch09:poisson-stencil}
\end{figure}


\paragraph{The conditioning grows with resolution.}
The eigenvalues of the one-dimensional operator are $\mu_{j} =
(4/h^{2})\sin^{2}(j\pi h/2)$ for $j = 1, \ldots, N$, so the largest is
near $4/h^{2}$ and the smallest near $\pi^{2}$ (the continuous
eigenvalue), giving a condition number $\kappa(\mat{T}) \approx
4/(\pi^{2}h^{2}) = \oom{N^{2}}$. The Kronecker sum inherits this: the
two-dimensional eigenvalues are the pairwise sums $\lambda_{jk} = \mu_{j}
+ \mu_{k}$, so $\kappa(\mat{L}_{2\mathrm{D}}) = \oom{N^{2}}$ as well.
Refining the grid to resolve finer detail, doubling $N$, quadruples the
condition number. This is the structural reason a Poisson solver on a
fine grid loses digits: the physics of the continuous problem, where
smooth low-frequency modes decay slowly and sharp high-frequency modes
decay fast, shows up in the matrix as a wide spread of eigenvalues, and
that spread is the condition number. The engineer reads the resolution
off the conditioning budget: at $N = 10^{3}$, $\kappa \approx 10^{6}$,
which double precision handles with ten digits to spare; at $N = 10^{6}$
in one dimension, $\kappa \approx 10^{12}$, which begins to bite.

\paragraph{Direct against iterative.}
For a two-dimensional grid the banded direct solve costs $\oom{n b^{2}}
= \oom{N^{2} \cdot N^{2}} = \oom{N^{4}}$ flops with the natural
ordering, because the bandwidth is $N$ and the matrix size is $N^{2}$; a
nested-dissection ordering brings this down to $\oom{N^{3}}$ and the
fill-in to $\oom{N^{2}\log N}$, the best a sparse-direct solver
achieves on this problem. Iterative methods do better asymptotically.
Conjugate gradient on the symmetric positive-definite Poisson matrix
converges in $\oom{\sqrt{\kappa}} = \oom{N}$ iterations, each an
$\oom{N^{2}}$ matrix-vector product, for an $\oom{N^{3}}$ total that a
good preconditioner cuts further. Multigrid, built specifically for this
elliptic operator, drives the residual down in $\oom{N^{2}}$ work, one
pass over the unknowns, by treating the wide eigenvalue spread head-on:
smooth the high-frequency error on the fine grid, transfer the
low-frequency error to a coarse grid where it is cheap to resolve, and
recurse. The wide conditioning that makes the problem hard for a
single-grid iteration is exactly what multigrid exploits, and the
Poisson problem is where the method was born.

\paragraph{What the vignette teaches.}
One partial differential equation exercises the whole chapter. The
discretisation is a linear system; its matrix is sparse, banded,
symmetric positive definite, and expressible as a Kronecker sum; its
conditioning grows as the square of the resolution and is read off the
one-dimensional spectrum for free; Cholesky is the factorisation the
structure invites; the fill-in is controlled by ordering; and above a
size threshold the iterative and multigrid methods take over because the
cubic wall of direct factorisation cannot be scaled. Volume IV builds
the numerical solution of partial differential equations on exactly this
foundation, and the algebra is the algebra of this chapter.

\begin{mastery}
Iterative solvers reach the systems direct factorisation cannot.
Rather than factor $\mat{A}$, they generate a sequence
$\vect{x}^{(0)}, \vect{x}^{(1)}, \ldots$ converging to the solution,
each step costing one or a few matrix-vector products, which for a
sparse matrix is $\oom{n}$ rather than $\oom{n^{3}}$. The
\engterm{stationary} methods split $\mat{A} = \mat{M} - \mat{N}$ and
iterate $\mat{M}\vect{x}^{(k+1)} = \mat{N}\vect{x}^{(k)} + \vect{b}$
with $\mat{M}$ cheap to invert: Jacobi takes $\mat{M}$ the diagonal,
Gauss--Seidel takes $\mat{M}$ the lower triangle. Both converge when
the iteration matrix $\mat{M}^{-1}\mat{N}$ has spectral radius below
one, guaranteed for diagonally dominant or symmetric positive-definite
$\mat{A}$, and they converge linearly, dropping the residual by a
fixed factor per step (Figure~\ref{fig:vol02:ch09:iterative-convergence}).
The \engterm{conjugate gradient} method, for symmetric positive-definite
$\mat{A}$, builds an optimal solution within the expanding Krylov
subspace $\spn(\vect{b}, \mat{A}\vect{b}, \mat{A}^{2}\vect{b},
\ldots)$ and converges with an error that shrinks like
$\bigl((\sqrt{\kappa} - 1)/(\sqrt{\kappa} + 1)\bigr)^{k}$. The
$\sqrt{\kappa}$, against the $\kappa$ that governs the direct-solve
error bound, is the structural reason a \engterm{preconditioner}
$\mat{M} \approx \mat{A}$ pays so steeply: replacing $\mat{A}\vect{x}
= \vect{b}$ with $\mat{M}^{-1}\mat{A}\vect{x} = \mat{M}^{-1}\vect{b}$
lowers $\kappa$, and the iteration count falls with its square root.
The engineering choice between direct and iterative is a memory and
conditioning trade: direct factorisation is reusable and
backward stable but its fill-in can exhaust memory; preconditioned
conjugate gradient and
its non-symmetric cousins (GMRES, BiCGSTAB) keep memory at $\oom{n}$
but need a preconditioner whose quality the engineer must engineer
\cite[ch.~32]{text:trefethen-bau1997}.
\end{mastery}

\begin{figure}[ht]
\centering
\begin{tikzpicture}[x=1.1cm, y=0.85cm, font=\small]

% Axes
\draw[->, thick] (0, 0) -- (10.6, 0)
  node[anchor=west, font=\scriptsize] {iteration $k$};
\draw[->, thick] (0, 0) -- (0, 6.6)
  node[anchor=south, font=\scriptsize] {$\log_{10}\norm{\vect{r}_k}_2$};

% x ticks
\foreach \k in {0,2,4,6,8,10} {
  \draw (\k, 0.08) -- (\k, -0.08) node[anchor=north, font=\scriptsize] {\k};
}
% y ticks: 0 down to -6, drawn upward as 0..6 then label negated
\foreach \y/\lab in {0/0, 1/-1, 2/-2, 3/-3, 4/-4, 5/-5, 6/-6} {
  \draw (0.08, \y) -- (-0.08, \y) node[anchor=east, font=\scriptsize] {\lab};
}

% Jacobi: slow linear decay (spectral radius ~ 0.8 => ~0.1 decade/iter)
\draw[very thick, engamber]
  (0, 0) -- (1, 0.32) -- (2, 0.64) -- (3, 0.97) -- (4, 1.29)
  -- (5, 1.61) -- (6, 1.93) -- (7, 2.25) -- (8, 2.58) -- (9, 2.90) -- (10, 3.22);
\node[engamber, anchor=west] at (7.1, 2.6) {Jacobi};

% Gauss-Seidel: roughly twice the rate of Jacobi here
\draw[very thick, engblue]
  (0, 0) -- (1, 0.62) -- (2, 1.25) -- (3, 1.87) -- (4, 2.50)
  -- (5, 3.12) -- (6, 3.74) -- (7, 4.37) -- (8, 4.99) -- (9, 5.62) -- (10, 6.24);
\node[engblue, anchor=south] at (8.2, 5.3) {Gauss--Seidel};

% Conjugate gradient: superlinear, faster late
\draw[very thick, engred]
  (0, 0) -- (1, 0.55) -- (2, 1.30) -- (3, 2.25) -- (4, 3.30)
  -- (5, 4.40) -- (6, 5.55) -- (7, 6.40);
\node[engred, anchor=west] at (4.4, 4.0) {conjugate gradient};

\node[anchor=north, font=\scriptsize, text width=10.5cm, align=center]
  at (5.0, -1.0) {
  Residual norm against iteration for a symmetric positive-definite
  system, schematic. Stationary iterations (Jacobi, Gauss--Seidel)
  drop the residual by a fixed factor per step set by the iteration's
  spectral radius. Conjugate gradient accelerates as it builds the
  Krylov subspace and terminates in at most $n$ steps in exact
  arithmetic.
};

\end{tikzpicture}
\caption[Convergence of stationary and Krylov iterations]{Schematic
residual-norm history for three iterative solvers on a symmetric
positive-definite system. The vertical axis is $\log_{10}$ of the
residual two-norm, so a straight descending line is geometric
convergence at a fixed rate per iteration. Jacobi and Gauss--Seidel
each converge linearly, with Gauss--Seidel typically about twice as
fast on the same matrix; conjugate gradient converges superlinearly,
its rate governed by $\sqrt{\kappa(\mat{A})}$ rather than
$\kappa(\mat{A})$, which is the reason a good preconditioner that
lowers $\kappa$ pays off so steeply. The slopes are illustrative; the
exact rates depend on the matrix spectrum.}
\label{fig:vol02:ch09:iterative-convergence}
\end{figure}


\begin{example}[One Gauss--Seidel sweep by hand]
Take the diagonally dominant system $\begin{psmallmatrix} 4 & 1 & 0
\\ 1 & 4 & 1 \\ 0 & 1 & 4 \end{psmallmatrix}\vect{x} =
\begin{psmallmatrix} 5 \\ 6 \\ 5 \end{psmallmatrix}$, whose exact
solution is $(1, 1, 1)^{\transpose}$. Starting from $\vect{x}^{(0)}
= \vect{0}$, Gauss--Seidel updates each component in turn using the
freshest values: $x_{1} = (5 - 0)/4 = 1.25$; then $x_{2} = (6 - 1.25
- 0)/4 = 1.1875$; then $x_{3} = (5 - 1.1875)/4 = 0.9531$. The first
sweep already lands within $20\%$ of the solution. A second sweep
gives $x_{1} = (5 - 1.1875)/4 = 0.9531$, $x_{2} = (6 - 0.9531 -
0.9531)/4 = 1.0234$, $x_{3} = (5 - 1.0234)/4 = 0.9941$, tightening
toward $(1, 1, 1)$. Each sweep is three scalar updates,
$\oom{\text{nonzeros}}$ work, with no factorisation stored. For this
$3 \times 3$ the direct solve is cheaper; the iterative advantage
appears only when $n$ is large and $\mat{A}$ is sparse, where the
sweep cost stays $\oom{n}$ while a direct factorisation's fill-in
would not.
\end{example}

\begin{mastery}
The conjugate-gradient method works only for symmetric
positive-definite matrices; the non-symmetric systems that arise from
convection, from non-self-adjoint operators, and from many
discretised flow problems need a different Krylov method, and the
standard choice is \engterm{GMRES} (generalised minimal residual). It
builds the same expanding Krylov subspace $\mathcal{K}_{k} =
\spn(\vect{r}_{0}, \mat{A}\vect{r}_{0}, \ldots, \mat{A}^{k-1}
\vect{r}_{0})$ from the initial residual, but instead of minimising
energy in the $\mat{A}$-inner-product (which is not defined for a
non-symmetric matrix) it minimises the ordinary residual norm
$\norm{\vect{b} - \mat{A}\vect{x}_{k}}_{2}$ over that subspace. The
minimisation is a small least-squares problem solved by the QR
factorisation of section 9.3 applied to the $(k+1) \times k$
Hessenberg matrix the Arnoldi process produces, so the QR machinery of
this chapter is the inner engine of GMRES. The cost is the method's
liability: GMRES stores the entire Krylov basis, so memory grows
linearly with the iteration count and the orthogonalisation cost grows
quadratically. Production codes therefore \engterm{restart} GMRES
every $m$ iterations, discarding the accumulated basis and beginning
afresh from the current iterate, trading guaranteed monotone
convergence for bounded memory. The restart parameter $m$ is an
engineering choice, and a poorly chosen $m$ can stall convergence
entirely, the non-symmetric counterpart of the preconditioner-quality
problem that governs conjugate gradient
\cite[ch.~35]{text:trefethen-bau1997}.
\end{mastery}

\begin{mastery}
When a matrix is large and only approximately low-rank, the classical
$\oom{mn^{2}}$ SVD is wasteful, because the engineer wants only the
top few singular triplets. \engterm{Randomised} linear algebra computes
them at a fraction of the cost by a deliberately probabilistic
sketch. To approximate the dominant rank-$r$ subspace of $\mat{A} \in
\R^{m \times n}$, draw a random Gaussian matrix $\mat{\Omega} \in
\R^{n \times (r + p)}$ with a small oversampling $p$, form the sketch
$\mat{Y} = \mat{A}\mat{\Omega}$ (one matrix-matrix product, cheap when
$\mat{A}$ is sparse or available only as an operator), orthonormalise
$\mat{Y}$ to a thin $\mat{Q}$ by QR, and project: $\mat{B} =
\mat{Q}^{\transpose}\mat{A}$ is small, and its SVD lifts back to an
approximate SVD of $\mat{A}$. The random projection captures the
dominant subspace with overwhelming probability, and the
approximation error is within a modest factor of the optimal
Eckart--Young error $\sigma_{r+1}$, with the failure probability
falling exponentially in the oversampling $p$. The method's value is
in the regime modern data analysis lives in: a matrix too large to
factor densely, available only through matrix-vector products, and
known to be well approximated by low rank. Randomisation converts the
intractable deterministic decomposition into a tractable probabilistic
one, with a probability of failure the engineer can drive below any
threshold by spending a few extra columns of the sketch
\cite{paper:v2ch10-eckart-young1936}.
\end{mastery}

\subsection{Singular value decomposition (recognition)}

Every real $m \times n$ matrix factors as

\[
\mat{A} = \mat{U}\mat{\Sigma}\mat{V}^{\transpose},
\]

where $\mat{U} \in \R^{m \times m}$ and $\mat{V} \in \R^{n \times n}$
are orthogonal, and $\mat{\Sigma} \in \R^{m \times n}$ is
``diagonal'' (the entries off the main diagonal are zero) with
non-negative entries $\sigma_{1} \geq \cdots \geq \sigma_{r} > 0$
on the diagonal, where $r = \rank \mat{A}$. The factorisation is the
\engterm{singular value decomposition}, or SVD
\cite[ch.~4]{text:trefethen-bau1997}. Volume II Chapter 10 develops
its derivation; this chapter introduces it at recognition level
because the project relies on it.

The SVD reveals exactly what the conditioning analysis hides. The
singular values $\sigma_{i}$ are the rank-revealing diagonals: a
small $\sigma_{i}$ means a near-singular direction, and the column
$\vect{v}_{i}$ of $\mat{V}$ associated with it is the input direction
the matrix nearly annihilates. A least-squares solution computed
with the SVD can truncate the smallest singular values explicitly
(the \engterm{truncated SVD}), which gives a regularised solution
whose conditioning the engineer controls. The QR and the normal
equations do not offer this control without further work.
Figure~\ref{fig:vol02:ch09:svd-geometry} draws the geometric
content: rotate, scale, rotate.

\begin{figure}[ht]
\centering
\begin{tikzpicture}[x=1.0cm, y=1.0cm, font=\small]

% Panel 1: input unit circle in R^2
\begin{scope}[shift={(0,0)}]
  \node[anchor=south, font=\small\bfseries] at (0, 2.0) {input: unit circle};
  \draw[->, thick] (-1.8, 0) -- (1.8, 0);
  \draw[->, thick] (0, -1.8) -- (0, 1.8);
  \draw[thick, engblue] (0,0) circle (1.5);
  \draw[->, very thick, engred] (0, 0) -- (1.5, 0) node[anchor=west, font=\scriptsize] {$\vect{v}_{1}$};
  \draw[->, very thick, engred] (0, 0) -- (0, 1.5) node[anchor=south, font=\scriptsize] {$\vect{v}_{2}$};
\end{scope}

% Arrow with V^T
\draw[->, thick] (1.9, 0) -- (2.6, 0) node[midway, above, font=\scriptsize] {$\mat{V}^{\transpose}$};

% Panel 2: rotated unit circle
\begin{scope}[shift={(4.2, 0)}]
  \node[anchor=south, font=\small\bfseries] at (0, 2.0) {rotate};
  \draw[->, thick] (-1.8, 0) -- (1.8, 0);
  \draw[->, thick] (0, -1.8) -- (0, 1.8);
  \draw[thick, engblue] (0,0) circle (1.5);
  \draw[->, very thick, engred] (0, 0) -- (1.5, 0) node[anchor=west, font=\scriptsize] {$\vect{e}_{1}$};
  \draw[->, very thick, engred] (0, 0) -- (0, 1.5) node[anchor=south, font=\scriptsize] {$\vect{e}_{2}$};
\end{scope}

\draw[->, thick] (6.1, 0) -- (6.8, 0) node[midway, above, font=\scriptsize] {$\mat{\Sigma}$};

% Panel 3: scaled ellipse, axes sigma_1 and sigma_2
\begin{scope}[shift={(8.6, 0)}]
  \node[anchor=south, font=\small\bfseries] at (0, 2.0) {scale};
  \draw[->, thick] (-2.4, 0) -- (2.4, 0);
  \draw[->, thick] (0, -1.4) -- (0, 1.4);
  \draw[thick, engblue] (0,0) ellipse (2.0 and 0.7);
  \draw[->, very thick, engred] (0, 0) -- (2.0, 0) node[anchor=west, font=\scriptsize] {$\sigma_{1}\vect{e}_{1}$};
  \draw[->, very thick, engred] (0, 0) -- (0, 0.7) node[anchor=south, font=\scriptsize] {$\sigma_{2}\vect{e}_{2}$};
\end{scope}

\draw[->, thick] (11.0, 0) -- (11.7, 0) node[midway, above, font=\scriptsize] {$\mat{U}$};

% Panel 4: final ellipse in output basis u_1, u_2
\begin{scope}[shift={(13.6, 0)}]
  \node[anchor=south, font=\small\bfseries] at (0, 2.0) {output ellipse};
  \draw[->, thick] (-2.4, 0) -- (2.4, 0);
  \draw[->, thick] (0, -1.6) -- (0, 1.6);
  \draw[thick, engblue, rotate=30] (0,0) ellipse (2.0 and 0.7);
  \draw[->, very thick, engred, rotate=30] (0, 0) -- (2.0, 0) node[anchor=west, font=\scriptsize] {$\sigma_{1}\vect{u}_{1}$};
  \draw[->, very thick, engred, rotate=30] (0, 0) -- (0, 0.7) node[anchor=south, font=\scriptsize] {$\sigma_{2}\vect{u}_{2}$};
\end{scope}

\node[font=\scriptsize, text width=14cm, align=center, anchor=north] at (6.5, -2.4) {
  $\mat{A} = \mat{U}\mat{\Sigma}\mat{V}^{\transpose}$ acts on the unit circle by
  rotating into the input singular basis, scaling each axis by $\sigma_{i}$, and
  rotating into the output singular basis. The condition number is the aspect
  ratio of the output ellipse: $\kappa(\mat{A}) = \sigma_{1}/\sigma_{n}$.
};

\end{tikzpicture}
\caption[The SVD as rotate, scale, rotate]{The singular value
decomposition $\mat{A} = \mat{U}\mat{\Sigma}\mat{V}^{\transpose}$
factors a linear map into three geometric operations: rotate the
input by $\mat{V}^{\transpose}$, scale axes by the singular values
$\sigma_{1} \geq \cdots \geq \sigma_{r} > 0$, and rotate the output
by $\mat{U}$. The unit ball maps to an ellipsoid whose principal
axes are $\sigma_{i}\vect{u}_{i}$. A small $\sigma_{i}$ flattens the
ellipsoid in the $\vect{u}_{i}$ direction; the input direction
$\vect{v}_{i}$ is the one the matrix nearly annihilates, and the
condition number $\kappa(\mat{A}) = \sigma_{1}/\sigma_{n}$ is the
aspect ratio of the output ellipsoid.}
\label{fig:vol02:ch09:svd-geometry}
\end{figure}


\begin{mastery}
The SVD reveals the rank and the near-rank of a matrix at the price of
the full spectral decomposition. When only the rank is wanted, and
wanted cheaply, \engterm{column-pivoted QR} is the working
alternative. The ordinary Householder QR of section 9.3 processes the
columns left to right; the pivoted variant, at each step, swaps in the
column of largest remaining norm before it reflects. The diagonal
entries of the resulting $\mat{R}$ then come out in decreasing
magnitude, $\abs{R_{11}} \geq \abs{R_{22}} \geq \cdots$, and a sharp
drop in that sequence marks the numerical rank: the columns before the
drop are independent to working precision, and the columns after it are
nearly dependent on them. The factorisation is $\mat{A}\mat{P} =
\mat{Q}\mat{R}$ with $\mat{P}$ the permutation of the pivots, and it
costs the same $\oom{mn^{2}}$ as unpivoted QR, an order of magnitude
below the SVD. The diagonal of $\mat{R}$ estimates the singular values
well enough for rank detection, subset selection (which columns of a
design matrix to keep), and the rank-revealing step inside a
least-squares solve that must handle a rank-deficient design matrix
without the full SVD. The SVD gives the exact singular values and the
optimal low-rank approximation; column-pivoted QR gives a fast, usually
adequate estimate of the same rank information, and the engineer chooses
between them by whether the exactness is worth the order-of-magnitude
cost \cite[ch.~5]{text:v2ch09-golub-vanloan2013}.
\end{mastery}

\begin{mastery}
A small SVD computed by hand shows where the singular values come
from. Take $\mat{A} = \begin{psmallmatrix} 3 & 0 \\ 4 & 5
\end{psmallmatrix}$. The right singular vectors and the squared
singular values are the eigenpairs of $\mat{A}^{\transpose}\mat{A} =
\begin{psmallmatrix} 25 & 20 \\ 20 & 25 \end{psmallmatrix}$, whose
eigenvalues solve $(25 - \lambda)^{2} = 400$, giving $\lambda_{1} =
45$ and $\lambda_{2} = 5$. The singular values are the square roots,
$\sigma_{1} = \sqrt{45} \approx 6.708$ and $\sigma_{2} = \sqrt{5}
\approx 2.236$, so $\kappa(\mat{A}) = \sigma_{1}/\sigma_{2} = 3$. The
right singular vectors are the eigenvectors of
$\mat{A}^{\transpose}\mat{A}$, $\vect{v}_{1} = (1, 1)^{\transpose}/
\sqrt{2}$ and $\vect{v}_{2} = (1, -1)^{\transpose}/\sqrt{2}$, and the
left singular vectors follow from $\vect{u}_{i} = \mat{A}\vect{v}_{i}/
\sigma_{i}$. The geometry the SVD encodes is now explicit: $\mat{A}$
maps the unit circle to an ellipse whose semi-axes have lengths
$\sigma_{1}$ and $\sigma_{2}$ along the directions $\vect{u}_{1}$ and
$\vect{u}_{2}$, the input directions $\vect{v}_{1}$ and $\vect{v}_{2}$
being the pre-images of those axes. The condition number is the
aspect ratio of that ellipse, the same quantity section 9.7 reads off
the singular values. The hand computation is feasible at $2 \times 2$
and is the route to internalising what the library routine returns;
beyond $3 \times 3$ the eigenvalue step itself needs iteration, which
is why the SVD carries the $\approx 22 n^{3}$ price tag of the flop
table.
\end{mastery}

\begin{mastery}
The SVD splits a matrix into rotate-scale-rotate; a companion
factorisation splits it into stretch-then-rotate, and it is the natural
language for problems that ask for the nearest rotation to a given map.
The \engterm{polar decomposition} writes any square matrix as $\mat{A}
= \mat{Q}\mat{P}$, with $\mat{Q}$ orthogonal and $\mat{P}$ symmetric
positive semidefinite. It falls straight out of the SVD: with $\mat{A}
= \mat{U}\mat{\Sigma}\mat{V}^{\transpose}$, group the factors as
$\mat{A} = (\mat{U}\mat{V}^{\transpose})(\mat{V}\mat{\Sigma}
\mat{V}^{\transpose})$, so $\mat{Q} = \mat{U}\mat{V}^{\transpose}$ is
orthogonal and $\mat{P} = \mat{V}\mat{\Sigma}\mat{V}^{\transpose}$ is
symmetric positive semidefinite. The reading is physical: $\mat{P}$ is
the pure stretch, its eigenvalues the singular values and its
eigenvectors the principal stretch directions, and $\mat{Q}$ is the
rigid rotation that follows. Continuum mechanics uses exactly this
split, where the deformation gradient factors into a stretch tensor and
a rotation, the polar decomposition theorem of finite-strain theory
that Volume III invokes. The orthogonal factor also answers a
best-approximation question the SVD certifies: $\mat{Q} = \mat{U}
\mat{V}^{\transpose}$ is the orthogonal matrix closest to $\mat{A}$ in
the Frobenius norm, which is why the polar factor is the tool for
orthogonalising a drifted rotation matrix, a rotation estimate in
robotics or spacecraft attitude that round-off has nudged off the
orthogonal manifold. Re-orthogonalising by the polar factor projects
back onto the nearest true rotation, where a Gram--Schmidt sweep would
return an orthogonal matrix that is not the nearest one
\cite[ch.~9]{text:v2ch09-golub-vanloan2013}.
\end{mastery}

\begin{mastery}
The incidence matrices of the levelling and resistor networks
generalise into one of the most-used objects in applied linear
algebra, the \engterm{graph Laplacian}. For a graph with $n$ nodes and
edge weights $w_{ij}$, the Laplacian is $\mat{L} = \mat{D} - \mat{W}$,
where $\mat{W}$ is the weight matrix and $\mat{D}$ is the diagonal
matrix of weighted degrees $D_{ii} = \sum_{j} w_{ij}$. The matrix is
symmetric positive semidefinite, the constant vector lies in its null
space (the floating-ground mode the resistor example grounded away),
and its spectrum encodes the graph's connectivity. The smallest
nonzero eigenvalue, the \engterm{algebraic connectivity} or Fiedler
value, is zero exactly when the graph is disconnected, and its
eigenvector, the \engterm{Fiedler vector}, sorts the nodes into the
two loosely connected halves: thresholding the Fiedler vector at zero
partitions the graph along its weakest cut. This is \engterm{spectral
clustering}, and it turns a combinatorial partitioning problem (which
is hard) into an eigenvalue problem (which the SVD machinery solves).
The same Laplacian governs heat diffusion on the graph, random walks,
the resistance distance between nodes, and the stiffness of a spring
network; recognising a matrix as a graph Laplacian tells the engineer
that its spectrum carries the connectivity and that its near-null
directions are the slow, large-scale modes. Volume IX uses the Fiedler
vector for mesh partitioning and load balancing on parallel solvers
\cite[ch.~9]{text:v2ch09-golub-vanloan2013}.
\end{mastery}

\begin{mastery}
The factorisations of this chapter solve algebraic systems; a glance
ahead shows the same apparatus solving differential ones. A linear
system of ordinary differential equations $\dot{\vect{x}} = \mat{A}
\vect{x}$ has the closed-form solution $\vect{x}(t) = e^{\mat{A}t}
\vect{x}(0)$, where the \engterm{matrix exponential} is defined by the
same power series as the scalar exponential, $e^{\mat{A}t} = \mat{I} +
\mat{A}t + \tfrac{1}{2}\mat{A}^{2}t^{2} + \cdots$. When $\mat{A}$ is
diagonalisable, $\mat{A} = \mat{V}\mat{\Lambda}\mat{V}^{-1}$, the
exponential decouples into scalar exponentials of the eigenvalues,
$e^{\mat{A}t} = \mat{V}e^{\mat{\Lambda}t}\mat{V}^{-1}$, so the
eigenstructure of $\mat{A}$ (Chapter 10) is what governs the dynamics:
eigenvalues with positive real part grow, negative decay, imaginary
oscillate. Computing the matrix exponential reliably is its own
numerical problem, and the naive power series is one of the recognised
ways to get it wrong, since cancellation in the series ruins the
accuracy for matrices with widely separated eigenvalues. The standard
production methods scale the matrix down, exponentiate by a Pad\'e
approximation, and square the result back up
\cite{text:v2ch10-higham-functions2008}. The preview matters here
because it closes the loop the book opened in Chapter 7: the
state-space form of a linear dynamical system is a matrix, its
stability is a spectral property of that matrix, and its time
evolution is a matrix function. Volume IX builds control and
estimation on exactly this footing.
\end{mastery}

\subsection{Floating-point overflow as a related failure}

Conditioning is one numerical failure mode; \engterm{overflow} is
another. A floating-point representation has a finite range. The
IEEE 754 double-precision format ranges from about $2.2 \times
10^{-308}$ to $1.8 \times 10^{308}$; the single-precision format
ranges from about $1.2 \times 10^{-38}$ to $3.4 \times 10^{38}$;
the 16-bit signed integer ranges from $-32768$ to $32767$. A
computation that produces a value outside the format's range
overflows. The IEEE float standard specifies that overflow yields
$\pm\infty$, and subsequent arithmetic propagates the infinity
deterministically; integer overflow has no such standard, and
languages and platforms differ in whether they wrap, saturate,
trap, or behave as undefined.

A matrix entry, an intermediate product, or a derived quantity that
overflows produces wrong results indistinguishable in the program's
output from correct ones. The discipline is to bound the magnitudes
that flow through the computation and to choose a representation
whose range covers the bound. Section 9.8 returns to this point
through a named case.

\section{Failure: ill-conditioned systems and what to refuse to solve}\pathtag{core}

An ill-conditioned linear system is not a system that fails to have
a solution. The unique solution exists, in exact arithmetic, when
the matrix is invertible. Ill-conditioning is a failure mode of the
\emph{computed} solution: the algebra promises a number, the
floating-point arithmetic delivers something else, and the engineer
who treats the delivered number as the algebraic answer is in
trouble.

\subsection{What to refuse}

A working list, in increasing severity:

\begin{itemize}[leftmargin=*]
  \item Refuse the normal-equations formulation when the design
    matrix has condition number above about $10^{7}$ in double
    precision. Use QR. Above about $10^{15}$, refuse QR as well
    and use the truncated SVD with an explicit cutoff stated in
    the report.
  \item Refuse a polynomial fit of degree higher than about five
    or six on data spread over an interval that does not include
    the origin, in the monomial basis. The Vandermonde
    conditioning blows up. Reformulate in the orthogonal-polynomial
    basis (Legendre, Chebyshev) or in the SVD-truncated form.
  \item Refuse to invert a matrix whose smallest singular value is
    below the rounding-error floor of the entries. The matrix is,
    for the precision in use, indistinguishable from a singular
    matrix.
  \item Refuse a numerical conversion between two representations
    whose ranges have not been verified against the operating
    envelope of the data the conversion handles. The Ariane 5
    case below makes this concrete.
  \item Refuse the report when the conditioning is not stated.
    A least-squares fit reported without its design-matrix
    condition number is incomplete.
\end{itemize}

The discipline is to make the conditioning visible, then to choose
the formulation whose conditioning the engineer can defend.

\subsection{The Ariane 5 case}

\paragraph{Technical mechanism.}
Ariane 5 Flight 501, the maiden flight of the European Space Agency's
new launcher, broke up about 39 seconds after main engine ignition
on 4 June 1996. The Lions Inquiry Board, reporting on 19 July 1996,
identified the cause as an unhandled software exception in the
inertial reference system (SRI), triggered during conversion of a
64-bit floating-point representation of an internal horizontal-bias
variable (BH) to a 16-bit signed integer; the result exceeded the
integer's range $[-32768, 32767]$ and the conversion overflowed
\cite{acc:ariane501-ifr}. The conversion occurred inside an
alignment function reused unchanged from Ariane 4. For Ariane 4,
the alignment function had been kept active for a few seconds after
lift-off to support a pre-launch-recovery requirement, and the
Ariane 4 trajectory kept BH within range. The Ariane 5 trajectory
developed horizontal velocity faster, and BH did not stay in range.
The variable was one of seven internal SRI variables for which
overflow protection had been considered; four were protected, three
(including BH) were left unprotected on the basis that the Ariane 4
trajectory analysis showed they could not go out of range. That
analysis was not redone for Ariane 5. The exception was unhandled;
the primary SRI shut down; the back-up SRI, running the same
software, shut down for the same reason for the same diagnostic
output. The on-board computer interpreted SRI diagnostic words as
flight data, computed extreme attitude corrections, commanded large
nozzle deflections, and the launcher broke up under aerodynamic
loads. The self-destruct system then activated. The underlying
numerical operation, real-to-integer conversion, is the simplest
possible numerical procedure: no matrix factorisation, no
conditioning analysis, no near-singularity. The failure is the
overflow failure mode of section 9.7, in its most elementary form.

\paragraph{Organisational context.}
The Inquiry Board's framing of the cause was a verification-domain
failure rather than a coding error. The implementation was correct
for the Ariane 4 envelope it had been verified against. What was
absent was a re-verification of the same implementation against the
Ariane 5 envelope when the software was reused. The board observed
that the Ariane 5 qualification programme had treated the reused
SRI software as already qualified by its Ariane 4 service history
rather than as new software whose flight envelope must be
re-verified \cite{acc:ariane501-ifr}. The Ariane 4 documentation
recorded that the variable BH could not exceed the 16-bit range on
the Ariane 4 trajectory; the documentation did not state, because
it had no reason to state, that this guarantee was a property of
the Ariane 4 trajectory rather than of the conversion itself. The
Ariane 5 verification then read the Ariane 4 guarantee as a
property of the SRI code and proceeded. The board's recommendations
were structural: re-analysis of internal variable ranges against
the new flight envelope, defensive programming for arithmetic
operations on flight-critical values, exception handling on all
arithmetic exceptions in flight-critical software, software
architecture review for reused components in qualification, and
loss of the assumption that service-history qualification on one
launcher transfers to another. The recommendations are
restatements, in the software-engineering register, of the
verification discipline of section 9.7.

\paragraph{Lesson for linear-systems modelling.}
The lesson generalises beyond software-conversion failures.
Whenever a piece of numerical code is ported across operating
envelopes (a control law, a Kalman filter, a finite-element solver,
a regression pipeline), the engineer verifies that the new
envelope's data fall inside the preconditions of every numerical
operation the code performs. The matrix-conditioning discipline of
section 9.7 is one instance of this verification: the precondition
is a bound on $\kappa(\mat{A})$ inside which a chosen factorisation
delivers a defensible answer; the operation is the solve; the
verification is the conditioning estimate. The Ariane 5 conversion
is another instance: the precondition is the integer range; the
operation is the conversion; the verification is the
range-analysis check against the operating envelope. The
discipline that prevents both failure modes is the same: name the
precondition, verify the input meets it, and refuse the operation
when verification fails. The chapter's principle (below) states
this in operative form.

\paragraph{Short-form companion cases.}
Two further failures sit beside the Ariane 5 case as standing
references for the same precondition-verification discipline,
without their own chapter section here. The Patriot missile
software failure of 25 February 1991 at Dhahran lost an Iraqi Scud
missile track because a 24-bit fixed-point clock drift accumulated
over $\approx 100\,\si{\hour}$ of continuous operation,
introducing an integration error of about $0.34\,\si{\second}$ in
the prediction of the Scud's location; the precondition violated
was a continuous-operation bound on the clock variable that
exceeded the qualification envelope. The Hubble Space Telescope
primary mirror failure of 1990 was a metrology-procedure
verification failure: the null corrector used to figure the mirror
was itself out of alignment by $1.3\,\si{\milli\meter}$ on a single
component, and the alignment was not cross-checked against an
independent measurement before the mirror was certified as the
correct figure \cite{acc:nasa-hubble-optical-systems-1990}. In
each case, an algebraic operation was performed inside a software
or measurement pipeline whose precondition the engineering process
had not verified against the new envelope. The discipline of
section 9.7 is the same in all three.

\subsection{The principle}

\begin{principle}[Refusal of unsafe numerical operations]
A numerical operation has a precondition: a range, a rank, a
conditioning bound, or a representation envelope inside which the
operation's output is defensible. The engineer states the
precondition, verifies the input meets it, and refuses the
operation when verification fails. A reported result without a
record of the verification is, for engineering purposes, not a
result. The discipline is independent of the algorithm and of the
problem's domain; the numerical conversion that destroyed Ariane 5
Flight 501 and the polynomial fit whose Gram matrix has a condition
number of $10^{14}$ are the same failure mode at different scales.
\end{principle}

\section{Project}\pathtag{core}

\begin{project}[Polynomial fit by three methods, with conditioning comparison]
The reader fits a polynomial of degree $p \geq 5$ to a noisy dataset
of their choice (the reader's own data from a Volume I project, or
a public dataset that is genuinely noisy and not synthetic), by
three methods, and produces a comparative analysis with a
1500-word reflection.

\textbf{Track}. Simulation with analysis. \textbf{Hazard class}: none.
\textbf{Tools}. Paper, pencil, and a programming environment with
linear-algebra primitives (LU, QR, SVD): Python with NumPy, MATLAB,
Julia, Octave, R, or equivalent.

\textbf{Dataset choice}. The reader chooses a dataset of $m \geq 30$
points $(x_{i}, y_{i})$, with $y_{i}$ noisy, that supports a smooth
underlying trend over the $x$-range. Public-data candidates: a
weather-station daily temperature record over a year, a household
electricity-meter reading over a month, a stock index closing price
over a quarter, a Volume I commute-time dataset, or any dataset the
reader has collected with a clear independent variable. The reader
records the dataset, its source, and the date range.

\textbf{Three formulations}. The reader computes the polynomial fit
of degree $p$ in the monomial basis $\{1, x, x^{2}, \ldots, x^{p}\}$,
producing the Vandermonde design matrix $\mat{V}$ of size $m \times
(p+1)$, by three methods:

\begin{enumerate}[leftmargin=*]
  \item \engterm{Normal equations}. Form $\mat{V}^{\transpose}\mat{V}$
    explicitly, form $\mat{V}^{\transpose}\vect{y}$, solve
    $\mat{V}^{\transpose}\mat{V}\hat{\vect{c}} = \mat{V}^{\transpose}\vect{y}$
    by LU. Record the condition number of $\mat{V}^{\transpose}\mat{V}$.
  \item \engterm{QR}. Compute $\mat{V} = \mat{Q}\mat{R}$ by the
    Householder algorithm (the standard library routine in any of
    the listed environments). Solve $\mat{R}\hat{\vect{c}} =
    \mat{Q}^{\transpose}\vect{y}$ by back substitution. Record the
    condition number of $\mat{V}$.
  \item \engterm{SVD}. Compute $\mat{V} = \mat{U}\mat{\Sigma}
    \mat{V}_{r}^{\transpose}$. Solve in the SVD form,
    $\hat{\vect{c}} = \mat{V}_{r}\mat{\Sigma}^{+}\mat{U}^{\transpose}
    \vect{y}$, where $\mat{\Sigma}^{+}$ is the pseudoinverse with
    a stated cutoff (any singular value below $\epsilon \cdot
    \sigma_{1}$ for $\epsilon = 10^{-12}$ is set to zero). Record
    the singular values and the cutoff applied.
\end{enumerate}

\textbf{Reporting}. The reader produces:

\begin{itemize}[leftmargin=*]
  \item A table of the three coefficient vectors $\hat{\vect{c}}$,
    one column per method.
  \item The three condition numbers, side by side.
  \item A plot of the data with the three fits overlaid (in any
    plotting environment).
  \item The residual norm $\norm{\mat{V}\hat{\vect{c}} -
    \vect{y}}_{2}$ for each method.
  \item A statement of the regime in which each method gives the
    most defensible answer.
\end{itemize}

\textbf{Reflection (1500 words)}. The reader writes a 1500-word
narrative addressing:

\begin{itemize}[leftmargin=*]
  \item Why the normal equations formulation is dangerous when the
    design matrix is ill-conditioned. The reader should refer to
    the squaring of the condition number and to a specific
    decimal-digit estimate from section 9.7's rule of thumb.
  \item What the SVD revealed that the normal equations
    concealed. The reader names at least one singular value of the
    Vandermonde matrix that fell near or below the cutoff and
    describes its associated input direction.
  \item How the polynomial fit would change if the reader rescaled
    $x$ to the interval $[-1, 1]$ before forming the design
    matrix, and why the rescaling is a defence against
    ill-conditioning rather than a mathematical change of the
    problem.
  \item Whether a polynomial of degree $p$ was the right model
    for the dataset, or whether the conditioning trouble pointed
    at a different model (a piecewise fit, a spline, a Fourier
    expansion). The reflection takes the conditioning seriously
    as a model-selection signal, not a numerical nuisance.
  \item What the reader would refuse to report from the three
    fits, in line with the refusal discipline of section 9.8.
\end{itemize}

\textbf{Deliverable}. The dataset, the code, the three fitted
coefficient vectors, the condition numbers, the residual norms,
the plot, and the 1500-word reflection. The general editor's
reference solution on a published-temperature dataset is available
online; the reader compares only after completing their own.

\textbf{Time}. Approximately one hour to choose and clean the
dataset, two to three hours to implement the three methods, one
hour to produce the table and plot, and three to four hours to
write the reflection. Total approximately seven to nine hours.

\textbf{Phases}.
\begin{enumerate}[leftmargin=*]
  \item \textit{Dataset selection and sanity check} (one hour).
    Choose the dataset; plot $(x_{i}, y_{i})$; compute $\bar{x}$,
    $\bar{y}$, the range of $x$, and the empirical noise estimate
    (the standard deviation of residuals from a low-degree fit).
    Record source and date.
  \item \textit{Reference implementation} (two to three hours).
    Set the polynomial degree $p$. Implement the three methods in
    a single script that consumes the dataset and produces a
    coefficient table, three condition numbers, and three residual
    norms. The reference script
    \path{code/least_squares_three_ways.py} is a working template;
    the deliverable is the reader's own implementation, not a
    transcription.
  \item \textit{Conditioning analysis} (one hour). Sweep the degree
    $p$ from $2$ to $10$ on the chosen dataset. At each $p$, record
    $\kappa(\mat{V})$, $\kappa(\mat{V}^{\transpose}\mat{V})$, and
    the residual norms for the three methods. Identify the degree
    at which the normal-equations coefficients diverge from QR.
  \item \textit{Rescaling and basis change} (one hour). Repeat the
    analysis with $x$ rescaled to $[-1, 1]$ and with a Chebyshev
    or Legendre basis. Compare conditioning gains.
  \item \textit{Reflection} (three to four hours). Address the five
    questions listed above.
\end{enumerate}

\textbf{Success criteria}.
\begin{itemize}[leftmargin=*]
  \item The three coefficient vectors are reported with their
    relative disagreement quantified as a percentage of the QR
    coefficient.
  \item The three condition numbers are reported and the rule of
    thumb is applied numerically to predict the digit loss for each
    formulation.
  \item The reflection states the regime in which each method is
    defensible and the regime in which each method must be refused.
  \item The deliverable identifies at least one singular value of
    the Vandermonde matrix that lies below or near the SVD cutoff
    and gives a one-sentence interpretation of its associated input
    direction.
\end{itemize}

\textbf{Common failure modes}.
\begin{itemize}[leftmargin=*]
  \item Forming $\mat{V}^{\transpose}\mat{V}$ for a degree-eight or
    higher fit on raw $x$ values in an interval that does not
    include zero, then reporting the resulting coefficients without
    a conditioning check. The reflection should identify this as
    the canonical Vandermonde pathology and refer to
    Figure~\ref{fig:vol02:ch09:vandermonde-conditioning}.
  \item Reporting residual norms only, without coefficient
    disagreement. Two formulations may agree on the residual to
    several digits while disagreeing on individual coefficients by
    orders of magnitude; the residual-only check fails to flag the
    conditioning problem.
  \item Treating the SVD cutoff $\epsilon$ as an algorithmic
    parameter rather than as an engineering choice. The cutoff
    encodes the engineer's judgement about which input directions
    the data can support; the reflection states the cutoff with the
    reasoning behind it.
  \item Choosing a polynomial when the dataset asks for a different
    model class. A sinusoidal trend (the temperature dataset is
    one) fits poorly to any polynomial degree the conditioning
    allows; the reflection should acknowledge the model-selection
    signal in the conditioning trouble.
\end{itemize}
\end{project}

\section{Exercises}

The exercises test calculation, derivation, estimation, simulation,
design, diagnosis, failure analysis, and judgment.

\subsection*{Calculation}\pathtag{core}

\begin{exercise}
Compute, by Gaussian elimination, the solution of $\mat{A}\vect{x}
= \vect{b}$ for $\mat{A} = \begin{psmallmatrix} 1 & 2 & -1 \\ 2 &
1 & 1 \\ -1 & 1 & 2 \end{psmallmatrix}$ and $\vect{b} = (3, 4,
1)^{\transpose}$. Record the LU factors $\mat{L}$ and $\mat{U}$
and verify $\mat{L}\mat{U} = \mat{A}$.
\end{exercise}

\paragraph{Solution sketch.}
Subtract $2 R_{1}$ from $R_{2}$ and add $R_{1}$ to $R_{3}$:
$\mat{U}^{(1)} = \begin{psmallmatrix} 1 & 2 & -1 \\ 0 & -3 & 3 \\
0 & 3 & 1 \end{psmallmatrix}$, transformed $\vect{b}^{(1)} = (3,
-2, 4)^{\transpose}$, multipliers $(2, -1)$ in column 1 of $\mat{L}$.
Add $R_{2}$ to $R_{3}$ to clear the $(3,2)$ entry: $\mat{U} =
\begin{psmallmatrix} 1 & 2 & -1 \\ 0 & -3 & 3 \\ 0 & 0 & 4
\end{psmallmatrix}$, $\vect{b}^{(2)} = (3, -2, 2)^{\transpose}$,
multiplier $-1$ at $(3,2)$ of $\mat{L}$. Back substitution:
$x_{3} = 2/4 = 0.5$, $x_{2} = (-2 - 3 \cdot 0.5)/(-3) = 7/6$,
$x_{1} = 3 - 2 \cdot 7/6 + 0.5 = 3 - 7/3 + 1/2 = 7/6$. Check:
$\mat{A}\vect{x} = (7/6 + 14/6 - 1/2, 14/6 + 7/6 + 1/2, -7/6 +
7/6 + 1)^{\transpose} = (3, 4, 1)^{\transpose}$. Factors: $\mat{L}
= \begin{psmallmatrix} 1 & 0 & 0 \\ 2 & 1 & 0 \\ -1 & -1 & 1
\end{psmallmatrix}$, $\mat{U}$ as above; verify $\mat{L}\mat{U} =
\mat{A}$ by direct multiplication.

\begin{exercise}
For the matrix $\mat{A} = \begin{psmallmatrix} 1 & 2 & 3 \\ 2 & 4
& 6 \\ 1 & 1 & 1 \end{psmallmatrix}$, find a basis for the column
space, the null space, and the row space. State the rank and
verify the rank-nullity theorem.
\end{exercise}

\paragraph{Solution sketch.}
Row reduce: $R_{2} \to R_{2} - 2 R_{1}$ gives a zero row;
$R_{3} \to R_{3} - R_{1}$ gives $(0, -1, -2)$. The reduced echelon
form is $\begin{psmallmatrix} 1 & 0 & -1 \\ 0 & 1 & 2 \\ 0 & 0 & 0
\end{psmallmatrix}$, so $\rank \mat{A} = 2$. Column-space basis:
the original first two columns $\{(1, 2, 1)^{\transpose}, (2, 4,
1)^{\transpose}\}$ (columns 1 and 2 are pivot columns; column 3 =
$-1 \cdot$ col 1 $+ 2 \cdot$ col 2 in REF, recovered as $(3, 6, 1)
= -(1,2,1) + 2(2,4,1) \cdot$). Row-space basis: any two
independent rows, e.g., $\{(1, 2, 3), (1, 1, 1)\}$. Null-space
basis: solve $x_{1} - x_{3} = 0$, $x_{2} + 2 x_{3} = 0$, giving
$(1, -2, 1)^{\transpose}$ (one-dimensional). Rank-nullity: $\rank
+ \dim \nul = 2 + 1 = 3 = n$.

\begin{exercise}
Compute the QR factorisation of $\mat{A} = \begin{psmallmatrix} 1
& 1 \\ 1 & 0 \\ 0 & 1 \end{psmallmatrix}$ by the Gram-Schmidt
procedure on the columns. Verify that $\mat{Q}^{\transpose}\mat{Q}
= \mat{I}_{2}$ and that $\mat{Q}\mat{R} = \mat{A}$.
\end{exercise}

\paragraph{Solution sketch.}
Column 1: $\vect{a}_{1} = (1, 1, 0)^{\transpose}$, $\norm{\vect{a}_{1}}_{2}
= \sqrt{2}$, $\vect{q}_{1} = (1, 1, 0)^{\transpose}/\sqrt{2}$.
Column 2: $\vect{a}_{2} = (1, 0, 1)^{\transpose}$, projection on
$\vect{q}_{1}$ is $\inner*{\vect{a}_{2}, \vect{q}_{1}} = 1/\sqrt{2}$,
so $\vect{a}_{2} - (1/\sqrt{2})\vect{q}_{1} = (1, 0, 1) - (1/2,
1/2, 0) = (1/2, -1/2, 1)$, norm $\sqrt{3/2}$. Thus $\vect{q}_{2} =
(1/2, -1/2, 1)/\sqrt{3/2} = (1, -1, 2)/\sqrt{6}$. Then
$\mat{R} = \begin{psmallmatrix} \sqrt{2} & 1/\sqrt{2} \\ 0 &
\sqrt{3/2} \end{psmallmatrix}$. Verify $\mat{Q}^{\transpose}\mat{Q}
= \mat{I}_{2}$ by inner products; $\mat{Q}\mat{R} = \mat{A}$ by
direct multiplication.

\begin{exercise}
Compute the determinant of $\mat{A} = \begin{psmallmatrix} 2 & 1
& 1 \\ 4 & 3 & 3 \\ 8 & 7 & 9 \end{psmallmatrix}$ by reading it
off the LU factorisation of the worked example in section 9.3.
Verify by cofactor expansion along the first row.
\end{exercise}

\paragraph{Solution sketch.}
From section 9.3, $\mat{U} = \begin{psmallmatrix} 2 & 1 & 1 \\ 0 &
1 & 1 \\ 0 & 0 & 2 \end{psmallmatrix}$ and $\mat{P} = \mat{I}$
(no row swaps), so $\det \mat{A} = \prod U_{ii} = 2 \cdot 1 \cdot
2 = 4$. Cofactor expansion: $\det \mat{A} = 2 (27 - 21) - 1 (36 -
24) + 1 (28 - 24) = 12 - 12 + 4 = 4$. The two computations agree.

\begin{exercise}
Compute the orthogonal projection of $\vect{b} = (1, 2, 3)^{\transpose}$
onto the column space of $\mat{A} = \begin{psmallmatrix} 1 & 0 \\
0 & 1 \\ 1 & 1 \end{psmallmatrix}$ using the projection formula
$\mat{A}(\mat{A}^{\transpose}\mat{A})^{-1}\mat{A}^{\transpose}
\vect{b}$. Verify that $\vect{b} - \proj\vect{b}$ is orthogonal
to each column of $\mat{A}$.
\end{exercise}

\paragraph{Solution sketch.}
$\mat{A}^{\transpose}\mat{A} = \begin{psmallmatrix} 2 & 1 \\ 1 & 2
\end{psmallmatrix}$, $(\mat{A}^{\transpose}\mat{A})^{-1} =
\frac{1}{3}\begin{psmallmatrix} 2 & -1 \\ -1 & 2 \end{psmallmatrix}$.
$\mat{A}^{\transpose}\vect{b} = (1 + 3, 2 + 3)^{\transpose} = (4,
5)^{\transpose}$. $\hat{\vect{x}} = \frac{1}{3}\begin{psmallmatrix}
2 & -1 \\ -1 & 2 \end{psmallmatrix}\begin{psmallmatrix} 4 \\ 5
\end{psmallmatrix} = \frac{1}{3}(3, 6)^{\transpose} = (1,
2)^{\transpose}$. Then $\proj_{\col(\mat{A})}\vect{b} =
\mat{A}\hat{\vect{x}} = (1, 2, 3)^{\transpose}$, so the residual
is zero: $\vect{b}$ already lies in $\col(\mat{A})$ (rows 1 and 2
of $\vect{b}$ determine $\hat{x}_{1}$ and $\hat{x}_{2}$, and row
3 of $\mat{A}\hat{\vect{x}}$ is $1 + 2 = 3 = b_{3}$). Orthogonality
holds trivially. A more instructive test case takes $\vect{b} =
(1, 2, 4)^{\transpose}$: the projection is still $(1, 2,
3)^{\transpose}$ but the residual $(0, 0, 1)^{\transpose}$ is
orthogonal to both columns.

\begin{exercise}
Compute the least-squares straight-line fit through the four
points $(0,1), (1,3), (2,4), (3,7)$ by setting up and solving the
normal equations. Compare the slope and intercept to the
closed-form Volume I Chapter 8 expressions.
\end{exercise}

\paragraph{Solution sketch.}
$\sum x_{i} = 6$, $\sum y_{i} = 15$, $\sum x_{i}^{2} = 14$,
$\sum x_{i} y_{i} = 0 + 3 + 8 + 21 = 32$, $m = 4$. Slope:
$\hat{\beta}_{1} = (4 \cdot 32 - 6 \cdot 15)/(4 \cdot 14 - 36) =
(128 - 90)/20 = 38/20 = 1.9$. Intercept: $\bar{x} = 1.5$,
$\bar{y} = 3.75$, $\hat{\beta}_{0} = 3.75 - 1.9 \cdot 1.5 = 3.75 -
2.85 = 0.9$. The normal equations route, $\mat{A}^{\transpose}
\mat{A}\hat{\vect{c}} = \mat{A}^{\transpose}\vect{b}$ with
$\mat{A}^{\transpose}\mat{A} = \begin{psmallmatrix} 4 & 6 \\ 6 &
14 \end{psmallmatrix}$ and $\mat{A}^{\transpose}\vect{b} = (15,
32)^{\transpose}$, yields the same $(0.9, 1.9)^{\transpose}$.

\begin{exercise}
For the matrix $\mat{A} = \begin{psmallmatrix} 1 & 1 \\ 1 & 1.0001
\end{psmallmatrix}$, compute $\mat{A}^{-1}$, the singular values
$\sigma_{1}, \sigma_{2}$, and the condition number $\kappa(\mat{A})$.
State how many decimal digits of accuracy the rule of thumb predicts
in the solution to $\mat{A}\vect{x} = \vect{b}$ given input data
accurate to $10^{-10}$.
\end{exercise}

\paragraph{Solution sketch.}
$\det \mat{A} = 1.0001 - 1 = 10^{-4}$. $\mat{A}^{-1} = 10^{4}
\begin{psmallmatrix} 1.0001 & -1 \\ -1 & 1 \end{psmallmatrix}$.
For singular values, $\mat{A}^{\transpose}\mat{A} = \begin{psmallmatrix}
2 & 2.0001 \\ 2.0001 & 2.00020001 \end{psmallmatrix}$ has trace
$\approx 4.0002$ and determinant $\approx 10^{-8}$, so eigenvalues
$\lambda_{1} \approx 4.0002$, $\lambda_{2} \approx 2.5 \times
10^{-9}$, giving $\sigma_{1} \approx 2.000$, $\sigma_{2} \approx
5.0 \times 10^{-5}$, $\kappa(\mat{A}) = \sigma_{1}/\sigma_{2}
\approx 4.0 \times 10^{4}$. Rule of thumb: $\log_{10}\kappa
\approx 4.6$. Data accurate to $10^{-10}$ carry $10$ digits;
$10 - 4.6 \approx 5.4$ digits in the solution. The reader should
not trust more than five digits of $\hat{\vect{x}}$ for this
near-singular matrix.

\begin{exercise}
Compute the angle between the vectors $\vect{u} = (1, 2, 2)^{\transpose}$
and $\vect{v} = (2, -1, 2)^{\transpose}$ using the inner-product
formula. Verify the Cauchy-Schwarz inequality numerically.
\end{exercise}

\paragraph{Solution sketch.}
$\inner*{\vect{u}, \vect{v}} = 2 - 2 + 4 = 4$. $\norm{\vect{u}}_{2}
= \sqrt{1 + 4 + 4} = 3$. $\norm{\vect{v}}_{2} = \sqrt{4 + 1 + 4}
= 3$. $\cos\theta = 4/9 \approx 0.444$, so $\theta \approx
\arccos(0.444) \approx 63.6^{\circ}$. Cauchy-Schwarz: $\abs{\inner*{
\vect{u}, \vect{v}}} = 4 \leq 9 = \norm{\vect{u}}_{2}
\norm{\vect{v}}_{2}$. Equality would hold only if $\vect{u}$ and
$\vect{v}$ were parallel; they are not.

\subsection*{Derivation}\pathtag{standard}

\begin{exercise}
Show that the column space and the null space of a matrix $\mat{A}
\in \R^{m \times n}$ satisfy $\dim \col(\mat{A}) + \dim \nul(\mat{A})
= n$. State, but do not derive, the analogous identity for the row
and left-null spaces.
\end{exercise}

\paragraph{Solution sketch.}
Row reduce $\mat{A}$ to its reduced row-echelon form $\mat{R}$.
Let $r$ be the number of pivot columns. The pivot columns of
$\mat{A}$ (not of $\mat{R}$) form a basis for $\col(\mat{A})$, so
$\dim \col(\mat{A}) = r$. The non-pivot columns yield $n - r$ free
parameters; each free parameter contributes one basis vector of
$\nul(\mat{A})$, so $\dim \nul(\mat{A}) = n - r$. Summing gives
$r + (n - r) = n$. The analogous identity is $\dim
\col(\mat{A}^{\transpose}) + \dim \nul(\mat{A}^{\transpose}) = m$;
row rank equals column rank equals $r$, the central content of the
four-subspaces picture (Figure~\ref{fig:vol02:ch09:four-subspaces}).

\begin{exercise}
Show that the matrix $\mat{P} = \mat{A}(\mat{A}^{\transpose}\mat{A})^{-1}
\mat{A}^{\transpose}$, when $\mat{A}$ has full column rank,
satisfies $\mat{P}^{2} = \mat{P}$ and $\mat{P}^{\transpose} =
\mat{P}$. Use only the algebraic properties of transpose and
inverse.
\end{exercise}

\paragraph{Solution sketch.}
Symmetry: $\mat{P}^{\transpose} = (\mat{A}^{\transpose})^{\transpose}
((\mat{A}^{\transpose}\mat{A})^{-1})^{\transpose}\mat{A}^{\transpose}
= \mat{A}((\mat{A}^{\transpose}\mat{A})^{\transpose})^{-1}
\mat{A}^{\transpose}$. Since $(\mat{A}^{\transpose}\mat{A})^{\transpose}
= \mat{A}^{\transpose}\mat{A}$, $\mat{P}^{\transpose} = \mat{P}$.
Idempotence: $\mat{P}^{2} = \mat{A}(\mat{A}^{\transpose}\mat{A})^{-1}
\mat{A}^{\transpose}\mat{A}(\mat{A}^{\transpose}\mat{A})^{-1}
\mat{A}^{\transpose} = \mat{A}(\mat{A}^{\transpose}\mat{A})^{-1}
\mat{A}^{\transpose} = \mat{P}$; the inner $\mat{A}^{\transpose}
\mat{A}$ cancels its inverse.

\begin{exercise}
Starting from $\norm{\mat{A}\vect{x} - \vect{b}}_{2}^{2}$, derive
the normal equations $\mat{A}^{\transpose}\mat{A}\hat{\vect{x}} =
\mat{A}^{\transpose}\vect{b}$ by setting the gradient with respect
to $\vect{x}$ to zero. State the assumption on $\mat{A}$ that
guarantees a unique minimiser.
\end{exercise}

\paragraph{Solution sketch.}
Expand: $f(\vect{x}) = \vect{x}^{\transpose}\mat{A}^{\transpose}
\mat{A}\vect{x} - 2\vect{x}^{\transpose}\mat{A}^{\transpose}\vect{b}
+ \vect{b}^{\transpose}\vect{b}$. The gradient is $\nabla f =
2\mat{A}^{\transpose}\mat{A}\vect{x} - 2\mat{A}^{\transpose}\vect{b}$.
Setting $\nabla f = \vect{0}$ gives the normal equations. The
Hessian is $2\mat{A}^{\transpose}\mat{A}$; the minimiser is unique
when this Hessian is positive definite, which holds if and only if
$\mat{A}$ has full column rank.

\begin{exercise}
Show that a square matrix $\mat{A}$ is invertible if and only if
$\det \mat{A} \neq 0$, by combining the multiplicative property
$\det(\mat{A}\mat{B}) = \det \mat{A} \cdot \det \mat{B}$ with
$\det \mat{I} = 1$.
\end{exercise}

\paragraph{Solution sketch.}
$(\Rightarrow)$ If $\mat{A}$ is invertible, $\mat{A}\mat{A}^{-1} =
\mat{I}$ and $\det \mat{A} \cdot \det \mat{A}^{-1} = \det \mat{I}
= 1$, forcing $\det \mat{A} \neq 0$. $(\Leftarrow)$ If $\det \mat{A}
\neq 0$, then the columns of $\mat{A}$ are linearly independent (a
linear dependence would collapse the image parallelepiped to lower
dimension, forcing $\det \mat{A} = 0$). For a square matrix,
linearly independent columns span $\R^{n}$, so $\mat{A}$ defines a
bijection and is invertible.

\begin{exercise}
Show that for a full-column-rank matrix $\mat{A}$, the Gram matrix
$\mat{A}^{\transpose}\mat{A}$ is symmetric and positive definite.
Use the symmetric and positive-definite properties to argue that
the Cholesky factorisation $\mat{A}^{\transpose}\mat{A} =
\mat{L}\mat{L}^{\transpose}$ exists.
\end{exercise}

\paragraph{Solution sketch.}
Symmetry: $(\mat{A}^{\transpose}\mat{A})^{\transpose} =
\mat{A}^{\transpose}\mat{A}$. Positive definiteness: for any
$\vect{x} \neq \vect{0}$, $\vect{x}^{\transpose}\mat{A}^{\transpose}
\mat{A}\vect{x} = \norm{\mat{A}\vect{x}}_{2}^{2} \geq 0$, with
equality if and only if $\mat{A}\vect{x} = \vect{0}$, which under
full column rank forces $\vect{x} = \vect{0}$. Existence of
Cholesky: a symmetric positive-definite matrix has all positive
eigenvalues; the constructive Cholesky algorithm proceeds column
by column, each diagonal pivot $\sqrt{L_{kk}^{2}}$ being positive
by induction on the leading principal submatrix's positive
definiteness, terminating with a real lower-triangular $\mat{L}$
such that $\mat{A}^{\transpose}\mat{A} = \mat{L}\mat{L}^{\transpose}$.

\begin{exercise}
Show that the condition number satisfies $\kappa(\mat{A}^{\transpose}
\mat{A}) = \kappa(\mat{A})^{2}$ for a full-column-rank $\mat{A}$.
Use the singular-value characterisation $\kappa(\mat{A}) =
\sigma_{1}/\sigma_{n}$ and the relation between the singular values
of $\mat{A}$ and those of $\mat{A}^{\transpose}\mat{A}$.
\end{exercise}

\paragraph{Solution sketch.}
The singular values of $\mat{A}^{\transpose}\mat{A}$ are the
squares of the singular values of $\mat{A}$: if $\mat{A} =
\mat{U}\mat{\Sigma}\mat{V}^{\transpose}$, then $\mat{A}^{\transpose}
\mat{A} = \mat{V}\mat{\Sigma}^{2}\mat{V}^{\transpose}$, which has
singular values $\sigma_{i}^{2}$ (it is symmetric positive definite,
so eigenvalues coincide with singular values). Therefore
$\kappa(\mat{A}^{\transpose}\mat{A}) = \sigma_{1}^{2}/\sigma_{n}^{2}
= \kappa(\mat{A})^{2}$.

\begin{exercise}
Show that the Householder reflector $\mat{H} = \mat{I} - 2\vect{v}
\vect{v}^{\transpose}/\norm{\vect{v}}_{2}^{2}$ is orthogonal and
symmetric, where $\vect{v}$ is any nonzero vector. State the
geometric action of $\mat{H}$ on a vector $\vect{x}$.
\end{exercise}

\paragraph{Solution sketch.}
Let $\alpha = 2/\norm{\vect{v}}_{2}^{2}$. Symmetry:
$\mat{H}^{\transpose} = \mat{I} - \alpha (\vect{v}\vect{v}^{\transpose})^{\transpose}
= \mat{I} - \alpha \vect{v}\vect{v}^{\transpose} = \mat{H}$.
Orthogonality: $\mat{H}^{\transpose}\mat{H} = \mat{H}^{2} = \mat{I}
- 2\alpha\vect{v}\vect{v}^{\transpose} + \alpha^{2}\vect{v}
(\vect{v}^{\transpose}\vect{v})\vect{v}^{\transpose} = \mat{I} -
2\alpha\vect{v}\vect{v}^{\transpose} + 2\alpha\vect{v}\vect{v}^{\transpose}
= \mat{I}$, using $\alpha \vect{v}^{\transpose}\vect{v} = 2$.
Geometry: $\mat{H}\vect{x} = \vect{x} - 2(\vect{v}^{\transpose}
\vect{x}/\norm{\vect{v}}_{2}^{2})\vect{v}$ is the reflection of
$\vect{x}$ through the hyperplane orthogonal to $\vect{v}$. The
component of $\vect{x}$ parallel to $\vect{v}$ is negated; the
component orthogonal to $\vect{v}$ is preserved.

\subsection*{Estimation}\pathtag{core}

\begin{exercise}
Estimate the floating-point operations required to solve a dense
linear system of size $n = 10^{4}$ by LU factorisation. Compare to
the operations required to solve the same system by inverting
$\mat{A}$ explicitly and multiplying. State the working rule of
thumb that makes explicit inversion the wrong choice.
\end{exercise}

\paragraph{Solution sketch.}
LU costs $\tfrac{2}{3}n^{3} = \tfrac{2}{3} \cdot 10^{12} \approx
6.7 \times 10^{11}$ flops for the factorisation plus $2 n^{2} = 2
\times 10^{8}$ flops for forward+back substitution. Total $\approx
6.7 \times 10^{11}$. Explicit inversion costs $\tfrac{2}{3}n^{3}$
for the LU plus $n$ back-solves to recover the columns of
$\mat{A}^{-1}$ ($n \cdot 2 n^{2} = 2 \times 10^{12}$ flops) plus
the matrix-vector multiply $\mat{A}^{-1}\vect{b}$ ($2 n^{2}$ flops).
Total $\approx 2.7 \times 10^{12}$, four times the cost. The rule
of thumb: never form $\mat{A}^{-1}$ explicitly; always solve
$\mat{A}\vect{x} = \vect{b}$ directly via factorisation. The
factor-of-four cost penalty is the small part; the larger penalty
is the numerical inaccuracy of $\mat{A}^{-1}$ relative to a direct
solve.

\begin{exercise}
Estimate the condition number of the $5 \times 5$ Hilbert matrix
$H_{ij} = 1/(i+j-1)$, given that the Hilbert matrix is the
canonical ill-conditioned example. State, by the rule of thumb,
how many decimal digits of accuracy a double-precision solution
to $\mat{H}\vect{x} = \vect{b}$ retains.
\end{exercise}

\paragraph{Solution sketch.}
Tabulated Hilbert condition numbers grow roughly as $\kappa(\mat{H}_{n})
\approx e^{3.5 n}$ for small $n$. For $n = 5$, $\kappa(\mat{H}_{5})
\approx 4.8 \times 10^{5}$ (working value to within a factor of
two). Rule of thumb: $\log_{10}\kappa \approx 5.7$. Double
precision carries 16 digits; the solution retains roughly $16 -
5.7 \approx 10$ digits, comfortable for engineering work. The
$10 \times 10$ Hilbert is the more relevant warning: $\kappa
\approx 1.6 \times 10^{13}$ leaves only two to three reliable
digits, which is the classic exam for any newly written solver.

\begin{exercise}
Estimate the condition number of the $11 \times 11$ Vandermonde
matrix on the equally-spaced points $x_{i} = (i-1)/10$ for $i = 1,
\ldots, 11$. State whether you would trust a degree-10 polynomial
fit on these points in double precision and which alternative
formulation you would recommend.
\end{exercise}

\paragraph{Solution sketch.}
By the monomial-basis rule of thumb $\log_{10}\kappa(\mat{V})
\approx 1.5 p$, the degree-10 fit on $[0, 1]$ has $\kappa
\approx 10^{15}$. The Gram matrix would have $\kappa \approx
10^{30}$, dwarfing double precision. The QR formulation, governed
by $\kappa(\mat{V})$, leaves about one digit of accuracy.
Recommendation: refuse the monomial basis. Rescale $x$ to $[-1,
1]$ (recovers about $3 \times 10 = 30$ orders of magnitude in
$\log_{10}\kappa$ as Figure~\ref{fig:vol02:ch09:vandermonde-conditioning}
shows) or switch to the Chebyshev basis on $[-1, 1]$ (keeps
$\kappa$ near unity through degree $10$).

\begin{exercise}
A least-squares regression has a design matrix with condition
number $10^{6}$. The data are accurate to four decimal digits.
Estimate the relative accuracy of the fitted coefficients in the
QR formulation and in the normal-equations formulation, and state
which is reportable.
\end{exercise}

\paragraph{Solution sketch.}
Data accuracy: $4$ digits, so $\norm{\Delta\vect{b}}/\norm{\vect{b}}
\approx 10^{-4}$. QR: governed by $\kappa(\mat{A}) = 10^{6}$, so
relative error of coefficients $\lesssim 10^{6} \cdot 10^{-4} =
10^{-2}$, i.e.\ two digits. Marginal but reportable, with the
caveat stated. Normal equations: governed by $\kappa(\mat{A}^{\transpose}
\mat{A}) = 10^{12}$, so relative error $\lesssim 10^{12} \cdot
10^{-4} = 10^{8}$, i.e.\ no digits. Not reportable. The engineer
reports the QR result with its design-matrix condition number
disclosed and explicitly refuses the normal-equations result.

\begin{exercise}
Estimate the dimension of the displacement space of a planar truss
with $n = 50$ joints, of which $10$ are pinned to the ground (each
removing two degrees of freedom) and $5$ are constrained to a
roller (each removing one degree of freedom). State the rank check
that turns this estimate into an exact count.
\end{exercise}

\paragraph{Solution sketch.}
Unconstrained displacement space: $2 n = 100$. Pin constraints
remove $2 \times 10 = 20$. Roller constraints remove $1 \times 5
= 5$. Estimated displacement dimension: $100 - 25 = 75$. The
constraint count is exact only when all constraints are linearly
independent; if two pins are collinear with a shared bar or a
roller is parallel to a pin's reaction direction, the constraint
matrix loses rank. The rank check is: assemble the constraint
matrix $\mat{C} \in \R^{25 \times 100}$ and compute its rank
numerically; if $\rank \mat{C} < 25$, the true displacement
dimension is $100 - \rank \mat{C} > 75$, and the structure has
hidden mechanisms.

\begin{exercise}
Estimate the wall-clock time to solve a dense linear system of
size $n = 5000$ on a single core that sustains $10^{10}$ flops per
second. Then estimate the time for $n = 50000$, using the cubic
scaling. State the largest $n$ for which a one-hour solve is
feasible at this throughput, and the implication for choosing
between dense and sparse formulations.
\end{exercise}

\paragraph{Solution sketch.}
LU on $n = 5000$: $\tfrac{2}{3} n^{3} \approx 8.3 \times 10^{10}$
flops, $\approx 8.3 \,\text{s}$ at $10^{10}\,\text{flop}/\text{s}$.
Cubic scaling for $n = 50000$: $\tfrac{2}{3} \cdot (5 \times 10^{4})^{3}
\approx 8.3 \times 10^{13}$ flops, $\approx 8.3 \times 10^{3}
\,\text{s} \approx 2.3\,\text{h}$. One-hour budget at $10^{10}
\,\text{flop}/\text{s}$ is $3.6 \times 10^{13}$ flops; solving
$\tfrac{2}{3} n^{3} = 3.6 \times 10^{13}$ gives $n \approx 38000$.
For larger $n$, dense LU is infeasible; the engineer must exploit
structure (banded, sparse, low-rank, Krylov-iterative) or accept
multi-hour solves. The cubic wall is the central reason finite-element
codes use sparse factorisations or preconditioned conjugate gradient.

\subsection*{Simulation}\pathtag{mastery}

\begin{exercise}
In a programming environment of the reader's choice, generate a
$100 \times 5$ Vandermonde matrix on $100$ equally spaced points
in $[0, 1]$. Compute its condition number. Then rescale the points
to $[-1, 1]$ and recompute. Quantify the conditioning improvement
and explain it in terms of the orthogonality of the column space.
\end{exercise}

\paragraph{Solution sketch.}
The reference script \path{code/condition_visualiser.py} at degree
$p = 4$ on $m = 100$ points gives $\kappa(\mat{V}_{[0,1]}) \approx
1.6 \times 10^{4}$ and $\kappa(\mat{V}_{[-1,1]}) \approx 70$, an
improvement of more than two orders of magnitude. The explanation:
on $[0, 1]$ all monomials $x^{k}$ are non-negative and increase
together, so their columns correlate strongly. On $[-1, 1]$ the
odd monomials are antisymmetric and the even monomials are
symmetric about zero; columns of opposite parity are exactly
orthogonal under the uniform inner product, and columns of the
same parity correlate less than on $[0, 1]$. The Chebyshev basis
(third column of the script) makes the orthogonality exact by
construction, taking $\kappa$ to near unity for any degree the
sampling supports.

\begin{exercise}
Simulate a least-squares fit of a degree-$10$ polynomial to $200$
points generated from $y = \sin(2\pi x) + \epsilon$ with
$\epsilon \sim \mathcal{N}(0, 0.05^{2})$ and $x \in [0, 1]$.
Implement the normal-equations and QR formulations. Compare the
fitted coefficients and the residual norms. Identify the
coefficient that disagrees most between the two methods and explain
the disagreement using the conditioning analysis of section 9.7.
\end{exercise}

\paragraph{Solution sketch.}
Typical run with the reference script
\path{code/least_squares_three_ways.py}: $\kappa(\mat{V}) \approx
10^{15}$ on $[0, 1]$ at degree $10$, so $\kappa(\mat{V}^{\transpose}
\mat{V}) \approx 10^{30}$, well past double precision. The residual
norms agree (both methods fit the data equally well in the 2-norm),
but the coefficient vectors diverge sharply, with the largest
disagreement typically in the highest-degree coefficient $c_{10}$
where the basis function $x^{10}$ is most correlated with the
lower powers and the conditioning trouble concentrates. The QR
coefficients are the trustworthy ones, governed by $\kappa(\mat{V})$;
the normal-equations coefficients are nonsense, governed by
$\kappa(\mat{V}^{\transpose}\mat{V})$. Rescaling $x$ to $[-1, 1]$
brings both methods into agreement to several digits and is the
default fix.

\begin{exercise}
Simulate a perturbation experiment: generate a $20 \times 20$
random matrix $\mat{A}$, compute the solution $\vect{x}$ of
$\mat{A}\vect{x} = \vect{b}$, then perturb $\mat{A}$ by a random
matrix of size $\norm{\Delta\mat{A}}/\norm{\mat{A}} = 10^{-8}$ and
recompute. Plot the relative error $\norm{\Delta\vect{x}}/
\norm{\vect{x}}$ against the condition number $\kappa(\mat{A})$
over $100$ random matrices and verify the sensitivity bound of
section 9.7.
\end{exercise}

\paragraph{Solution sketch.}
Typical run: most random Gaussian $20 \times 20$ matrices have
$\kappa(\mat{A})$ between $10$ and $10^{4}$; with $\norm{\Delta
\mat{A}}/\norm{\mat{A}} = 10^{-8}$ the predicted relative error in
$\vect{x}$ is $\kappa \cdot 10^{-8}$, ranging from $10^{-7}$ to
$10^{-4}$. The scatter plot of $\log_{10}(\norm{\Delta\vect{x}}/
\norm{\vect{x}})$ against $\log_{10}\kappa$ shows a band of points
lying near or below the line $y = x - 8$, with occasional
outliers below the bound (the bound is tight only in the worst
direction). To force higher conditioning the engineer can
construct $\mat{A} = \mat{Q}_{1}\mat{D}\mat{Q}_{2}^{\transpose}$
with $\mat{Q}_{i}$ random orthogonal and $\mat{D}$ diagonal with a
controlled $\sigma_{1}/\sigma_{n}$.

\begin{exercise}
Implement the Householder QR algorithm from scratch on a $6 \times
4$ matrix and verify the result against the library routine. The
goal is to see, by working through the algorithm, why Householder
reflections are numerically stable while classical Gram-Schmidt
is not.
\end{exercise}

\paragraph{Solution sketch.}
The Householder algorithm, column $k$: form $\vect{x}$ as the
sub-column $\mat{A}_{k:m, k}$, compute the reflector vector $\vect{v}
= \vect{x} - \sgn(x_{1})\norm{\vect{x}}_{2}\vect{e}_{1}$ (note the
sign chosen to avoid cancellation), build $\mat{H}_{k} = \mat{I} -
2 \vect{v}\vect{v}^{\transpose}/\norm{\vect{v}}_{2}^{2}$, and
update $\mat{A}_{k:m, k:n} \to \mat{H}_{k}\mat{A}_{k:m, k:n}$. After
$n = 4$ reflections, $\mat{A}$ is upper-triangular and $\mat{R}$
sits in its upper part; $\mat{Q} = \mat{H}_{1}\mat{H}_{2}
\mat{H}_{3}\mat{H}_{4}$ implicitly. Verify $\norm{\mat{Q}\mat{R} -
\mat{A}}_{F} \lesssim 10^{-15} \cdot \norm{\mat{A}}_{F}$ against
\path{numpy.linalg.qr}. The stability follows from orthogonality of
$\mat{H}_{k}$: every step is a unitary transformation that does not
amplify floating-point error. Classical Gram-Schmidt loses
orthogonality in the columns of $\mat{Q}$ because the projection
subtractions are not unitary; the error compounds across columns.

\subsection*{Design}\pathtag{standard}

\begin{exercise}
Design a procedure for fitting a polynomial of degree $p$ to a
noisy dataset that explicitly checks the conditioning of the
design matrix and chooses among the normal-equations, QR, and
truncated-SVD formulations based on a stated condition-number
threshold. State the thresholds used and justify them with the
rule of thumb of section 9.7.
\end{exercise}

\paragraph{Discussion.}
A defensible procedure: (1) rescale $x$ to $[-1, 1]$ before forming
the design matrix; (2) compute $\kappa(\mat{V})$ once (via SVD or
the library routine); (3) if $\kappa(\mat{V}) < 10^{3.5}$, the
Gram condition number is below $10^{7}$ and normal equations are
defensible at four to five digits in double precision; use
normal equations for speed; (4) if $10^{3.5} \leq \kappa(\mat{V})
< 10^{7.5}$, use QR; (5) if $\kappa(\mat{V}) \geq 10^{7.5}$,
compute the SVD and apply truncation at $\epsilon = 10^{-12}$
times $\sigma_{1}$, then report the truncation cutoff and the
number of singular values retained. The thresholds follow the
rule of thumb: each factor of $10$ in $\kappa$ costs one digit;
sixteen digits in double precision, minus the data's precision,
gives a budget that the squared (Gram) conditioning consumes
twice as fast. Always report $\kappa(\mat{V})$ and the chosen
method with the fit.

\begin{exercise}
Design a verification check for a numerical-conversion operation
between two representations (e.g., 32-bit float to 16-bit integer)
in a flight-software context. The check should state the
operation's precondition (range), the verification method, and
the action when verification fails. Reference the Ariane 5 case
of section 9.8 in the design rationale.
\end{exercise}

\paragraph{Discussion.}
Precondition: the input value $x$ satisfies $-32768 \leq x \leq
32767$ over the entire flight envelope at the operation's call
sites, with margin chosen by the safety case (a $10\%$ margin on
the trajectory's worst-case excursion is a defensible default).
Verification method: (1) static analysis over the flight envelope
to derive an upper bound on $\abs{x}$ from the propulsion,
guidance, and trajectory models; (2) runtime range check
immediately before conversion that asserts $\abs{x} \leq 32767$;
(3) on assertion failure, route to the safe-state handler rather
than allowing the IEEE float to convert to undefined-on-overflow
integer behaviour. The action on verification failure must not be
``log and continue,'' which is what the Ariane 5 software did in
effect by raising an unhandled exception that propagated to SRI
shutdown. The defensible action is graceful degradation to a
fallback computation that does not require the converted value, or
explicit transition to a defined safe mode. The rationale references
\cite{acc:ariane501-ifr}: the failure mode is the absence of a
range-check on the conversion's input under a new flight envelope.

\begin{exercise}
Design a sparse-matrix linear solver for a finite-element
stiffness problem of size $n = 10^{6}$ with bandwidth $b = 100$.
State the storage requirement, the operation count, and the
factorisation method (banded LU, Cholesky, or sparse-direct). Mark
this exercise mastery if the chapter has not introduced sparse
factorisation; provide the storage and complexity scaling
otherwise.
\end{exercise}

\paragraph{Discussion.}
Storage: a banded matrix of size $n$ and bandwidth $b$ stores
$(2 b + 1) n$ entries, approximately $200 \times 10^{6} = 2
\times 10^{8}$ doubles, or $1.6 \,\text{GB}$. Banded LU
factorisation costs $\oom{n b^{2}} = 10^{10}$ flops, about one
second on a contemporary core; the fill-in is bounded by
$\oom{n b}$. If the stiffness matrix is symmetric positive
definite (it is, for elastic problems), banded Cholesky halves
the storage and the flop count and is the production choice. For
$b$ larger or the structure less band-like, a sparse-direct
solver with nested-dissection ordering (MUMPS, PARDISO, SuiteSparse)
trades factorisation cost for fill-in control. Iterative methods
(preconditioned conjugate gradient) become competitive when $n$
exceeds $10^{7}$ or memory becomes the binding constraint.

\subsection*{Diagnosis and reverse engineering}\pathtag{core}

\begin{exercise}
A colleague reports a least-squares fit with residual norm
$10^{-12}$ on data measured to three decimal digits. Identify
the three things that should make the engineer suspicious of the
report and propose follow-up questions.
\end{exercise}

\paragraph{Solution sketch.}
The three signals. First, the residual is far smaller than the data
noise floor: $10^{-12}$ on $10^{-3}$-precision data means the fit
has driven through the noise, which is overfitting. Second, no
condition number is reported; the residual alone hides whether the
fit is meaningful or whether the design matrix has consumed all the
degrees of freedom. Third, the model class is not stated; a
polynomial of degree $m - 1$ on $m$ data points interpolates
exactly, residual zero, with no predictive value. Follow-up
questions: (a) what is the model degree relative to the sample
size; (b) what is $\kappa(\mat{V})$ of the design matrix; (c)
what is the residual on a held-out sample; (d) which factorisation
was used.

\begin{exercise}
A linear system $\mat{A}\vect{x} = \vect{b}$ in a control
application produces a solution that changes by a factor of two
when the right-hand side changes by one part in $10^{4}$. Diagnose
the conditioning regime of $\mat{A}$ and propose two reformulations
that would stabilise the computation.
\end{exercise}

\paragraph{Solution sketch.}
By the sensitivity bound, $\norm{\Delta\vect{x}}/\norm{\vect{x}}
\approx \kappa(\mat{A}) \cdot \norm{\Delta\vect{b}}/\norm{\vect{b}}$.
Observed $\norm{\Delta\vect{x}}/\norm{\vect{x}} \approx 1$ at
$\norm{\Delta\vect{b}}/\norm{\vect{b}} = 10^{-4}$ implies
$\kappa(\mat{A}) \approx 10^{4}$. The system is moderately
ill-conditioned. Two reformulations. (1) Regularise: replace
$\mat{A}\vect{x} = \vect{b}$ by $(\mat{A}^{\transpose}\mat{A} +
\lambda\mat{I})\vect{x} = \mat{A}^{\transpose}\vect{b}$ (Tikhonov,
or ridge in regression), with $\lambda$ chosen near $10^{-4}
\norm{\mat{A}}_{2}^{2}$ to cap the conditioning. (2) Reduce to
the dominant singular components: compute $\mat{A} = \mat{U}
\mat{\Sigma}\mat{V}^{\transpose}$ and truncate singular values
below a cutoff $\epsilon \sigma_{1}$ (truncated SVD, the
chapter's preferred refusal). Both reformulations declare the
problem's effective rank explicitly rather than letting the
floating-point arithmetic discover it.

\begin{exercise}
A polynomial fit of degree $8$ to $40$ data points produces wildly
oscillating coefficients despite a small residual. Identify the
conditioning failure that produces oscillation, propose two
alternative formulations (rescaling, orthogonal-polynomial basis,
truncated SVD) and state which would be the right choice without
running the computation.
\end{exercise}

\paragraph{Solution sketch.}
The conditioning failure: Vandermonde at degree $8$ on raw $x$
values in an interval that does not include zero has $\kappa
\approx 10^{12}$ in the monomial basis. The Gram matrix is past
double-precision usable conditioning. Large oscillations in the
coefficients with small residual are the signature: cancellation
in the linear combination $\mat{V}\hat{\vect{c}}$ reproduces the
data while individual coefficients are nonsense. First alternative,
rescaling to $[-1, 1]$: $\kappa$ drops by about $3 \times 8 = 24$
orders of magnitude in the rule-of-thumb extrapolation, recovering
six to eight digits per coefficient. Second alternative, the
Chebyshev basis on the rescaled interval: keeps $\kappa$ near
unity at every degree. Without running the computation, the
right default is the Chebyshev basis: orthogonal-polynomial
formulations were designed for this problem and are the canonical
fix for polynomial-fit conditioning trouble.

\begin{exercise}
A finite-element analysis of a truss returns a stiffness matrix
$\mat{K}$ with $\det \mat{K} = 10^{-30}$. Identify the physical
meaning of the small determinant, diagnose whether the structure
is genuinely under-constrained or merely poorly numbered, and
propose the diagnostic check that distinguishes the two.
\end{exercise}

\paragraph{Solution sketch.}
The physical meaning of $\det \mat{K} \to 0$ is a rigid-body mode:
a displacement pattern $\vect{u}$ for which $\mat{K}\vect{u} =
\vect{0}$, meaning the structure can move with no internal strain.
The diagnostic distinction. Genuine under-constraint: the
rigid-body modes are physical (the structure is missing supports,
or a hinge has been left in). Poor numbering: the small determinant
is an artefact of bad row scaling (the diagonal entries differ by
$10^{20}$ because some rows have units of $\si{\newton\per\meter}$
and others $\si{\newton\meter\per\radian}$). The diagnostic check
is the SVD of $\mat{K}$: compute the singular vectors associated
with the smallest singular values. If those vectors describe
sensible rigid-body translations or rotations, the structure is
genuinely under-constrained. If they describe coordinate
combinations that follow the row-unit scaling pattern, the matrix
is poorly numbered and a row-equilibration preconditioner
($\mat{D}^{-1}\mat{K}\vect{x} = \mat{D}^{-1}\vect{b}$, with
$\mat{D}$ diagonal of row norms) restores defensible conditioning.

\subsection*{Failure analysis}\pathtag{core}

\begin{exercise}
Argue, in 200 to 300 words, why the Ariane 5 Flight 501 numerical
conversion should have been refused, given the verification
discipline this chapter develops. Identify the specific
precondition that the conversion's input violated and the
verification procedure that would have detected the violation
during qualification rather than at $H_{0} + 39 \,\si{\second}$.
\end{exercise}

\paragraph{Solution sketch.}
The precondition: $\mathrm{BH} \in [-32768, 32767]$ at every
call site of the conversion across the operational flight
envelope. The Ariane 5 trajectory developed a horizontal velocity
profile that took $\mathrm{BH}$ above $32767$ at about $H_{0} +
36.7\,\si{\second}$, violating the precondition. The verification
procedure: a range analysis of every internal variable in the
SRI software against the Ariane 5 worst-case envelope, using the
trajectory profiles already produced by the launcher's flight
dynamics team. A two-page tabulation of $(\text{variable, upper
bound, lower bound, conversion target})$ for the seven internal
variables in the alignment function would have flagged BH (and the
other two unprotected variables) as outside their conversion
targets. The discipline this chapter develops names that
tabulation as a verification artefact: the conversion is a numerical
operation with a precondition, and the precondition must be
verified against the new envelope when the code is ported. The
Inquiry Board's recommendations (range analysis against the new
envelope, defensive arithmetic on flight-critical paths, exception
handling that does not propagate to system shutdown) read as
restatements of the same principle in the software-engineering
register. The refusal was available; what was missing was the
process step that connected the trajectory analysis to the SRI
software qualification.

\begin{exercise}
A regression analysis on a clinical dataset reports a coefficient
estimate $\hat{\beta} = 2.4$ with standard error $0.3$. A
replication on a slightly different sample produces $\hat{\beta}
= -1.1$ with standard error $0.4$. Argue, in 200 to 300 words,
whether the disagreement is consistent with conditioning trouble
in the design matrix or with a real difference in the underlying
populations, and identify the diagnostic that would distinguish
the two.
\end{exercise}

\paragraph{Solution sketch.}
The reported difference $(-1.1) - 2.4 = -3.5$ with combined
standard error $\sqrt{0.3^{2} + 0.4^{2}} = 0.5$ gives $z = -7$:
statistically incompatible under the null of identical populations
and identical conditioning. Two competing explanations. Conditioning:
if the design matrix has $\kappa \approx 10^{4}$ or higher, small
changes in sample composition (which change the design matrix
columns) can swing the fitted coefficient by orders of magnitude.
The standard errors reported by the regression software, computed
from the variance-covariance matrix, may themselves be unreliable
when the Gram matrix is ill-conditioned. Real population difference:
the two samples are drawn from genuinely different distributions
(different recruitment, different intervention era, different
measurement protocol), and the coefficient really does change
sign. The diagnostic: compute $\kappa(\mat{X})$ for both design
matrices and the variance inflation factor for the coefficient of
interest in each. If $\kappa > 10^{3}$ or $\mathrm{VIF} > 10$,
conditioning trouble is in play and the reported standard errors
are not trustworthy. If both design matrices are well conditioned
($\kappa < 10^{2}$, $\mathrm{VIF} < 5$), the disagreement reflects
a real population difference and the next investigation is into
the sampling protocol or the unmodelled covariates that vary
between the two cohorts.

\begin{exercise}
A simulation of a chemical reactor, using a $1000 \times 1000$
linear system at every time step, begins to produce nonsense
after $1000$ time steps although the early steps appeared
correct. Argue, in 200 to 300 words, whether the trouble is
likely to be ill-conditioning of the time-step matrix, error
accumulation across steps, or a genuine instability of the
underlying physical model, and propose the diagnostic that
identifies which.
\end{exercise}

\paragraph{Solution sketch.}
Three candidate causes, ranked by their signature. (1)
Time-step-matrix conditioning: $\kappa$ near the precision limit
makes each solve lose digits, with the loss compounding linearly
or super-linearly across steps. Signature: errors grow steadily
from step one but stay small until step $\approx 1000$ when the
cumulative digit loss exhausts double precision. (2) Error
accumulation: each step adds $\oom{\epsilon_{\text{mach}}}$
round-off; over $1000$ steps this is $\oom{1000 \epsilon_{\text{mach}}}
\approx 10^{-13}$, far too small to cause visible failure on its
own. Likely a contributing factor, not the primary cause. (3)
Physical instability: the underlying reactor model has a positive
real eigenvalue (autocatalytic kinetics, exothermic runaway) whose
exponential growth saturates the numerical scheme at a definite
time. Signature: the failure time depends on the initial
perturbation and matches a predictable $\exp(\lambda t)$ growth.
Diagnostic: (a) compute $\kappa(\mat{A}_{\text{step}})$ at the
first step; if $\kappa < 10^{6}$, conditioning is not the
primary cause. (b) Repeat the simulation with a perturbed initial
condition and check whether the failure time shifts in the
predicted exponential-growth way. (c) Linearise the reactor
equations and compute the dominant eigenvalue; if positive, the
physical model is genuinely unstable in the operating regime, and
the simulation is correctly reporting the instability rather than
failing. The fix differs by cause: precondition the matrix,
shorten the step, or refuse the simulation as outside the model's
domain of validity.

\subsection*{Judgment}\pathtag{standard}

\begin{exercise}
A vendor reports a polynomial calibration of degree $9$ for a
sensor, with the coefficients given to $12$ significant digits.
Argue, in one paragraph, whether the precision is meaningful
and what the engineer should ask the vendor before adopting the
calibration.
\end{exercise}

\paragraph{Discussion.}
Twelve significant digits is not meaningful for a degree-$9$
polynomial fit unless the design matrix's condition number is
disclosed and below about $10^{4}$. By the rule of thumb, a
degree-$9$ Vandermonde fit on unrescaled data over a finite
interval has $\kappa \gtrsim 10^{13}$, which leaves at most three
trustworthy coefficient digits in double precision. The reported
twelve digits are an artefact of the floating-point representation,
not of the calibration itself. The engineer asks the vendor: (a)
$\kappa(\mat{V})$ of the design matrix; (b) the basis used
(monomial, Chebyshev, B-spline); (c) the data interval and whether
$x$ was rescaled; (d) the residual norm and the out-of-sample
prediction error; (e) the factorisation method. A calibration with
no answers to (a) and (b) is not adoptable for traceable
measurement.

\begin{exercise}
A senior colleague proposes solving a $10^{6} \times 10^{6}$ dense
linear system by computing $\mat{A}^{-1}$ explicitly and
multiplying. Argue, in one paragraph, why this is the wrong
approach, and propose the right approach for the most likely use
case (multiple right-hand sides on the same matrix, single
right-hand side, structured matrix).
\end{exercise}

\paragraph{Discussion.}
A $10^{6} \times 10^{6}$ dense matrix stores $10^{12}$ doubles, or
$8 \,\text{TB}$, before the engineer attempts any factorisation;
the inverse would require equal storage and a further $\oom{n^{3}}
= 10^{18}$ flops to form. Even on a $10^{15}$-flop-per-second
exascale machine the inversion is a multi-thousand-second task,
and the resulting $\mat{A}^{-1}$ carries the squared conditioning
of the original solve, so the matrix-vector multiply $\mat{A}^{-1}
\vect{b}$ is less accurate than a direct solve. The right approach
depends on the use case. Multiple right-hand sides on the same
matrix: compute the LU factorisation once and reuse it (the cost
of the second and subsequent solves is $\oom{n^{2}}$, not $\oom{n^{3}}$).
Single right-hand side: direct LU solve, no inverse formed.
Structured matrix (sparse, banded, Toeplitz, low-rank, hierarchical):
use a structure-exploiting solver whose complexity is sub-cubic;
a $10^{6} \times 10^{6}$ dense system is almost always actually
sparse and the structure is the leverage.

\begin{exercise}
An engineer is offered two implementations of a least-squares
solver: the normal-equations formulation, which is faster, and
the QR formulation, which is slower but more accurate. The
application is a real-time signal-processing pipeline whose
design matrix has condition number $10^{4}$. Argue, in 200 to
300 words, which implementation the engineer should choose and
state the conditioning threshold above which the choice would
flip.
\end{exercise}

\paragraph{Discussion.}
At $\kappa(\mat{A}) = 10^{4}$ the Gram matrix has $\kappa \approx
10^{8}$. By the rule of thumb, the normal-equations solution loses
about $8$ decimal digits, leaving $16 - 8 = 8$; QR loses about
$4$, leaving $12$. Both are well above what a single-precision
implementation would deliver and both are far above what a real-time
signal-processing pipeline typically requires (six to eight digits
in the fitted coefficients suffices for almost any audio, radar,
or sensor-fusion application). The engineer's call is between
roughly twice the speed (normal equations, $\tfrac{2}{3}n^{3}$
flops versus QR's $\tfrac{4}{3}n^{3}$) and four extra digits.
For a $10\,\si{\milli\second}$ real-time budget at $n \approx 100$,
the choice favours normal equations. The threshold flips when
$\kappa(\mat{A}) > 10^{7.5}$: the Gram conditioning passes
$10^{15}$, leaving normal equations with no usable digits, and
the engineer must use QR or refuse the formulation. The reflection
should also note that the choice is reversible: a production
pipeline can monitor $\kappa(\mat{A})$ at runtime and switch
methods if the conditioning drifts.

\begin{exercise}
A research group reports a linear-algebraic claim with no
condition-number information and no statement of which
factorisation was used. Argue, in 200 to 300 words, what the
engineer should ask the group before treating the claim as
load-bearing for design, and what disclosures the engineer should
demand for any future submission.
\end{exercise}

\paragraph{Discussion.}
The engineer asks: (1) what is the condition number of the
underlying matrix (in the relevant norm, with the chosen
factorisation); (2) which factorisation was used (LU, Cholesky,
QR, SVD), with the implementation library and version; (3) what
the residual or backward error was, in the same norm; (4) whether
any singular values were below a cutoff and what cutoff was
applied; (5) what the precision of the floating-point representation
was and whether iterative refinement was used. Without (1) and
(2), the reported numerical result has no defensible accuracy
estimate; refuse the claim as load-bearing for design. For future
submissions the engineer demands: a numerical-disclosure section
in the paper or report that records the five answers, with the
provenance of the matrix (synthetic, measured, simulated) and the
seed of any random generation. The disclosures are not an
imposition; they are the minimal documentation a reader needs to
reproduce the claim and to assess whether the reported precision is
real. The discipline that section 9.8 develops, refusal of an
unsafe numerical operation, generalises to refusal of an
undocumented numerical claim.

\begin{exercise}
A reader of this chapter argues that, since the SVD reveals
everything the conditioning analysis hides, every linear-algebra
problem should be solved by the SVD. Argue, in one paragraph, why
this is wrong as a default, identifying both the cost argument
(SVD is more expensive than QR or LU) and the structure argument
(many problems have structure that LU or QR exploits and SVD
discards).
\end{exercise}

\paragraph{Discussion.}
The SVD costs $\oom{m n^{2} + n^{3}}$ flops with a constant
several times that of QR and an order of magnitude above LU. For
a square $n \times n$ system, LU at $\tfrac{2}{3}n^{3}$ delivers
the answer; SVD at $\approx 22 n^{3}$ delivers the answer plus
the full spectrum the engineer rarely needs. The default-SVD
position pays the spectral cost on every solve. The structure
argument is the stronger one. Sparse matrices: LU and Cholesky
preserve much of the sparsity (with fill-in control by reordering);
SVD destroys it. Banded matrices: banded LU is $\oom{n b^{2}}$;
SVD ignores the bandedness. Symmetric positive-definite matrices:
Cholesky is half the cost of LU; SVD is far more. Toeplitz,
circulant, and hierarchical matrices: each has a sub-cubic
specialised solver; SVD discards the structure. The defensible
default is to use the cheapest factorisation that exploits the
matrix's structure and to escalate to SVD only when the
conditioning analysis requires the spectral information.
